% !TEX program = pdflatex
% ---------------------------------------------------------------------------
% Auto-generated by code/analysis/scripts/merge_final_reviews.py
% Do not edit by hand: regenerate instead.
%
% Section bodies are inlined verbatim from
%   code/analysis/latex/agentic/<track>/overall/paper/sections/
%   knowledge/overall/paper/main.tex
%   frontier/overall/paper/main.tex
% ---------------------------------------------------------------------------
\documentclass[11pt]{article}
\usepackage[utf8]{inputenc}
\usepackage[T1]{fontenc}
\usepackage{lmodern}
\usepackage[margin=1in]{geometry}
\usepackage{csquotes}
\usepackage[style=authoryear,backend=biber,sorting=nyt,maxbibnames=99,giveninits=true,uniquename=init]{biblatex}
\usepackage[hidelinks]{hyperref}
\addbibresource{../references.bib}
\title{Tile Drainage: Agent-Generated Review\\[0.3em]\large State of Knowledge and Research Frontier}
\author{}
\date{}
\begin{document}
\maketitle
\tableofcontents
\clearpage


\part{Part I --- State of Knowledge}
% ===== track: knowledge -- Tile Drainage: State of Knowledge =====

% ----- knowledge/00_abstract.tex -----
\section{Abstract}

Artificial subsurface (tile) drainage is among the most consequential
anthropogenic modifications of the field-to-stream pathway, draining roughly
$12$--$17\%$ of the world's cultivated land and more than $22$~Mha in the
United States alone \parencite{Valipour2014,Valayamkunnath2020}. This review consolidates the state of knowledge on
drained-land hydrology, water quality, agronomy, and management, emphasising
well-supported findings, mechanisms, and quantified magnitudes over open
questions. Its unifying conclusion is that drainage is fundamentally a
hydrological intervention: by intercepting water that would otherwise recharge
groundwater or leave as surface runoff and routing it rapidly to streams, tiles
commonly carry $40$--$65\%$ of streamflow in humid settings, dominate field-scale
water export, elevate low flows and baseflow, shorten travel times, and
homogenise catchment response \parencite{King2014,Schilling2019}. This hydrology governs the water-quality
outcomes that follow. Relative to undrained or surface-drained land, subsurface
drainage reduces runoff, sediment, and phosphorus but opens an efficient conduit
for nitrate, the dominant and best-quantified export; nitrate loads are
hydrologically controlled, strongly seasonal, and concentrated in high-flow
episodes, with concentrations routinely near or above the $10~\mathrm{mg~L^{-1}}$
drinking-water limit. The presumption of a filtering ``clean'' subsurface pathway
has been overturned: a small, hydraulically active volume of macropores and
cracks delivers particulate phosphorus and sediment to drains within hours \parencite{Pease2018}, so
tiles frequently dominate edge-of-field phosphorus loads \parencite{Pease2018,Sharpley1995}. Against these losses,
a tiered hierarchy of in-field, drainage-design, and edge-of-field practices
produces quantified but scale-attenuated reductions---cover crops clustering
around $40$--$60\%$ nitrate-load reduction \parencite{Malone2018,Malone2014b,Malone2023}, alongside controlled drainage,
bioreactors, wetlands, and saturated buffers governed by residence time,
temperature, and loading \parencite{Christianson2013,Jaynes2008}---so meeting downstream targets requires stacking
complementary measures. Drainage also reliably raises yields, reclaims saline
and waterlogged land, and returns favourable economic and social benefits, while
validated process models, geophysics, and remote sensing underpin design and
mapping. Because independent methods, climates, and scales converge on the same
mechanisms and comparable magnitudes, these conclusions can be stated with
confidence, and recognising drainage water quality as a design criterion is
their consolidated implication.

% ----- knowledge/01_introduction.tex -----
\section{Introduction}

Artificial subsurface (tile) drainage is among the most consequential
anthropogenic modifications of the field-to-stream pathway in agricultural
landscapes. By lowering the water table and routing excess soil water rapidly
to outlets, drainage converts poorly drained and waterlogged land into some of
the world's most productive and intensively farmed cropland. Its footprint is
correspondingly large: roughly $12$--$17\%$ of cultivated land is artificially
drained worldwide \parencite{Valipour2014}, about three-quarters of surveyed
United States producers report drained cropland \parencite{Morton2015}, nearly
$90\%$ of Finnish farmers rate subsurface drainage as important---ahead of
irrigation \parencite{PeltonenSainio2021}, and remote-sensing inventories now
place the drained United States extent at more than $22$~Mha
\parencite{Valayamkunnath2020}. Because tile networks intercept water that would
otherwise recharge shallow groundwater or leave as surface runoff and deliver it
directly to streams, drainage behaves as a master variable of drained-land
hydrology and water quality: in humid settings tiles commonly carry a large and
measurable fraction of streamflow and dominate field water export
\parencite{King2014,Williams2018,Schilling2019}. Understanding this engineered
re-plumbing of the landscape is therefore central to managing both agricultural
productivity and the quality of the waters that receive drainage.

The hydrologic changes wrought by drainage are double-edged. Relative to
undrained or surface-drained land, improved subsurface drainage generally
reduces surface runoff, peak flow, sediment, and particulate phosphorus, yet it
opens an efficient conduit for nitrate, coupling tile-drained cropland to
hypoxia in the Gulf of Mexico and to the eutrophication of inland waters
\parencite{FAUSEY1995,SKAGGS1994,THOMAS1995,Gramlich2018,Raymond2012}. The
once-common presumption that the soil profile filters subsurface water into a
``clean'' pathway has been overturned: a small volume of hydraulically active
macropores delivers phosphorus, sediment, and conservative tracers to drains
within hours of application or rainfall \parencite{Pease2018}. At the
same time drainage remains a foundational production technology that raises
yields on excess-water-stressed soils, reclaims saline and waterlogged land, and
returns favourable financial and social benefits \parencite{Singh2022b,Yang2025}.
Reconciling these competing outcomes---water-quality protection against
agronomic and economic gain---is the defining management challenge of drained
agriculture.

The knowledge synthesised here rests on several mutually reinforcing lines of
evidence. Continuously monitored plots, fields, and watershed outlets quantify
flows and loads; tracer, isotope, and geochemical studies partition flow
pathways and fingerprint contaminant sources; controlled columns, tanks, and
lysimeters isolate mechanisms; long-term records spanning decades reveal trends
and legacies; and a mature toolkit of process-based and data-driven models,
together with near-surface geophysics and remote sensing, simulates and locates
buried networks
\parencite{Smith2011,Singh2019b,Gowda2012,Allred2020}. Where these independent
approaches converge, findings can be stated with strong confidence, and it is
that consolidated evidence---rather than open questions---that this review sets
out to establish.

This paper is the state-of-knowledge component of a two-part synthesis. It
consolidates what the evidence has established about agricultural subsurface
drainage---the consolidated findings, mechanisms, and quantified outcomes that
are well supported---and how strongly each is supported; the open questions,
contested results, and emerging frontiers are treated in a companion frontier
component. The organisation follows the causal logic of drained-land systems.
Because hydrology is the master control on everything that follows, we begin
with drainage flow pathways, the modified flow regime, and its scaling. We then
turn to contaminant loss and its mechanisms---nitrate export and its biophysical
controls; phosphorus and sediment loss and the preferential transport that
drives it; and the further dimensions of pesticides, microbes, trace elements,
and greenhouse gases. Management is treated next along the pathway from field to
outlet: in-field nutrient and cropping practices, drainage design and
water-table management, and edge-of-field treatment technologies with their
watershed-scale integration. We then consolidate the agronomic, reclamation, and
socio-economic benefits of drainage, and close with the modeling, measurement,
and mapping toolkit on which the physical and management conclusions rest.
Throughout, quantified magnitudes are developed within their respective thematic
sections; the aim of this synthesis is to bring those threads together into a
coherent account of what is now known about tile drainage.

% ----- knowledge/02_drainage-hydrology-flow-pathways-regime-and-scal.tex -----
\section{Drainage Hydrology: Flow Pathways, Regime, and Scaling}

Subsurface drainage reshapes the field and catchment water balance in ways that are now well quantified across observational, tracer, and modelling lines of evidence. Tile networks intercept water that would otherwise recharge shallow groundwater or return as surface runoff and deliver it rapidly to streams, so drained landscapes carry a large and measurable share of streamflow, a distinctive flow regime, and a homogenised, scale-invariant response. This section establishes that physical hydrology---the flow pathways, the modified regime, and its scaling---on which the later water-quality and management findings depend.

\subsection{Tile Drainage as a Dominant Flow Pathway}

Watershed-scale studies consistently establish that tile drains route a large fraction of precipitation and streamflow. In an Ohio watershed, 21\% of precipitation (206~mm) was recovered through tiles, accounting for 41\% of annual and 47\% of monthly discharge \parencite{King2014}, and equivalent-medium modelling of an Iowa watershed attributed 97~mm of tile drainage---13\% of precipitation and 39\% of streamflow \parencite{Thomas2016}. In the Boone River, end-member mixing and \textsc{SWAT} bracketed the tile contribution at roughly 46--54\% of discharge, with \textsc{SWAT} estimating about 53.7\% (Boone) and 65.8\% (Lyons Creek) of annual yield \parencite{Schilling2019}; snowmelt-dominated modelling placed tile flow near 16\% of annual precipitation (about 37\% of yield), with tiling reducing surface runoff and increasing subsurface flow and yield \parencite{Rahman2014}. Field-scale export can be even more tile-dominated: tile discharge accounted for $83\pm20$\% of Ohio water export \parencite{Williams2018}, and most discharge from an Indiana closed depression left through tiles rather than surface flow \parencite{Williams2019}. Across humid settings the drain share commonly falls in the 40--65\% range.

Partitioning between tile and surface pathways varies systematically with hydrologic condition and drainage area. In Iowa, tile drainage averaged 30\% of streamflow during events but 61\% during non-event periods; the event ratio $Q_T/Q$ rose with drainage area while the non-event ratio fell, and annual $Q_T/Q$ increased log-linearly from 0.29 at 6~km$^2$ to 0.39 at 45~km$^2$ ($R^2=0.98$) \parencite{Thomas2016,King2014}.

Tile contributions peak when soils are wet and evapotranspiration low. Ohio tile discharge was strongly seasonal (median monthly summer 0.9~mm versus winter 23.4~mm) \parencite{King2014}, Boone River drainage dominated flow from March through June \parencite{Schilling2019}, and simulations indicated the largest baseflow effect from March to July---over 40\% more baseflow in drained watersheds and about 60\% more in April--June \parencite{Schilling2008}. In cold regions tile flow tracks snowmelt, with a seasonal absence from May to September \parencite{Kompanizare2024}, and water-table position relative to tile depth gates drainage: in a well-drained Iowa block the water table lay below tile depth 72\% of the time, steepening gradients during wet periods \parencite{Logsdon2020}.

\subsection{Modification of the Flow Regime}

\subsubsection{Low Flows and Baseflow}

A robust, repeated finding is that tile drainage elevates low flows and the baseflow fraction. Iowa modelling showed drained landscapes carried far higher baseflow (150~mm, 72\%, at 24~m spacing) than undrained (52~mm, 25\%), with drainage acting primarily on the baseflow portion of the hydrograph \parencite{Schilling2008}. Ohio simulations found tiles greatly increased low flow while decreasing higher flows less, raised annual average daily flow about 5\%, and lowered the probability of the water table entering the root zone from about 50\% to about 5\% \parencite{Rathore2025}. Routing analysis reinforced that tiles increase low flows, decrease intermediate flows, and minimally affect the largest flows \parencite{Sloan2017}. In Latvia, drainage raised mean summer--autumn runoff by 132--186\% and winter low-water runoff by 53--286\%, with annual mean, minimum, and maximum discharge up 46\%, 158\%, and 26\% \parencite{Apsite2017}. Channelized Ohio headwaters showed far fewer zero-discharge days (8 versus 87.5~d~yr$^{-1}$ in natural streams), inverting part of the accepted headwater paradigm \parencite{King2009}.

\subsubsection{Peak Flows and Flood Response}

Effects on the largest peaks are consistently modest, but scale- and condition-dependent for smaller ones. Routing modelling concentrated peak-flow reductions at intermediate scales (100--1000~km$^2$), falling from 26\% at the watershed to 11\% at the basin outlet for one event, with negligible field-scale effect on the largest floods; placement mattered, with draining headwaters producing the greatest reduction \parencite{Sloan2017}. Field-scale Iowa scenarios showed no significant difference in mean annual peak discharge across six soil types, indicating drainage chiefly alters pathway rather than the largest peaks \parencite{Sloan2016}. Ohio simulations found tiles could enhance event peaks for sub-threshold or dry-condition events but reduced them for larger events, cutting streamflow in 23 of the 25 largest events \parencite{Rathore2025}. Under expanded tiling, smaller peaks became more frequent while larger peaks declined, steeper rising limbs implying that reduced large-event peaks arose from reduced volumes rather than altered timing \parencite{Rahman2014}. Distributed modelling similarly found drainage decreases peak flow while increasing recession volume, reproducing observed behaviour with only 21\% of the area tiled \parencite{Fonley2021}.

\subsubsection{Flashiness, Drought, and Long-Term Trends}

Flashiness responds to drainage in a soil- and slope-dependent way. A global meta-analysis of 464 watersheds found overall flashiness (Richards--Baker Index) similar regardless of drainage extent but diverging by setting: in flat watersheds ($<$2\% slope) medium/high drainage raised the index, whereas in sloping watersheds low extent had the highest, with drainage extent correlating positively with agricultural land use and clay content and negatively with slope \parencite{Zhou2025}. Field-scale Iowa simulations showed drainage decreased flashiness for low-conductivity soils and increased it for high-conductivity soils, homogenising response with an optimal spacing near 15~m \parencite{Sloan2016}, while Ohio modelling found overall flashiness decreased with tiles, minimised near 50~m spacing \parencite{Rathore2025}. Across 30 Nordic--Baltic catchments, flashiness related negatively to baseflow index and positively to drained share and intensity, and computing it from hourly rather than daily discharge increased it nearly threefold \parencite{Deelstra2014,Deelstra2014b}. At the regime scale, drainage significantly increased streamflow-drought duration and intensity while reducing frequency, each 10\% increase in tile-drained area lengthening drought duration by 0.03~d, raising intensity 12\%, and lowering frequency 0.10 events~yr$^{-1}$ \parencite{Adelsperger2024}. Indiana trend analysis found increasing low- and medium-flow trends with weakened precipitation linkage in heavily drained areas \parencite{Kumar2009}, although a synthesis of seven Midwestern rivers concluded that inter-annual variability in streamflow and baseflow was controlled primarily by current- and previous-year precipitation \parencite{Baeumler2020}.

\subsection{Scale Dependence, Homogenization, and Travel Times}

Artificial drainage homogenises catchment response and reduces the scale dependence of hydrologic signatures. Non-tiled Iowa watersheds show decreasing area-normalized peak discharge and increasing baseflow ratio with scale, whereas tiled basins behave scale-invariantly; within-basin variability in recession coefficients declines as drained fraction rises, and the tiled/non-tiled distinction disappears above roughly 2200~km$^2$ \parencite{BolandBrien2014}. In Minnesota's Redwood basin, mean travel time fell 41\% (77 to 45~d), year-to-year variability was reduced, and travel time became independent of catchment size above about 200~km$^2$ \parencite{DaneshYazdi2016}, with streamflow dynamics becoming more linear after 1976 despite little change in precipitation \parencite{FoufoulaGeorgiou2015}. Isotope work in a UK catchment found short mean transit times ($<$1~yr, faster in artificial drains) and concluded that drainage homogenised distributed responses \parencite{DimitrovaPetrova2020}, and an Indiana watershed behaved as an effectively homogeneous unit describable by a single exponential travel-time distribution with a 2--3~d mean \parencite{Basu2010}; even so, tile contributions remained detectable beyond 16{,}000~km$^2$ in the Boone River \parencite{Schilling2019}.

Modelling and tracers quantify how drainage shortens groundwater travel times. In Iowa, raising drainage density shortened mean travel time and raised the tile contribution to baseflow to an upper bound of 37\%, while increasing channel incision from 0.3 to 2.7~m lowered mean travel time near-linearly and raised the tile contribution from 30\% to 54\% \parencite{Schilling2012}. In Bear Creek, mean groundwater travel time fell from 5.57 to 1.07~yr at high drainage intensity (versus 181.6~yr presettlement), cutting riparian-buffer residence time to 15.5--36.7~d \parencite{Schilling2015}. Analytical solutions confirm that under tile drainage groundwater levels are lower and decline faster, total groundwater discharge to the river increases, and the tile contribution grows with density, though direct groundwater discharge during recession is reduced \parencite{Liang2015}; recession of heavily tiled watersheds is correspondingly more linear \parencite{Schilling2008}.

Runoff generation is governed by pathway-activation thresholds. In an Indiana closed depression, rainfall alone related linearly to cumulative flow with no threshold, but a clear threshold emerged when antecedent soil-moisture deficit and rainfall were summed; tile flow generally began earlier, peaked earlier, and lasted longer than surface flow \parencite{Williams2019}. Tile celerity there was controlled by antecedent wetness---preferential flow mattering more when dry and stored pre-event water dominating when wetter---with groundwater the primary tile source (mean transit time 245~d) \parencite{Williams2021}. In near-level Canadian vertisols, overland flow generally activated before tile flow, and direct preferential flow into tiles was uncommon \parencite{Kokulan2019}. Across 795 US--Canadian events, 43\% produced both surface and tile flow and 57\% tile flow only, surface flow contributing a median 12\% among dual-pathway events \parencite{Macrae2019}; in Kansas tile-outlet terraces, baseflow contributed 38--51\% of discharge \parencite{Stops2022}.

Tracer and geochemical studies partition tile flow into fast preferential and slower matrix components whose balance shifts with season and wetness. In a French field, dry-condition rainwater transferred directly through macropores whereas saturated conditions were dominated by pre-event soil water; surface and tile drainage transferred 44\% and 20\% of annual effective rainfall, with 65\% of infiltrating water not captured by tiles \parencite{Gaillot2023}, and catchment monitoring found the fast preferential peak contributing about 15\% of water and sediment fluxes \parencite{Grangeon2021}. Ohio master-recession analysis apportioned annual tile discharge into 9\% quick-old, 39\% quick-new, 42\% slow-old, and 10\% slow-new components \parencite{Nazari2021}, preferential flow accounted for a higher fraction in clay than loam \parencite{Pluer2020}, and Swedish work attributed autumn events to preferential flow and spring events to matrix flow \parencite{Salazar2011}. Flow is highly concentrated in time: in Ohio, surface runoff occurred $41\pm42$ and tile drainage $240\pm69$~d~yr$^{-1}$, with 80\% of tile discharge delivered in 21\% of the year \parencite{Williams2018}.

\subsection{Repartitioning of Water Relative to Undrained Land}

Relative to undrained or surface-drained conditions, improved subsurface drainage generally reduces surface runoff, peak flow, sediment, and phosphorus even as it increases nitrate delivery \parencite{FAUSEY1995,SKAGGS1994,THOMAS1995,Gramlich2018}. In Oregon, drainage reduced watershed runoff and sediment yield by about 65\% and 55\%, eliminated baseflow, and lowered peak flow \parencite{ISTOK1983}; Albanian sites recorded surface-runoff reductions of 30--56\% \parencite{Grazhdani1996}; and Minnesota modelling found reduced runoff, soil loss, and sediment yield, especially under no-till \parencite{Maalim2013}. In North Carolina, improved drainage cut surface runoff by 192~mm~yr$^{-1}$ (66\%) while raising total outflow by 10\% and strongly attenuating peak runoff (101 to 28~mm~day$^{-1}$) \parencite{KONYHA1992}.

At the field-water-balance scale, the tile drain, not the soil surface, is usually the principal conduit in humid settings. In Indiana, surface runoff was only 1--10\% of annual precipitation versus 26--45\% for subsurface drainage (simulated balance 7\% runoff, 34\% drain flow, 51\% evapotranspiration), with all ponding generated by saturation excess \parencite{Saadat2020}. Iowa pothole modelling found tile discharge and evapotranspiration accounting for 52.4\% and 47.5\% of precipitation \parencite{Amado2018,Martin2019}, and UK field drains were the dominant pathway transferring runoff and most sediment to the stream, responding in every monitored event \parencite{Deasy2009}. The direction of these shifts is texture-dependent: on sandy loams subsurface drainage minimised runoff and total-phosphorus loss whereas the reverse held on clay loam, phosphorus speciation being mainly a function of soil texture \parencite{Eastman2010}. Mechanistically, good drainage increases infiltration and lowers solute concentrations in surface runoff \parencite{Tong2016,Tao2017}, and by lowering the water table it limits saturation-excess runoff and surface sealing \parencite{Augeard2005}.

% ----- knowledge/03_nitrate-export-magnitudes-and-biophysical-contro.tex -----
\section{Nitrate Export: Magnitudes and Biophysical Controls}

Nitrate is the overwhelmingly dominant nitrogen species leaving artificially
drained cropland, and its export is one of the best-quantified consequences of
subsurface drainage. Across continuously monitored plots, drainage districts,
and watershed outlets, three findings recur with strong support: the
subsurface pathway carries the great majority of nitrogen loss, the magnitude
of that loss is governed first by hydrology, and the underlying
concentration--discharge behaviour reflects rapid, macropore-mediated delivery
from largely organic-derived nitrate sources that are only partly attenuated by
denitrification. This section establishes those magnitudes and biophysical
controls---the ``what and why'' of nitrate export---and couples load to flow
specifically for nitrate; broader flow-regime behaviour is treated in the
hydrology section, and all management levers are deferred to the management and
design sections.

\subsection{Subsurface drainage as the dominant nitrogen pathway}

\subsubsection{Pathway dominance}

Multiple independent studies establish that in tile-drained landscapes the
subsurface pathway carries most of the nitrogen. In intensively monitored
Norwegian catchments, subsurface drainage carried $76$--$91\%$ of runoff and
$87$--$95\%$ of total-N (TN) loss, with drainage TN concentrations $2$--$4$
times those in surface runoff and nitrate-N constituting $71$--$85\%$ of TN
\parencite{Bechmann2014}. New Zealand mole-pipe systems captured about $62\%$
of surplus rainfall as drainage and routed roughly $80\%$ of soluble N, while
surface runoff dominated P, sediment, and \emph{E.\ coli}
\parencite{Monaghan2016}. At watershed scale, tile drainage accounted for
$56\%$ of annual discharge and $44$--$82\%$ (mean $62\%$) of NO$_3$-N export in
Ohio \parencite{Williams2015}, averaged $57\%$ of streamflow and dominated
nitrate export in Indiana \parencite{Grose2022}, and in Iowa the largest drains
contributed $70$--$87\%$ of exported water and nitrate \parencite{Coupe2020}.
In flat Illinois terrain, subsurface stations lost
$13$--$30$~kg~ha$^{-1}$~yr$^{-1}$ of nitrate-N ($16$--$26\%$ of applied N), with
more than $70\%$ of samples exceeding the $10$~mg~L$^{-1}$ maximum contaminant
level (MCL), while surface runoff was negligible \parencite{Dattamudi2020}.

\subsubsection{Loads, concentrations, and MCL exceedance}

Reported magnitudes span a wide but coherent range. Field- and watershed-scale
nitrate-N yields commonly reach or exceed $20$--$60$~kg~ha$^{-1}$~yr$^{-1}$ in
maize--soybean systems
\parencite{Gentry2009,Gentry2014,Dattamudi2020,Tomer2008}, and tile
concentrations frequently approach or exceed the $10$~mg~L$^{-1}$
drinking-water limit. In three Iowa drainage districts tile water averaged
about $13$~mg~L$^{-1}$ NO$_3$-N and exceeded the standard $85$--$90\%$ of the
time \parencite{Schilling2012b}, and across three catchments the flow-weighted
mean was $17.1$~mg~L$^{-1}$ with $74.2\%$ of samples above the MCL
\parencite{Ikenberry2014}. New Zealand tile drains exported
$14$--$109$~kg~ha$^{-1}$~yr$^{-1}$ TN ($58$--$90\%$ nitrate)
\parencite{Tanner2011}; Illinois riverine and tile yields averaged $30$--$32$
and ranged to $57$--$59$~kg~ha$^{-1}$~yr$^{-1}$ \parencite{Gentry2014}; and very
high concentrations characterise some European and irrigated systems, with
Italian drainage averaging $94$--$216$~mg~NO$_3$~L$^{-1}$ \parencite{ROSSI1991}
and Polish water usually exceeding $70$~mg~L$^{-1}$ \parencite{Dragon2016}. A
key methodological finding is that interannual variability in load greatly
exceeds that of concentration: across 29 Iowa watersheds, flow-weighted
concentration ($0.6$--$15.5$~mg~L$^{-1}$) was a more robust indicator of
management change than load or yield \parencite{Danalatos2022}. Early Iowa work
already characterised nitrate loss as ``highly variable and strongly dependent
on the amount of water lost'' \parencite{BAKER1975}, and long-term records
document real declines: German drainage N discharge fell from $33$ to
$10$~kg~N~ha$^{-1}$~yr$^{-1}$ between the 1980s and 1990s, driven mainly by a
reduced N surplus \parencite{Hirt2005b}.

\subsubsection{Fertilizer input, hydrology, and soil-N legacy}

At regional scale, fertilizer input and watershed hydrology jointly regulate
export. About $34\%$ of applied fertilizer N was exported to streams in the
Mississippi basin, wetter watersheds exported a larger fraction, and tile
drainage acted as a hydrologic modifier that increases throughput and shortens
residence time, raising nitrate loss to the Gulf of Mexico
\parencite{Raymond2012}. Direct comparison confirmed the causal link: nitrate
load from an unfertilized Minnesota field was $58\%$ of that from a fertilized
field \parencite{Ghane2016b}. Multivariate models combining water yield, prior
corn yield, and fertilizer sales explained $96\%$ of annual nitrate-yield
variation \parencite{Gentry2014}. Losses stem partly from asynchrony between
mineralization and crop uptake and from large mineralizable soil-N pools, so
land-management effects may only become visible after years to decades owing to
soil-N legacies \parencite{Cambardella1999,Bauwe2020}.

\subsection{Hydrologic control of nitrate export and seasonality}

The single most robust cross-scale finding is that nitrate load is controlled
by drainage discharge. Annual nitrate-N loss relates strongly to annual flow,
with differences among drain spacings explained by flow
\parencite{Kladivko2021,Bauwe2020}; in Italy NO$_3$ loss correlated with
drainage water ($r=0.93$) but only weakly with precipitation ($r=0.45$)
\parencite{ROSSI1991}. Because export tracks flow, interannual variability is
driven principally by precipitation and drainage volume rather than by
concentration \parencite{Danalatos2022,BAKER1975}.

High flows dominate export. More than half of Iowa catchment nitrate export
occurred during the highest $10\%$ of daily flows \parencite{Ikenberry2014},
and $79$--$95\%$ of Indiana NO$_3$-N export occurred at the 60th--100th flow
percentile \parencite{Speir2021}. Because loads track flow they are strongly
seasonal, concentrating in the non-growing season when evapotranspiration is
low and drainage is high: about $70\%$ of annual drainage and nitrate load
occurred in the November--April fallow season in Indiana
\parencite{Kladivko2021,Kladivko2004}, and $81\%$ of Ohio tile discharge
occurred in the nongrowing season, so nitrate and TN loads were greater then
despite higher growing-season concentrations under corn \parencite{King2016}.
Short episodes can dominate: $21$~kg~N~ha$^{-1}$ was lost in only $11$ days
during snowmelt in Japan \parencite{Hayashi1999}. The broader flow-regime
behaviour underlying this coupling is treated in the hydrology section.

\subsection{Concentration--discharge behaviour, sources, and denitrification}

\subsubsection{Concentration--discharge behaviour and rapid response}

High-frequency records reveal characteristic concentration--discharge (C--Q)
signatures. Annual nitrate export is largely chemostatic---yields rise
roughly proportionately with runoff---and persistent chemostasis has been
documented in Illinois and Ohio tile drainage
\parencite{Speir2021,Hu2024,Williams2015}. Storm-event hysteresis is
predominantly counterclockwise ($77\%$ of Indiana storms; $>75\%$ of a
long-term Ohio record), indicating dominance of distal nitrate sources during
delivery \parencite{Speir2021,Liu2022b}. Analysis concluded that intensive
subsurface drainage, rather than deep groundwater, was the dominant factor
shaping nitrate C--Q relationships \parencite{Liu2022b}. Tile nitrate responds
within hours of water inputs: after irrigation it rose sharply within hours
then dropped abruptly, with preferential flow intercepted by the drain
originating close to and above the tile \parencite{Mohanty1998}, and bypass
(macropore) flow carried most leached N in Japan \parencite{Hayashi1999}.
Threshold behaviour is common: tile discharge began once the soil-moisture
deficit exceeded $-17\pm3$~mm and the water table rose above $0.84\pm0.03$~m
\parencite{Williams2022}, and antecedent moisture modulated discharge and
solute pulses \parencite{Trentman2020,Kennedy2012}.

\subsubsection{Sources, transformation, and denitrification}

Dual-isotope studies indicate that drainage nitrate is often derived largely
from organic sources (manure and soil organic matter) via nitrification, with
denitrification not always dominant. In Ontario, $\delta^{15}$N and
$\delta^{18}$O relationships lacked the 2:1 denitrification slope, pointing to
organic sources and continued nitrification through all seasons
\parencite{Smith2011,Smith2011b}, and rapid nitrification of applied ammonium
($>90\%$ within $72$~h) supplies ready nitrate for leaching \parencite{Hu2024}.
Where denitrification is active it is detectable: along Polish flow paths
nitrate declined while sulfate, excess N$_2$, $\delta^{15}$N, and
$\delta^{18}$O increased \parencite{Dragon2016}, and nested French watersheds
implied $34$--$82\%$ nitrate retention dominated by underground denitrification
\parencite{Billy2011,Billy2013}. In Iowa drainageways, denitrification in
slowly permeable soils could remove essentially all upland groundwater nitrate,
functioning as a natural treatment system
\parencite{Schilling2007,Cambardella1999}. Yet where drains intercept water
before it reaches these attenuating zones, export is high: tile discharge is
dominated by pre-event water---three-component separation gave $65\%$
groundwater and $29\%$ shallow soil water---yet tile drains carried the highest
NO$_3$-N of any source and contributed over $99\%$ of watershed nitrate export
\parencite{Williams2022,Grose2022}.

% ----- knowledge/04_phosphorus-and-sediment-loss-and-preferential-tr.tex -----
\section{Phosphorus and Sediment Loss and Preferential Transport}

\subsection{Overview}

Phosphorus (P) moves in far smaller mass than nitrogen, yet even at low
concentrations it drives eutrophication, and the long-standing presumption
that the soil profile filters subsurface water into a ``clean'' P pathway
has been overturned. A geographically diverse body of monitoring, isotopic
and radionuclide tracing, soil micromorphology, dye and smoke-injection
experiments, and process modelling now supports a well-consolidated
picture. Subsurface drainage typically captures the largest share of field
water and, despite lower P concentrations than surface runoff, frequently
dominates edge-of-field P \emph{loads}, often as particulate P (PP)
delivered by preferential flow
\parencite{Pease2018,VanEsbroeck2016,Uusitalo2001,Laubel1999}. Dissolved P
release is tightly coupled to soil P status, degree of P saturation, and
subsoil sorption, with legacy P sustaining loads at watershed scale
\parencite{Sharpley1995,Beauchemin2003,Jarvie2017}. A tiny, hydraulically
active macropore volume---earthworm burrows, cracks, root channels, and
installation scars---governs rapid solute and sediment delivery and
reorganises soil structure over years to decades
\parencite{Shipitalo2004,Villholth1998,Montagne2008}. Losses are highly
concentrated in time, dominated by the non-growing season and a few large
events \parencite{VanEsbroeck2016,Brendel2019}. This section synthesises
these established magnitudes and mechanisms; P-removal technologies and the
P responses of controlled drainage are treated in the edge-of-field and
design sections, respectively.

\subsection{Phosphorus Loss Pathways and Water Partitioning}

Tile drainage commonly captures the largest share of field water export:
subsurface discharge accounted for $81\pm23\%$ of annual discharge in the
Eastern Corn Belt \parencite{Pease2018} and $78$--$90\%$ in Ontario
\parencite{VanEsbroeck2016}, and installing drainage re-routes water from
surface- to drainage-dominated flow \parencite{Turtola1995,BENGTSON1995}. A
recurring distinction is that surface runoff usually carries higher P
\emph{concentrations} whereas tile drainage---carrying more water---can
dominate annual \emph{loads}. In Ohio, flow-weighted DRP and TP
concentrations were greater in surface runoff, yet subsurface discharge
supplied $71\%$ of DRP and $69\%$ of TP load \parencite{Pease2018}; the same
contrast appears in long-term Swedish monitoring \parencite{Norberg2022}.
Tiles exported $40$--$77\%$ of TP and $19$--$67\%$ of DRP load in Ontario
\parencite{VanEsbroeck2016}, $21$--$52\%$ of TP in Wisconsin
\parencite{Madison2014}, and $65.4\%$ of annual TP at a southern-Ontario
field \parencite{Qi2018}. Models corroborate the partitioning, attributing
$66\%$ of DRP and $76\%$ of TP loss to subsurface drainage against observed
$70\%$ and $80\%$ \parencite{Grewal2025,Sadhukhan2019b}.

Tile dominance is not universal. In snowmelt-affected clay landscapes of the
northern Great Plains, overland flow contributed $72$--$89\%$ of runoff and
$>95\%$ of SRP and TP because tiles were decoupled from the surface by
frozen ground \parencite{Kokulan2019b,Kokulan2021}, and surface runoff also
prevailed under some tillage, pasture, and flooded-paddy systems
\parencite{Zhao2001,Melland2008,DarziNaftchali2016}.

\subsection{Particulate versus Dissolved Phosphorus and Bioavailability}

A central finding is that PP is frequently the dominant form in tile
drainage, contrary to the earlier presumption that the profile filters
subsurface water to dissolved-dominated. In Finnish clays PP was the
grand-mean $92\%$ of TP in both runoff and drainage \parencite{Uusitalo2001};
at Rothamsted's Broadbalk drains PP averaged $54\%$ of TP and was strongly
related to particulate matter \parencite{Lu2003,Lu2004}; and PP dominated
catchment losses in Denmark and elsewhere
\parencite{Grant1996,Brendel2019,vanderSalm2012}. Dissolved forms
nonetheless dominate in snowmelt-driven, grassland, high-STP, and
organic-soil systems: tile SRP averaged $60\%$ of TP across Michigan sites
\parencite{Clement2017}, and DRP dominated fertilised continuous-grass
losses \parencite{Zhang2015b}. Scottish work showed molybdate-reactive P was
only $42\%$ of subsurface TP loss, underscoring the need to measure all P
forms \parencite{Hooda1999}.

Because algal availability, not total mass, governs eutrophication risk,
bioavailable-P estimates matter: desorbable PP and DRP were of equal
importance to algal-available P in Finnish clays
\parencite{Uusitalo2001,Uusitalo2003}, and only $8$--$18\%$ of grassland
lysimeter TP was molybdate-reactive \parencite{Bilotta2008}.

\subsection{Soil Controls on Phosphorus Release and Legacy Phosphorus}

Dissolved P loss is tightly coupled to soil P status. Loss in surface runoff
and drainage related strongly to soil P content
($r^2 = 0.85$--$0.96$ for surface, $0.86$ for subsurface)
\parencite{Sharpley1995}, and higher soil-test P (STP) raised DRP in both
tiles and runoff \parencite{Osterholz2020}. Threshold behaviour is common: a
change point in STP marks where extractable P and DRP rise more steeply
\parencite{McDowell2001b}, and rainfall sharply increased subsurface DRP
where STP exceeded $41$~mg~kg$^{-1}$ \parencite{Ni2020}. Where soils sorb P
weakly, dissolved losses rise: A-horizon degree of P saturation (DPS) alone
explained $83\%$ of tile TP variance, rising to $95$--$96\%$ with
water-extractable P and sorption grouping \parencite{Beauchemin2003}, and a
DPS threshold near $0.25$ predicted sandy-soil SRP \parencite{Siemens2004}.
The highest-leaching profiles combined large labile topsoil P pools with
low-sorption subsoils, so subsoil characteristics must be included in risk
assessment \parencite{Liu2024b}. Subsurface P losses are generally small in
fine-textured soils at moderate STP but become significant in sandy,
organic, and very high-P soils from long-term fertilisation or manuring
\parencite{Sims1998}. Sequential fractionation links pools to loss, with P
lost to water associated with short-term ($<24$~h) exchangeable pools tied to
Al, Fe, and Ca \parencite{McDowell2001}.

Legacy P is manifest at watershed scale: cumulative SRP loads in Lake Erie
tributaries rose $94$--$253\%$ after 2002 even as cropland P balances were
drawn down, $63$--$66\%$ of the increase attributable to greater delivery
consistent with increased tile drainage and connectivity
\parencite{Jarvie2017}.

Organic and drained peat soils are distinctive high-risk sources, where
wetting/drying, organic-matter degradation, and redox chemistry drive
transport \parencite{Grenon2021}. Flooding organic soils increased P release
$4$--$8$-fold, and organic-P mineralisation can supply up to
$160$~kg~P~ha$^{-1}$~yr$^{-1}$ \parencite{Sharpley1995}. Very large filterable
P loads were measured from an Organic soil ($86.7$~kg~P~ha$^{-1}$) versus a
Podzol ($1.7$) \parencite{McDowell2015b}, and degraded drained peatlands were
the highest-risk eutrophication source in a German catchment, governed by
seasonal, redox-mediated remobilisation \parencite{Tiemeyer2009,YliHalla2022}.
Storm-triggered ochre flocs from tile drains carried $27$--$141$~mg~L$^{-1}$
TP with Fe and P co-distributed \parencite{Zimmer2016}.

\subsection{Preferential Flow, Macroporosity, and Drainage-Induced Soil Structure}

\subsubsection{Biopores as the dominant conduits}

Rapid delivery of water and surface-applied solutes to drains is governed by
a small volume of continuous macropores rather than uniform matrix flow. In a
Finnish clay, at $10$~cm above the tile movement occurred exclusively through
earthworm burrows ($106$ dyed burrows~m$^{-2}$), the base of the plough layer
acted as a hydraulic barrier, and infiltration above the tile lines was about
twice that between them and correlated with \textit{Lumbricus terrestris}
numbers \parencite{Shipitalo2004}. Worm burrows were the only structures at
tile depth in a Swiss loam \parencite{Stamm2002}, and roughly $50$ continuous
burrows~m$^{-2}$ reached from surface to $90$~cm at a German site, their
abundance the most important precondition for preferential flow of isoproturon
\parencite{Zehe2001}. Smoke-injection and dye studies confirmed continuous
surface-to-drain pathways clustered within a $\sim$1-m band over the drain,
their number increasing as the water table dropped
\parencite{Nielsen2015,Petersen2012}. The hydraulically active pore volume is
tiny: field analysis indicated an active macroporosity of only $\sim0.17\%$ of
soil volume, yet this fraction produced abrupt drainage concentration
increases by factors of $20$--$400$ within hours \parencite{Villholth1998},
and macropore area fractions peaked below $0.009$, correlating strongly with
saturated conductivity \parencite{Frey2011}.

\subsubsection{Spatial organisation and static structure}

Two organising gradients recur: macroporosity and biopore abundance decline
sharply with depth, whereas lateral variation relative to the drain is real
but weaker. In a Finnish clay, macroporosity fell strongly with depth but
differed little with distance, though earthworm burrows were $44\%$ more
numerous above a drain and field saturated conductivity varied with drain
position \parencite{Alakukku2010,Frey2011}. Reviews report higher
infiltration and lower bulk density above drain lines and runoff reductions of
$34$--$55\%$ relative to undrained soils \parencite{Montagne2009}. Sustained
drainage consistently lowers bulk density and raises air-filled porosity and
gaseous diffusivity: Ohio drainage decreased bulk density and increased
air-filled pores \parencite{HUNDAL1976}, and relative gas-diffusion gains
correlated with corn yield \parencite{Nakajima2014}. Consistent with an
aeration mechanism, drainage raised soil oxygen, mineralisation, and
nitrification \parencite{GROFFMAN1984}, and dynamic aeration indicators
outperformed static bulk density in explaining productivity across a large
Canadian survey \parencite{Caron2025}. The effect is context-dependent,
however: some comparisons found lower bulk density under undrained or
subirrigated soils \parencite{LAL1993,Baker2004}, drainage reduced surface
organic matter and soil organic carbon in some Ohio studies
\parencite{FAUSEY1992,Kumar2014}, and in the most compacted soils aeration
control by drainage becomes impossible \parencite{SMEDEMA1993}.

\subsubsection{Hydraulic evolution, mole drainage, and pedogenesis}

Drainage and associated loosening raise infiltration and saturated hydraulic
conductivity, and these gains often develop over years. Fields drained
$>15$~years earlier had saturated conductivity $7.58$~cm~h$^{-1}$ faster than
recently drained fields, with preferential-flow development the likely driver,
implying that soil physical and biological properties evolve after
installation \parencite{Frankl2023,JafariTalukolaee2016}. Where matrix
permeability is inadequate, mole drainage and subsoiling create the effective
pore network: up to $90\%$ of flow occurred through mole drains on UK heavy
clays \parencite{HARRIS1993}, and moled structure persists as open cavities
for $3$--$7$ years and thereafter as zones of increased permeability
\parencite{MULLER1988,PLAMENAC1988,Castanheira2010}. Over years to decades,
redirected water fluxes accelerate pedogenesis: subsurface drainage
intensified eluviation and clay translocation, removed swelling clays and
$>$one-third of Fe oxides near the drain, and redistributed Mn oxides, with
eluviation restricted to the first metre either side of the drain and
redox-driven changes extending several metres
\parencite{Montagne2010,Montagne2008,Montagne2013,Mercier2000}. About
$70$~years of drainage formed an argic horizon in boreal Finland, a change
judged probably partly irreversible \parencite{YliHalla2009}.

\subsubsection{Macropore transport of solutes and manure constituents}

These structural findings manifest as rapid, connectivity-limited preferential
transport. Conservative tracers reach drains within minutes to hours: chloride
breakthrough increased by factors of $20$--$400$ within hours
\parencite{Villholth1998}; bromide appeared within $1$~h yet cumulative
capture was slow ($8\%$ in $48$~h, $\sim95\%$ only after months)
\parencite{Frey2012,Kladivko2023}. Liquid-manure constituents reached drains
within $5$~min, and soluble P moved through artificial macropores with little
sorption, identifying macropore flow as the primary rapid-P mechanism
\parencite{Stamm2002,Geohring2001}. These behaviours are captured only by
models representing two-domain flow or a direct macropore--drain connection: a
two-dimensional mobile--immobile model matched bromide dynamics far better than
single-porosity formulations \parencite{Köhne2006}; adding a macropore domain
was essential to match drainage chloride \parencite{Villholth1998b}; an
``express fraction'' of $\sim2\%$ reproduced immediate tracer breakthrough in
RZWQM \parencite{Fox2004,Fox2007b}; and a macropore subroutine in APEX raised
DRP-load performance from unsatisfactory to very good, with median annual
macropore-flow and macropore-P contributions of $48\%$ and $43\%$
\parencite{Ford2017}. Fine-textured soils generated more macropore flow than
coarse-textured soils, so soil texture explained connectivity-driven seasonal
variation in flow partitioning \parencite{Poon2021,Holbak2021}.

\subsection{Temporal Dynamics of Phosphorus and Sediment Loss}

Phosphorus loss is highly concentrated in time. Most annual loss occurs in the
non-growing season, which accounted for $83$--$97\%$ of runoff, $84$--$100\%$
of DRP, and $67$--$98\%$ of TP loss in Ontario
\parencite{VanEsbroeck2016,Qi2018}. Event dominance is stark: more than $75\%$
of DRP and $66\%$ of TP export occurred in the top $25\%$ of flows, and a
single 6-day event supplied $46.6\%$ of TP, $84.0\%$ of DRP, and $81.0\%$ of
suspended solids in one subwatershed \parencite{Brendel2019}.
Concentration--discharge behaviour differs by pathway: tile TP and DRP
increased with discharge (chemodynamic) while surface runoff was more
chemostatic, and the tile DRP:TP ratio fell at high flow, indicating a growing
PP contribution \parencite{Osterholz2020}. In cold regions, snowmelt and
winter thaws dominate exports, predominantly in dissolved form---limiting the
value of particulate-focused practices \parencite{Baulch2019,Clement2017}. A
drought-then-deluge sequence in 2021 produced unprecedented P export ($43\%$
of a 12-year DRP load in nine months) \parencite{Millar2025,Mumbi2025b}.

Preferential flow reconciles substantial subsurface PP loss with the filtering
expected of matrix flow. $^{137}$Cs tracing shows drainage particulate matter
originates from topsoil (median grain sizes $2$--$5$~$\mu$m) translocated
within hours \parencite{Laubel1999,Uusitalo2001}, and PP loss is
episodic---nearly one-third of annual PP loss occurred during sampled storms
\parencite{Laubel1999,Grant1996}. Because of this episodicity, discrete
sampling underestimated annual P loss by $26\%$ to $>100\%$
\parencite{Grant1996}. At catchment scale, source-fingerprinting attributed
$28$--$55\%$ of sediment yield to field drains, with topsoil supplying
$70.9$--$88.6\%$ of the load and drains contributing over $50\%$ of the annual
sediment budget in one catchment \parencite{Walling2002,Chapman2005}.

\subsection{In-Field Controls on Phosphorus Loss}

Long-term P inputs raise subsurface P losses, generally in the order
continuous grass $>$ rotation $>$ continuous corn, with fertilised grass
shifting from a P sink to a source \parencite{Zhang2015b}. Swine-manure
compost produced much greater DRP, PP, and TP than controls
\parencite{Zhang2015}, and about $60\%$ of DRP losses from a no-till
cover-crop site followed fall surface fertilisation, highlighting incidental
loss \parencite{Macrae2023,Mumbi2025b}. Across the MANAGE database, generally
less than $2\%$ of applied P was lost in drainage \parencite{Christianson2016}.
Tillage involves a trade-off between dissolved and particulate forms: no-till
and conservation tillage were associated with higher subsurface DRP, whereas
tillage increased PP and sediment
\parencite{Ni2020,Macrae2023,Koskiaho2002}. In snowmelt regions, zero tillage
and perennial forages can \emph{increase} dissolved P despite reducing erosion
\parencite{Baulch2019}. Overall, effective P control depends on matching the
right practice to site hydrology, connectivity, weather, and legacy P, and on
recognising these dissolved--particulate trade-offs
\parencite{King2015,Ni2020,Baulch2019}.

% ----- knowledge/05_in-field-nutrient-and-cropping-management.tex -----
\section{In-Field Nutrient and Cropping Management}

Agronomic management of the crop and its nitrogen supply is the principal
in-field lever on drainage nitrate, complementary to the drainage-system and
edge-of-field practices treated elsewhere. A large plot-, field-,
watershed-, and regional-scale literature converges on a consistent picture:
drainage nitrate concentration and load are governed jointly by nitrogen (N)
input, crop and rotation choice, the amount and seasonal timing of drainage
flow, and soil processing of N, so that no single practice is universally
decisive but several act in predictable directions
\parencite{Dinnes2002}. Under conventional maize and maize--soybean systems
flow-weighted NO$_3^-$-N routinely exceeds the $10$~mg~L$^{-1}$
drinking-water standard---more than $70\%$ of site-years in one twelve-state
synthesis \parencite{Liu2021}, and corn, corn--soybean, and soybean
site-years averaging $14.0$, $13.5$, and $12.1$~mg~L$^{-1}$ across a
continental database \parencite{Hertzberger2019}. The underlying rate--loss
and hydrologic mechanisms are developed in the nitrate-export section; here
the focus is on how rate, source, timing, cover crops, rotation, and tillage
alter that loss and at what cost.

\subsection{Nitrogen Rate, Source, and Residual Soil Nitrate}

Drainage nitrate loss rising with N application rate is the most consistently
documented in-field relationship, and it is paired with a yield plateau near
the economic optimum. In Iowa production fields, yearly nitrate mass loss was
greater for high- than for medium- and low-N treatments, with concentrations
of $5$--$45$~mg~N~L$^{-1}$ \parencite{Jaynes2001}; four-year losses ranged
$22$--$68$~kg~N~ha$^{-1}$ and rose with rate in both continuous corn and
corn--soybean \parencite{Helmers2012b}. Increasing N above the agronomic
optimum yields diminishing returns: at $168$--$252$~kg~N~ha$^{-1}$ in
Illinois maize yield differed little while cumulative loads were
statistically similar \parencite{Li2023b}, echoing early Iowa work that
identified a single $150$~kg~ha$^{-1}$ application as a favorable compromise
\parencite{Azevedo1997}; a $10\%$ rate cut lowered yield without significantly
reducing loss at one site, underscoring site and hydrologic dependence
\parencite{Gentry2024}. Regional syntheses formalize the pattern: maize yield
followed a quadratic response peaking near $193$~kg~N~ha$^{-1}$ while drainage
N loss increased exponentially with rate, and yield-scaled loss rose sharply
once N surplus exceeded $\sim23$~kg~N~ha$^{-1}$, identifying zero surplus as a
practical target \parencite{Zhao2016}. N rate was the strongest predictor of
annual concentration ($R^2=0.59$) with no yield gain above
$150$--$224$~kg~ha$^{-1}$ \parencite{Hertzberger2019}, the most cost-effective
high-rate scenarios also produced the highest losses \parencite{Liu2021}, and
across syntheses each unit rise in rate produced $\sim0.12$~kg~NO$_3^-$-N
loss, rate explaining $27\%$ of variability
\parencite{Chatterjee2020b,Christianson2015b,Ritter2015}.

Residual soil nitrate (RSN) links rate to leaching risk and creates
legacies. RSN was minimally affected below the economic optimum N rate (EONR)
and increased sharply only above it \parencite{Bandura2025}; continuous corn
accumulated up to $225$~kg~ha$^{-1}$ RSN in dry years, flushed to drainage in
subsequent wet years \parencite{Randall2001}, and a very high fall rate
continued to affect nitrate and residual N for more than four years
\parencite{Jaynes2015}. Encouragingly, corn yields rose over three decades at
nearly constant rate while drainage nitrate stayed flat, indicating improved
use efficiency but no loss reduction from yield gains alone
\parencite{Clifford2023}; the leaching breakpoint lies above the agronomic
optimum, with a larger buffer in rotations \parencite{Pasley2021}.

Nitrogen source matters, with manure generally increasing loss relative to
inorganic fertilizer at comparable rates. Liquid swine manure produced
greater concentrations and losses than UAN with no yield difference
\parencite{Bakhsh2005}, liquid-hog manure raised soil and drainage nitrate
above inorganic fertilizer with equivalent yields \parencite{Elmi2005}, and
annually repeated fall manure produced the highest losses, timing of manure
addition controlling loss more than rotation \parencite{HernandezRamirez2011}.
Applying manure before both crops rather than only corn raised loads
\parencite{Dougherty2020}, and fall manure leached more than spring manure
\parencite{Bakhsh2007,Rekha2011}. Organic sources generally boosted yield but
increased losses \parencite{Liu2021}, though one meta-analysis found no
significantly greater dissolved N load \parencite{Christianson2015b}.

\subsection{Nitrogen Timing, Placement, and Enhanced-Efficiency Products}

Shifting N from fall to spring or in-season application generally lowers
drainage nitrate, but effects are smaller and less consistent than rate or
cropping-system changes. In Minnesota rotations fall N without nitrapyrin
produced the highest losses, with reductions relative to fall of
$13$--$18\%$ (spring, split, fall+nitrapyrin) \parencite{Randall2003,Randall2005};
fall N produced greater concentration and load than spring N in Illinois
\parencite{Gentry2024}, and spring or side-dress N cut concentrations while
maintaining yield \parencite{Pittelkow2017,Vetsch2019}. Minnesota work found
spring application raised use efficiency $>20\%$ and reduced tile losses
$\sim36\%$ relative to fall \parencite{Dinnes2002}, though the fall-versus-spring
effect was often smaller than weather, drainage volume, and rate
\parencite{Liu2021}, and out-of-season and pre-plant timings had higher
concentrations than at-plant and side-dress timings \parencite{Hertzberger2019}.

Split, side-dress, and placement effects are mixed. Split fertilization
reduced export only at lower rates \parencite{Liu2021} and was not
significantly different from non-split in a continental synthesis
\parencite{Hertzberger2019}; side-dress at V2--V12 gave no consistent
difference \parencite{Jaynes2013}, and midseason splits sometimes raised
concentration and RSN \parencite{Jaynes2006,Bandura2025}. Split-applied urea
modestly improved water quality \parencite{Waring2022}, and remote-sensing
variable-rate side-dress averaged a $40\%$ reduction over 15 years, mainly via
$33\%$ lower input and avoidance of peak flow, while maintaining yield
\parencite{Wilson2020b}. Placement was inconsistent: injection and
incorporation sometimes had higher concentrations than surface application
\parencite{Hertzberger2019}, and a localized-compaction applicator gave $16\%$
lower concentrations than a spoke injector \parencite{Bakhsh2010}. Among
enhanced-efficiency products, nitrapyrin generally did not reduce nitrate or
improve yield alone \parencite{Vetsch2019,Pittelkow2017} but reduced fall-N
losses $10$--$18\%$ in Minnesota rotations \parencite{Randall2003,Randall2005};
polymer-coated products were more robust, with ESN lowering combined N$_2$O,
NH$_3$, and NO$_3^-$ losses $47$--$49\%$ below urea while maintaining yield,
the largest benefit in wet periods \parencite{Menegaz2025,Nash2015,Nash2015b}.
Reduced-rate summer fertigation cut concentration and load $\sim16\%$ without
reducing silage yield \parencite{Gamble2018}, and fall anhydrous ammonia with
nitrapyrin matched spring application environmentally \parencite{Kaur2024,Kaur2024c}.

\subsection{Cover Crops: Reductions, Mechanisms, and Yield}

Cover crops---particularly fall-planted winter cereal rye---are among the
most consistent and largest in-field reducers of drainage nitrate. The
clearest, most repeatable signal is reduced concentration. In a long-term
Iowa study rye cut concentrations $48\%$ over five years and oat $26\%$,
though neither significantly reduced cumulative loads because annual drainage
varied \parencite{Kaspar2012}; rye reduced concentrations $>30\%$ even at low
biomass \parencite{Waring2020}, and reviews report reductions of
$13$--$94\%$ \parencite{Ward2018}. Where rye is well established, large load
reductions follow: split N plus rye reduced flow-weighted concentration and
load $35\%$ and $37\%$ without yield penalty \parencite{PrezaFontes2021}; rye
added to spring and fall systems cut load $40\%$ and $47\%$, reaching
$61$--$64\%$ in a wet year \parencite{Ruffatti2019}; a no-till rye system
lowered losses $59$--$67\%$ \parencite{OBrien2022}; rye contributed a $43\%$
decrease in cumulative four-year load \parencite{Gentry2024}; and Indiana
watershed monitoring found field-scale reductions of $58$--$99\%$, pooling to
$65\%$ \parencite{Vincent2025}. Multi-year and northern-margin studies show
meaningful but often statistically marginal cumulative effects: rye reduced
cumulative drainage $12.1\%$ and N loss $20.4\%$ over six Iowa years, with
significant reductions only in favorable years \parencite{MartinezFeria2016};
$11\%$ discharge and $13\%$ loss reductions in Minnesota \parencite{Strock2004};
and in a manured silage system rye lowered concentration from $53$ to
$39$~mg~L$^{-1}$ but loads remained high, warranting additional practices
\parencite{Krueger2013}.

The mechanism is assimilation of soil mineral N into biomass during the
fall--spring fallow, shifting nitrate out of leaching-vulnerable depths.
Cereal rye and radish accumulated $130$--$250$~kg~N~ha$^{-1}$ and moved
NO$_3$-N from environmental (20--80~cm) to agronomic depths, high-C/N
overwintering rye best protecting N in wet springs \parencite{Lacey2015};
living cover consistently lowered soil and vadose-water nitrate
\parencite{Bosompemaa2021,Strock2004,Qi2011b}. Effectiveness is therefore
strongly biomass- and weather-dependent: biomass and reductions were weak
under poor conditions \parencite{PrezaFontes2021,OBrien2022}, cereal rye
biomass varied $1.4$--$5.9$~Mg~ha$^{-1}$ (highest under late termination)
\parencite{Gentry2025,Birru2023}, and success depended on early establishment,
achieved in only one of four years after corn in one Minnesota study
\parencite{Strock2004,Bawa2023}. Regional modeling explained $>95\%$ of
site-to-site variation, rye being more effective in warmer, less-wet southern
locations \parencite{Malone2014b}. Effects on drainage \emph{volume} are
modest and often non-significant, so load reductions are driven mainly by
concentration: rye cut May--July drainage $21\%$ via enhanced spring
evapotranspiration but annual drainage only $9\%$ (non-significant)
\parencite{Qi2010b,Qi2011b}, and in the 2012 drought rye did not deplete soil
water at the cash crop's expense \parencite{Daigh2014b}.

Rye cover cropping generally imposes no consistent yield penalty, especially
for soybean \parencite{Strock2004,Waring2020,Qi2011b}, with projected-climate
differences of only $\sim0.1$~Mg~ha$^{-1}$ and no continuous-corn penalty in
simulation \parencite{Malone2020,Birru2023}. Documented exceptions concentrate
in maize, drought years, and silage---a $6\%$ average maize reduction (34\% in
2012) \parencite{MartinezFeria2016} and $16\%$ lower corn silage after manured
rye \parencite{Krueger2013}. Management modulates outcomes: earlier fall
sowing is consistently more effective (mid-September sowing reduced losses
$45\%$ versus October dates) \parencite{Feyereisen2006b}, and delayed spring
termination raises biomass with limited yield cost \parencite{Birru2023,Chatterjee2025}.
In harvested-rye systems, fertilizing the rye boosts biomass and energy with
only a small N-loss cost \parencite{Malone2018,Malone2023}. Synthesizing
across studies, well-established rye clusters around $40$--$60\%$ N-loss
reduction with $25$--$50\%$ concentration reduction
\parencite{Malone2018,Chatterjee2025,Malone2014b,Malone2023}; Iowa
climate-change projections had rye reduce annual loads from $47.5$ to
$18.9$~kg~N~ha$^{-1}$ ($60\%$) under baseline and $63.6$ to
$17.0$~kg~N~ha$^{-1}$ ($73\%$) under future climate, meeting load targets for
all six climate models \parencite{Malone2020}, and watershed models estimate
$15$--$34\%$ reductions in tile nitrate yield
\parencite{David2015,Makhtoumi2025,Gupta2023}.

\subsection{Crop Rotation, Diversification, and Perennial Systems}

Cropping system exerts strong control on drainage nitrate, generally
exceeding tillage and often rivaling N rate. Continuous corn consistently lost
more nitrate than corn--soybean or perennial systems, leading authors to judge
it environmentally unsustainable \parencite{Weed1996,Kanwar1997}; rotations
produced greater yield and lower yield-scaled loss and a larger buffer above
the agronomic optimum \parencite{Zhao2016,Pasley2021}. Combining a shift to
corn--soybean with reduced rate and a cover crop lowered concentration from
$28$ to $8$~mg~L$^{-1}$ over 14 years \parencite{Ritter2015}, and more diverse
rotations delivered large reductions with maintained yield---a three-year
rotation averaged $4.2$ versus $8.0$~mg~L$^{-1}$ for a control, exceeding
$10$~mg~L$^{-1}$ in only $1\%$ of samples \parencite{Gentry2025,Ahmed2007b}.

Perennial and prairie systems are the most effective in-field option. Prairie
concentrations were $\le1.4$~mg~L$^{-1}$, $2$--$195$ times below row crops
\parencite{Daigh2015b}; restored prairie grass reduced losses to near rainfall
levels \parencite{HernandezRamirez2011}; and alfalfa and CRP plots had
$0.3$--$4$ versus $14$--$40$~mg~L$^{-1}$ for row crops, with four-year losses
of $\sim200$~kg~ha$^{-1}$ for corn versus $4$--$7$ for perennials
\parencite{Randall2001,Waring2024}. The benefit can be short-lived on
conversion back to row crops as RSN rises \parencite{Huggins2001}. Converting
corn--soybean area to alfalfa or wheatgrass reduced watershed losses
$78$--$81\%$ \parencite{Wilson2023}, and bioenergy perennials also reduce
loss \parencite{Ferin2023,Woodbury2018,Guo2018}. Diversified rotations, cover
crops, and perennials are recognized as complementary reducers
\parencite{Dinnes2002}.

\subsection{Tillage and Residue Management}

Tillage exerts weak and inconsistent control over drainage nitrate relative to
precipitation, cropping system, and N rate. Over 11 Minnesota years drain flow
was higher under no-till but nitrate losses only slightly higher under
conventional tillage, and moldboard-plow and no-till losses were similar
($43$ and $41$~kg~ha$^{-1}$), cropping system mattering far more
\parencite{RANDALL1995,Randall2001}. In Iowa no-till and ridge-till usually
had the lowest losses but rotation had greater influence
\parencite{Weed1996,Kanwar1997}, regional syntheses found no significant
tillage difference in N load \parencite{Zhao2016}, and conservation tillage
had the highest concentration ($16.1$~mg~L$^{-1}$) \parencite{Hertzberger2019,Wang2016b}.
Where reduced tillage builds soil health its nitrate effect remains equivocal:
no-till lowered loads only in some water years and reduced yields
\parencite{Masri2024}, though it raised soil organic carbon and lowered EONR
\parencite{Cassman2025,Waring2024}. Residue and stover management had limited
effect on nitrate: $50\%$ stover harvest reduced simulated loss only
$1.1$~kg~ha$^{-1}$ \parencite{Malone2019}, residue removal had little impact on
concentration or yield \parencite{Dougherty2020,Guo2018}, and continuous corn
with residue removal plus rye achieved losses similar to other row crops
\parencite{Daigh2015b,Farahani2020}.

\subsection{Soil Testing, Use Efficiency, Regional Aggregation, and Loss-Pathway Tradeoffs}

Watershed-scale soil-test-based programs confirm in-field mechanisms at
landscape scale. Adopting the late-spring nitrate test (LSNT) across a 400-ha
Iowa subbasin reduced concentration by at least $30\%$ via reduced fertilizer
use while matching yields \parencite{Jaynes2004,Karlen2005,Malone2010}, and
soil-test programs generally reduced RSN and loss versus yield-goal management
\parencite{Dinnes2002,Ritter2015}. Alternative farming practices reduced
drainage $41\%$ and concentration from $17.2$ to $8.2$~mg~L$^{-1}$, cutting
losses $59$--$62\%$ \parencite{Oquist2007}, and because $85\%$ of a Minnesota
watershed's load came from tile drainage, substantial perennialization was
needed for a $30$--$50\%$ reduction \parencite{Wilson2023,Makhtoumi2025}. Use
efficiency framing is increasingly emphasized: partial factor productivity was
highest for side-dress and lowest for very high fall N \parencite{Jaynes2015},
and EONR is proposed as a sustainability standard \parencite{Bandura2025};
adoption of cover crops, controlled drainage, and bioreactors remains
constrained by financial and social barriers \parencite{David2015}.

Scaling site results to the drained landscape indicates regionally significant
potential. For the five-state Corn Belt, winter rye was estimated to reduce
nitrate loss to the Mississippi by $\sim20\%$, with adoption on
$34$--$81\%$ of agricultural land and removal costs of
US\$$2.24$--$3.25$~kg$^{-1}$~N \parencite{Kladivko2014}; a statewide Illinois
analysis found rye reduced tile nitrate loss $\sim30\%$
\parencite{Gupta2023b}, and a five-state energy-cover-crop assessment
estimated a $216$~million~kg~yr$^{-1}$ reduction ($27\%$ of the assumed Gulf
load), with nine hotspot clusters on $19\%$ of the area contributing $48\%$ of
the reduction \parencite{Malone2023}.

Finally, mitigation practices redistribute N among loss pathways---leaching,
N$_2$O, and NH$_3$---so single-pathway evaluation can mislead. In-season N plus
rye cut leaching $37\%$ but raised annual N$_2$O $26\%$, leaving total N loss
and environmental-damage cost statistically unchanged \parencite{PrezaFontes2023}.
Yet split N plus rye had the lowest dissolved N$_2$O and nitrate loads, and no
consistent nitrate--N$_2$O tradeoff emerged across systems, showing stacked
practices need not trade one pathway for another
\parencite{Oliveira2025,OBrien2022}; ESN lowered N$_2$O, NH$_3$, and NO$_3^-$
simultaneously \parencite{Menegaz2025}. N$_2$O rose nonlinearly with rate above
a threshold near $120$~kg~N~ha$^{-1}$, and each unit of fertilizer-N produced
$0.01$~kg~N$_2$O-N alongside $0.12$~kg~NO$_3^-$-N loss, subsurface-drained
fields having the highest mean N$_2$O emission
\parencite{Wang2016b,Chatterjee2020b,Fontes2019}. The consistent conclusion
across scales is that combining two or more practices---rate optimization,
in-season timing, cover crops or diversified rotation, and drainage
management---reduces nitrate loss most cost-effectively, since in-field
management alone is often insufficient \parencite{Dinnes2002,Singh2023}.

% ----- knowledge/06_drainage-design-and-water-table-management.tex -----
\section{Drainage Design and Water-Table Management}

Drainage design and active water-table management are the primary levers
through which practitioners set both the volume of subsurface outflow and the
loads it carries. Because solute export in drained cropland is governed first
by hydrology, the geometric choices of spacing and depth and the operational
choice of when to hold or release water translate directly into edge-of-field
flow and nutrient loss. This section synthesises what the evidence has
established about those controls: how spacing and depth trade flow against
loss, how controlled drainage and drainage water management reduce loads in
proportion to the flow they retain, how capture-and-reuse systems recover both
water and quality benefits, and how enhanced configurations and controlled
drainage perform under projected climates.

\subsection{Drain Spacing and Depth as Controls on Flow and Loss}

Drain spacing and depth are well-established design controls on both drainage
volume and nitrate loss. More closely spaced drains produce greater drainage
and greater nitrate-N loss than wider spacings, whereas nitrate
\emph{concentrations} generally do not vary consistently with spacing---load
differences are driven by flow rather than by concentration
\parencite{Kladivko2021,Kladivko2004}. Calibrated models consistently show that
shallower and wider-spaced drains reduce tile flow and nitrate loss at modest
yield cost: reducing drain depth from $1.5$ to $0.9$~m lowered simulated
nitrate loss from $34.0$ to $29.4$~kg~ha$^{-1}$, and widening spacing from $27$
to $60$~m lowered it from $33.8$ to $28.4$~kg~ha$^{-1}$
\parencite{Moriasi2013,Nangia2010c,Davis2000,Hwang2024}. The design tension is
that reducing subsurface loss shifts water to surface runoff: for phosphorus,
increasing spacing reduced subsurface volumes and tile TP loss (more strongly
on finer soils) while increasing surface runoff and its associated TP load
\parencite{Morrison2013}.

These relationships motivate an optimum-intensity principle: use the minimum
subsurface intensity required to relieve excess-water stress, fit the system to
the soil, and apply water-table control where maximum intensity is unnecessary
\parencite{Breve1998}. Deeper drains often raise profit but substantially
increase nitrate losses, and simulation showed that widening spacing to $100$~m
cut total nitrate-N loss by $20\%$ ($37$ to $29$~kg~ha$^{-1}$) at small profit
cost \parencite{ElSadek2002b,Breve1998}; deep drains are frequently unnecessary
for salinity control \parencite{Ritzema2011}. Shallow drainage reproduces much
of the flow- and load-reduction benefit, reducing Iowa subsurface volume by
$46$--$58\%$ and nitrate by $29$--$49\%$ generally without reducing yield
\parencite{Helmers2012,Schott2017}. A multi-state analysis found shallow drains
gave quicker water-table drawdown, reduced yield variability, and $3.4$~cm less
total flow than deep drains, though with more surface runoff and a lower
benefit--cost ratio \parencite{Ghane2021b}, and field comparisons found shallow
systems had the lowest nitrate concentration and loading among configurations,
sometimes at a yield penalty whereas controlled drainage trended toward higher
yields \parencite{Smith2019,Smith2011b}.

\subsection{Controlled Drainage and Drainage Water Management}

Controlled drainage (CD)---raising an outlet control structure to hold the
water table between drainage events---is among the best-established
conservation-drainage practices. A synthesis found CD reduced drainage outflow
by $17\%$ to more than $80\%$ depending on soil, drainage intensity, control
strategy, and location, with nitrogen-loss reduction often about the same
percentage as the volume reduction \parencite{Skaggs2010}. A global
meta-analysis of 84 publications and 1788 paired site-years confirmed that CD
and controlled drainage with subirrigation (CDSI) reduce subsurface discharge
and nutrient losses while modestly increasing yield; CDSI reduced nitrate and
TN more, CD reduced TP more, and CDSI (but not CD) significantly increased
surface runoff and surface nitrate loss \parencite{Wu2025}. Simulations across
48 Midwest locations found CD reduced drainage and N losses everywhere, with
efficacy strongly climate-dependent \parencite{Youssef2018}.

Field studies quantify these effects concretely. Quebec CDSI cut drainage
volume $24\%$, nitrate concentration $25\%$, and nitrate loss $43\%$
\parencite{Ng2000}; a long-term Canadian study found CD reduced total nitrate
loss by $\sim58\%$ and increased yield $3\%$ without elevated greenhouse-gas
emissions \parencite{Madramootoo2024}; and nine years of North Carolina CD
reduced outflow $33\%$ and nitrate export $30\%$ while raising yields
$10$--$11\%$ \parencite{Poole2018}. Growing-season CD reduced water flux
$\sim60\%$, nitrate $51\%$, ammonium $58\%$, DRP and TP $66\%$, and
\emph{E.\ coli} $76\%$, with watershed-scale deployment cutting stream fluxes
of water, N, and P \parencite{Sunohara2016,Sunohara2015}. In Illinois, managed
systems always exported less nitrate than free drainage, with flow reductions
of $35$--$96\%$ and load reductions of $37$--$79\%$ (averaging $61\%$)
\parencite{Cooke2012}.

The mechanism is consistently one of flow--load coupling: syntheses report
drainage water management reducing tile flow by $\sim46\%$ and nitrate loads by
$\sim45$--$48\%$ \emph{mainly by reducing discharge rather than by enhancing
denitrification}, which plays only a secondary role
\parencite{Ross2016,Ward2018}. Two trade-offs qualify the benefit. First,
because reduction operates through volume, drainage management often fails to
lower nitrate \emph{concentrations} even where it lowers loads
\parencite{Grenon2023,Wesström2014}. Second, CD can increase surface DRP and TP
that partly offset subsurface benefits \parencite{King2022,Wu2025}. Timing is
decisive: DRAINMOD analysis identified early March as the most effective
control window (reductions up to $100\%$) whereas April starts resembled free
drainage \parencite{Kesicka2023c}. Greenhouse-gas responses are mixed but
generally not adverse \parencite{Madramootoo2024,Smith2019}.

\subsection{Drainage Water Recycling and Subirrigation}

Drainage water recycling (DWR)---capturing and storing drainage for
supplemental irrigation---benefits both production and downstream water
quality. Yield gains reach $61\%$ for corn and $31\%$ for soybean, with
Canadian N-loss reductions of $40$--$70\%$ and P-loss reductions of
$12$--$36\%$, and modelled downstream N-delivery reductions of $24$--$37\%$
\parencite{Hay2021}. A North Carolina reservoir reduced TN load $47\%$, TP
$30\%$, and sediment $87\%$, with hydraulic retention time the key control,
although the reservoir was a net source of particulate P
\parencite{Moursi2023,Moursi2022}. Resilience is a defining advantage:
multi-location synthesis of 53 site-years showed DWR is more resilient than
free drainage, protecting yield during both drought and excess precipitation
and reducing yield variability by $22\%$ \parencite{Willison2021}. Free
drainage and drainage water recycling raised corn yield $20\%$ and $26\%$
respectively over a nondrained control (recycling gains reaching $48\%$ in
individual years) by avoiding drought losses \parencite{Kaur2021}, and North
Carolina recycling
raised corn and soybean by $8\%$ and $17\%$ with subirrigation raising the
water table $15$~cm \parencite{Moursi2024}. Adoption economics indicate crop
insurance is not a major barrier \parencite{Sellars2022}, and in arid systems
integrated on-farm drainage management reused all but $0.7\%$ of applied water
\parencite{Ayars2014,Kosari2024}.

Subirrigation---delivering water upward through the drains---similarly lowers
losses and stabilises supply. It reduced soil-profile nitrate by up to $50\%$
and drainage nitrate concentrations by $55$--$73\%$, with water-quality
benefits accruing even when outflow occasionally exceeded free drainage,
implicating enhanced denitrification \parencite{Elmi2004,Elmi2002,Noory2009}.
CDSI raised groundwater levels, supplied $170$--$410$~mm of infiltration, and
increased Dutch yields by up to $64\%$
\parencite{deWit2022,deWit2024,FAUSEY1995,Almen2021,Madramootoo2007}. Its
performance, however, is strongly site-dependent, governed by lateral hydraulic
conductivity: across three Ohio sites subirrigation improved relative yield to
$87$--$93\%$ and $80$--$91\%$ at two sites but had almost no effect at Defiance
owing to very low conductivity \parencite{Gunn2018,Gunn2019}. Where hydraulics
are favourable, CD and subirrigation kept soil water optimal in both wet and
dry years \parencite{Niaghi2019}, and in claypan soybean systems subirrigation
improved yield, reduced interannual variability, and allowed earlier planting
\parencite{Nelson2011}; combining drainage water management with foliar
pyraclostrobin produced synergistic soybean gains up to $1070$~kg~ha$^{-1}$
\parencite{Nelson2011b}.

\subsection{Enhanced Configurations and Persistence under Climate Change}

Beyond conventional single-level layouts, engineered configurations extend the
same design principles. Multi-level (bilevel) systems cut the time the water
table exceeded $0.5$, $0.75$, and $1.0$~m by $84\%$, $67\%$, and $41\%$ while
producing much lower shallow-drain salinity \parencite{Hornbuckle2007};
analytical theory shows that adding shallow drains between deep drains raises
the drainage rate and allows wider deep-drain spacing
\parencite{Hornbuckle2012}, with managed deep drains occupying an intermediate
position \parencite{Ayars2006b,Christen2001b}. Engineered enhancements alter
flow further: a permeable filter increased discharge (about $1.9$ times
conventional) and accelerated drawdown while reducing runoff
\parencite{Tao2016,Tao2017}; combined open-ditch-plus-subsurface designs
improved waterlogging removal \parencite{Ren2021,Tao2025}; siphon-type drainage
raised flow rates $27$--$90\%$ \parencite{Ren2023}; artificial macropore tubes
generated flow where none occurred and desalinated the profile
\parencite{Chen2021}; and supplementary drainage of a clay field raised mean
annual discharge $2.6$-fold and cut topsoil runoff by $75\%$
\parencite{Häggblom2019}.

A robust and management-relevant finding is that the benefits of controlled
drainage persist, and in some seasons strengthen, under projected climates. In
Ohio, conventional discharge was projected to decline while CD annual
reductions remained near $4$--$6\%$ and summer growing-season reductions
increased, indicating greater warm-season value \parencite{Pease2017}. Michigan
simulations projected free-drainage discharge rising $20\%$ by mid-century while
CD continued to reduce discharge by $63$--$72\%$, with future and historical
performance statistically indistinguishable \parencite{BinShokrana2023}.
Finnish scenarios projected CD effects remaining similar to the historical
interval \parencite{Salla2022}, and Canadian projections of warmer, wetter
conditions found CD reducing subsurface drainage more strongly than free
drainage, though rising temperature drove dry-stress yield reductions that
motivate the capture, storage, and reuse of drainage water
\parencite{Ndulue2024,Gunn2018}. Chinese and Polish studies likewise found that
appropriate outlet-blocking strategies could reduce future runoff and maintain
groundwater at present levels \parencite{Awad2021,Sojka2020}. Together these
results indicate that the flow--load coupling underpinning water-table
management is durable: because controlled systems act on discharge, their
proportional benefits carry forward into a wetter, warmer future.

% ----- knowledge/07_edge-of-field-treatment-reactive-media-and-water.tex -----
\section{Edge-of-Field Treatment, Reactive Media, and Watershed-Scale Integration}

Once nutrients enter tile drains, in-field management can no longer reach
them, so a family of edge-of-field practices intercepts the load before it
reaches surface waters. Two decades of microcosm and column studies,
pilot- and full-scale installations, and global syntheses have established
constructed wetlands, denitrifying woodchip bioreactors, reactive
phosphorus filters, saturated buffers, and in-channel structures among the
best-characterised measures for artificially drained cropland. Their
performance is governed by a small, recurring set of controls---hydraulic
residence time (HRT), temperature, flow, loading, pH, and cumulative
loading---and their reductions must be interpreted against the fraction of
drainage actually treated and against downstream, watershed-scale targets.
Controlled drainage and drainage water recycling, treated in the design
section, complement these practices.

\subsection{Constructed wetlands}

Constructed wetlands treating drainage are extensively documented nitrate
sinks in which denitrification dominates. Across French watersheds, nitrate
removal averaged $\sim42\%$, of which denitrification supplied $93\%$, and
devoting $\sim1\%$ of the contributing area can achieve $50\%$ removal, with
hydraulic residence time the priority design control \parencite{Tournebize2017}.
Field systems span a coherent range: Iowa wetlands reduced incoming nitrate
$49\%$ \parencite{Drake2018}, a Minnesota wetland removed $67\%$ over four
years \parencite{Gordon2021}, and Illinois wetlands retained $36$--$37\%$ of
N, implying $\sim450$~ha to cut watershed N loading $46\%$
\parencite{Kovacic2006}. Broader syntheses place average total-N removal near
$39\%$ \parencite{Wang2025b}, with roughly $40$--$50\%$ N removal generally
achievable at wetland-to-catchment area ratios of $1$--$4\%$ and only minor
gains beyond \parencite{Mendes2021}. Because performance is a strong function
of hydraulic loading, residence time, temperature, and inlet concentration,
no single-number design is adequate, and higher influent loads raise removal
\emph{rate} but lower percentage efficiency
\parencite{Kadlec2012,Steidl2019,Mendes2021}. Nitrate removal is reliable, but
phosphorus performance is variable and wetlands sometimes export more TP than
they receive \parencite{Tanner2005,Tanner2011,Messer2021}, with plant harvest
an important P route in some systems \parencite{Lenhart2016}. Paired
surface-flow wetland plus woodchip systems act as consistent nutrient sinks
(TN removal $8$--$51\%$, TP $0$--$80\%$), the woodchip component contributing
more removal than the sedimentation basin while occupying a smaller footprint
than a standalone wetland \parencite{Pugliese2020}.

\subsection{Woodchip bioreactors: performance, controls, and mechanisms}

\subsubsection{Field-scale nitrate removal}

Denitrifying woodchip bioreactors reduce nitrate dependably across climates
and designs, with a magnitude that spans a wide but well-explained range.
Consolidating full-size systems worldwide, \textcite{Christianson2021} report
a mean nitrate-N load reduction of $40\pm26\%$ (median $46\%$) and removal
rates of $7.2\pm9.6$~g~N~m$^{-3}$~d$^{-1}$, removing an estimated $20$--$40\%$
of annual nitrate-N loss in the Midwest. The largest compilation---36
site-years at ten full-size Illinois bioreactors---found a median annual
removal of $49$~kg~N per bioreactor, a median edge-of-field efficiency of
$17\%$, and removal in \emph{every} year \parencite{Christianson2026}.
Controlled and cold-climate studies fill in the range: an Illinois field
experiment averaged $63\%$ ($20\%$ at $2$-h to $98\%$ at $8$-h HRT;
$11.6$~g~m$^{-3}$~d$^{-1}$) \parencite{Bell2015}; a Quebec system cut
flow-weighted nitrate-N load $99\%$ and total-N $72$--$73\%$
\parencite{Husk2017}; four South Dakota bioreactors reduced loads $39$--$89\%$
at $5.2$--$26.5$~g~N~m$^{-3}$~d$^{-1}$ \parencite{Thapa2023}; and a New Zealand
system rose from $31\%$ ($2.8$-d HRT) to $93\%$ ($6.4$-d)
\parencite{Rivas2020}. Woodchip \emph{denitrification walls} are equally well
supported: an Iowa wall reduced annual nitrate loss from $54.5$ to
$24.5$~kg~ha$^{-1}$ over nine years with no significant trend
\parencite{Moorman2010}, a wall cut flow-weighted nitrate and mass loss
$\geq55\%$ with no yield penalty and no land taken out of production
\parencite{Jaynes2008}, and a rye cover crop plus wall reduced N leaching
$58$--$59\%$ \parencite{Rogovska2023}. Not all systems meet targets:
underloaded or low-input bioreactors removed only $\sim9.5\%$
\parencite{Bock2018} or $6$--$28\%$ of edge-of-field nitrate \parencite{Maxwell2022}.

\subsubsection{Controls: hydraulics, temperature, and loading}

The dominant, repeatedly quantified control is HRT, which trades off against
the fraction of flow captured: efficiency rises with HRT
\parencite{Bell2015,Rivas2020}, while in the ten-bioreactor dataset annual
flow-weighted HRT correlated with flow captured ($\rho=-0.67$) and
treated-water removal efficiency ($\rho=0.60$)
\parencite{Christianson2026}. Column work shows removal rate per unit wood
rising with flow while efficiency falls \parencite{Greenan2009}, and higher
influent loads raise the removal \emph{rate} but lower percentage efficiency
\parencite{Rivas2020,Christianson2026,Zhang2021b}. Temperature exerts a
strong, consistent effect: exponential functions of water temperature and HRT
explained $73\%$ of the variance in load reduction in Illinois
\parencite{Bell2015}; South Dakota removal fell $\sim30\%$ below
$12\,^{\circ}$C \parencite{Thapa2023}; columns dropped from $96\%$ at
$21.5\,^{\circ}$C to $29\%$ at $10\,^{\circ}$C \parencite{Soupir2018}; and
Michaelis--Menten $V_{max}$ fell sharply from $22$ to $5\,^{\circ}$C
\parencite{Kouanda2021,Zhang2021b}. Kinetically, removal is zero-order when
nitrate is non-limiting and shifts to first-order once nitrate becomes
limiting \parencite{Hua2016,Kouanda2021}.

\subsubsection{Mechanisms, carbon, and microbial ecology}

Denitrification is firmly established as the dominant pathway, accounting for
$>98\%$ of transformed nitrate in microcosms, with immobilisation and DNRA
minor \parencite{Greenan2006,Greenan2009}. Woodchip carbon sustains this
capacity for years: wood raised denitrification potential $31$- to
$4000$-fold over subsoil, supported denitrifier populations $>10^8$~g$^{-1}$,
and gave in situ nitrate half-lives of $0.64$--$0.81$~d
\parencite{Moorman2010}; DNRA and anammox occur secondarily
\parencite{Hartfiel2022,Mendes2021}. Cellulose is the most rapidly consumed
compound, and wood decomposition is depth-dependent (up to $75\%$ loss near
the inlet/top, $<13\%$ at depth), governing longevity
\parencite{Hartfiel2022,Moorman2010,Ghane2018}. Alternative substrates can
raise performance: corncob (CC75) removed more N than woodchips ($61.2\%$ vs
$37.9\%$) at lowest cost \parencite{Law2023}, corn-stover blends enhanced
removal \parencite{Oladeji2023}, and agricultural residues can outperform
woodchips for removal though woodchips give better hydraulics
\parencite{Greenan2006,Kohn2023,Zhang2021b}. Communities are shaped by
substrate, depth, season, and temperature \parencite{Hartfiel2022}, with
cold-active denitrifiers (\textit{Cellulomonas}, \textit{Pseudomonas}) and
nitrate-reducing fungi identified as key taxa
\parencite{Anderson2020,Jang2019,Aldossari2021}; bioaugmentation shows
conditional promise but is limited by washout
\parencite{Jeglot2022,Wang2025b}, and denitrification function is generally
robust to water-level disturbance \parencite{Hathaway2017}.

\subsubsection{Design, sizing, and enhancement}

Sizing balances flow treatment against efficiency: using $10$--$20\%$ of
peak flow at $6$--$8$-h design retention treats the majority of annual volume,
and longer, narrower beds outperform short, wide ones
\parencite{Christianson2013}. Load-duration analysis showed beds sized to a
$0.5$-d HRT can remove $44$--$52\%$ of base-flow nitrate and $20$--$30\%$ of
total annual loss while occupying $<0.27\%$ of watershed area
\parencite{Moorman2015}. Hydraulically, high-aspect-ratio designs approach
plug flow, whereas wide designs suffer dispersion and short-circuiting
\parencite{deOliveira2023,Dougherty2020b}; expanding cross-sectional area
improves flow capture \parencite{Maxwell2022}. Enhancement strategies extend
or intensify treatment: pumping supplies additional nitrate-laden water and
lengthens the season \parencite{Hartfiel2023,DíazGarcía2025}, electrical
stimulation raises removal at a cost \parencite{Law2018}, and heated
tourmaline reduces clogging while raising removal \parencite{Ma2021}. A
catchment-scale three-bed system removed $713$~kg nitrate-N ($12.5\%$ of
catchment load) but required a sediment-exclusion system during high flow
\parencite{Feyereisen2023}.

\subsection{Pollution swapping and greenhouse-gas by-products of bioreactors}

Non-ideal conditions can generate N$_2$O, CH$_4$, sulfate reduction, and
dissolved-organic-carbon (DOC) losses, but reported risks are generally small
\parencite{Hartfiel2022}. Nitrous-oxide emissions are consistently a tiny
fraction of nitrate removed: $0.003$--$0.028\%$ of N denitrified in columns
\parencite{Greenan2009}; no significant difference from controls at an Iowa
wall (indirect emission equivalent to tile drainage without the bioreactor)
\parencite{Moorman2010}; and an average field emission factor of $0.07\%$,
well below the $0.75\%$ IPCC comparison \parencite{Bock2018}. Across five
Illinois bioreactors, only a small fraction of removed nitrate was emitted as
N$_2$O and modelled watershed-scale indirect emissions were \emph{lower} with
a bioreactor than without \parencite{Brunton2024}. Methane and sulfate
responses are more variable: some bioreactors are notable CH$_4$ sources
(soil covers lower surface CH$_4$) \parencite{Brunton2024}, yet a dual-chamber
system's landscape-scale impact was a $<1\%$ increase in CO$_2$e, with
pollution swapping more likely above $50$-h HRT \parencite{Hartfiel2023b}.
Sulfate reduction under non-N-limiting conditions is common, increases with
HRT, and occurred in $83\%$ of Illinois runs
\parencite{Bell2015,Povilaitis2020,Rivas2020}. DOC leaches mainly at
start-up as an initial flush that then declines
\parencite{Abusallout2017,Rivas2020,Greenan2009}; it is dominated by
high-molecular-weight fractions that are disinfection-by-product precursors
but is removable by steel-chip filters or coagulation \parencite{Abusallout2017}.

\subsection{Reactive media for phosphorus removal}

Because woodchips remove little phosphorus---and often leach it at
start-up---dedicated reactive filter media have been developed to retain the
dissolved reactive phosphorus (DRP) that moves through drains in solution
\parencite{Christianson2021}. The evidence converges on a two-family model:
iron- (Fe-)based media (iron oxyhydroxides, zero-valent iron, iron filings,
steel by-products) retain phosphate mainly by strong inner-sphere adsorption,
whereas calcium- (Ca-)rich media (blast-furnace and steel slag, limestone,
lime-amended pellets) retain it by Ca-phosphate precipitation driven by
dissolved $\mathrm{Ca^{2+}}$ and elevated pH; sorption is the preferred
edge-of-field mechanism, and Fe-based media are generally more suitable
because they bind phosphate stably over a wider pH range
\parencite{Mendes2022,Li2020}.

\subsubsection{Media classes and mechanisms}

Iron-based materials are the best-performing class, routinely removing
$\geq95\%$ of phosphate in batch and $>90\%$ in most columns with negligible
desorption \parencite{Allred2014,Allred2012,Allred2017,Mendes2022}. Iron
filings reached the highest capacity of $4.95$~mg~g$^{-1}$ and steel slag
$1.68$~mg~g$^{-1}$, versus only $0.10$--$0.13$~mg~g$^{-1}$ for natural
minerals---an order-of-magnitude advantage for industrial by-products
\parencite{Sellner2019b}. Granular iron-oxide media retained SRP below
detection far longer than steel by-products \parencite{Hauda2020}, Fe-based
CFH media were most robust under variable flow \parencite{Pugliese2023}, and
iron-activated alumina reached $94.4\%$ removal \parencite{Kumari2025}. Ca-rich
slag and cement-kiln blends achieve $>98\%$ removal at low-to-moderate
loading \parencite{Agrawal2011,Sanford2016}, and lime-amended fly-ash pellets
reached $1.98$~mg~P~g$^{-1}$ \parencite{Li2017}, but their performance is
governed by a finite supply of reactive $\mathrm{Ca^{2+}}$: phosphorus is
removed in a moving front as $\mathrm{Ca^{2+}}$ and pH-buffering capacity are
consumed, downstream slag being more efficient than upstream
\parencite{Wang2020d}, and field removal collapses once treated-water pH falls
below $8.5$, with $100\%$ removal sustained only to $\sim500$~mg~P~kg$^{-1}$
slag \parencite{Penn2020}. Consistent with the two mechanisms, lower pH favours
Fe-based adsorption whereas higher pH favours Ca-based precipitation
\parencite{Mendes2022}, and finer particle sizes and coexisting-ion chemistry
(sulfate, bicarbonate, and DOC reduce uptake) modulate performance
\parencite{Li2020,Sellner2019b,Pugliese2023}. Natural minerals such as zeolite
give lower, modest removal ($21$--$58\%$) unless modified \parencite{Agrawal2011,Sadeghi2021}.

\subsubsection{Kinetics, longevity, and field structures}

Phosphate uptake is rapid, permitting short contact times---$10$--$15$~s
sufficed for near-complete removal from spiked drainage
\parencite{Allred2017}---and Fe-based media exhibit low desorption and are
regenerable, with a NaOH or KOH flush restoring capacity
\parencite{Allred2014,Allred2017,Kumari2025,Sellner2019b}. Field structures
confirm both feasibility and the laboratory pH and retention-time limits: a
tile-drain structure removed $19.2$~kg of dissolved P ($55\%$ of input)
\parencite{Penn2020}, an aluminium-treated steel-slag ditch structure removed
$\sim27\%$ of dissolved P and $50\%$ of total P \parencite{Shedekar2020}, and
a compact activated-alumina structure treating the entire flow removed $30\%$
of the dissolved-P and $29\%$ of the nitrate-N load
\parencite{Penn2024}. Field removal typically underperforms the design basis
\parencite{Shedekar2020}, and not all trials succeed---a Texas filter blend
failed to reduce phosphate (while removing chlorothalonil)
\parencite{Agrawal2011b}. Coupling P media with woodchips addresses the
woodchip P deficit: placing iron both upstream and downstream of woodchips
enhances nitrate and orthophosphate removal
\parencite{Goodwin2015,Hua2016,Li2018b,Oladeji2023}, and mixed-media
bioreactors that initially release P become sinks within a year, cutting SRP
$97\%$ \parencite{Husk2018,Zhang2024b}. The central caveat is hydraulic: a
stacked woodchip bioreactor removed $93\%$ of the nitrate entering it, but
because $93\%$ of tile discharge bypassed it, this was only a $5.4\%$ load
reduction---whereas a P-removal structure treating the full flow removed
nearly ten times more nitrate \parencite{Penn2024}.

\subsection{Saturated buffers, denitrification walls, ditches, and in-channel practices}

In-channel and edge structures add interception where drains discharge.
Saturated riparian buffers, which redistribute tile flow laterally through
the riparian soil, offer the largest climate-mitigation potential per unit
area among conservation-drainage practices \parencite{Hall2025}. Riparian
zones remove nitrate efficiently by denitrification---groundwater nitrate fell
from $7$--$8$ to $<0.1$~mg~L$^{-1}$ within $16$--$47$~m
\parencite{JACOBS1985,Schilling2022}---and redirecting tile flow into a buffer
removed nitrate at $1.01\pm0.07$~mg~N~L$^{-1}$~m$^{-1}$, with chloride
confirming denitrification rather than dilution
\parencite{Jaynes2014,Carstensen2021}. Across six saturated buffers, all
removed nitrate from tile-outlet water ($13$--$179$~kg~N~yr$^{-1}$ over
$3.4$--$40.5$~ha; $8$--$84\%$ effectiveness) at a mean cost of
US\$$2.94$~kg$^{-1}$~N, making them simple, cheap, and quick to install
\parencite{Jaynes2019}. Denitrification beds followed Michaelis--Menten
kinetics with very low greenhouse-gas emissions \parencite{Ghane2015}.
In-channel practices contribute additional reductions: a low-grade weir cut
nitrate load $58\%$ \parencite{Strock2025b}, a two-stage ditch reduced TP load
$40\%$ \parencite{Hodaj2017}, and blind inlets and subsurface fertiliser
placement reduced DRP $65$--$72\%$
\parencite{Kröger2013,Wilson2019,Sharpley2019}. Stacking practices multiplies
benefit: a combined drainage-management--bioreactor--saturated-buffer system
reduced nitrate export $78$--$88\%$ \parencite{Singh2024b}. A key limit is
that where tile drains route water directly to streams they bypass the
riparian zone, curtailing natural attenuation \parencite{Puckett2005}.

\subsection{Comparative effectiveness, cost, and watershed-scale integration}

\subsubsection{Load-reduction benchmarks and cost}

Comparative syntheses give consistent benchmarks: denitrifying bioreactors,
wetlands, cover crops, rotation, and nutrient management each reduced
subsurface nitrate loads $23$--$40\%$ \parencite{Yuan2022}, and nine BMPs in
the Lower Mississippi Alluvial Valley delivered $15$--$100\%$ nutrient
reductions \parencite{Kröger2012}. A US Corn Belt review of 119 records found
controlled drainage, water-quality wetlands, bioreactors, drainage water
recycling, and saturated buffers all effective for nitrate but with differing
co-benefits: wetlands gave the highest nitrate reduction plus P removal and
habitat value, drainage water recycling the greatest yield benefit, and
bioreactors effective nitrate removal but limited extra benefits
\parencite{Mitchell2023}. Global synthesis places bioreactor flow-season
removal at $14$--$66\%$, wetlands near $39\%$ TN, and buffers and controlled
drainage at $\sim30$--$50\%$ annual N-load reduction \parencite{Wang2025b}. On
cost, wetlands rank most cost-effective ($2.04$--$2.91$~\$~kg$^{-1}$~N),
followed by two-stage ditches and cover crops
\parencite{Roley2016,Messer2021}; bioreactor installation runs
$<$\$$5{,}000$--$27{,}000$ at $<$\$$2.50$--$\sim$\$$20$~kg$^{-1}$~N, with
corncob media, pumped systems, and spring bioreactors cheaper still
\parencite{Christianson2021,Law2023,Hartfiel2024,Easton2019}. For phosphorus,
edge-of-field practices retain P chiefly by particulate sedimentation
($27$--$100\%$) but depend strongly on residence time and can flip from sink
to source during events, so effective P control requires matching practice to
site hydrology, connectivity, and legacy P
\parencite{Mendes2020,Kröger2013,Singh2020c,King2015,Ni2020}.

\subsubsection{Scale attenuation, required reductions, and practice-stacking}

Because tile drainage is the primary pathway controlling watershed water and
nitrogen export, downstream targets translate to large required reductions at
the tile outlet---about $55\%$ for one Iowa river
\parencite{Schilling2012b,Ikenberry2014}. Concentrations decline with drainage
area as export dilutes downstream, and substantial in-network retention
($34$--$82\%$) occurs via underground denitrification, though in-stream
processing immediately below tile outlets offers little potential because its
timing is poorly matched to seasonal loading
\parencite{Schilling2012b,Billy2011,Billy2013}. This scale attenuation of
in-field effects means edge-of-field benefits measured at a single tile can be
diluted at the watershed outlet, and meeting large goals requires
\emph{combining} in-field and edge-of-field practices. In the Mississippi
basin, nitrogen-management practices alone were insufficient, but N-management
plus N-removal practices reached the $45\%$ target while converting under
$6\%$ of cropland when targeted \parencite{McLellan2015}, and LiDAR-based
placement achieved $34$--$39\%$ nitrate reduction while removing $<5\%$ of
cropland \parencite{Tomer2015,Kvítek2023}. Long-term records confirm that
stacked BMPs work, with Czech outlet nitrate declining over 2004--2018 and
combined practices cutting load $60.3\%$
\parencite{Fraindová2022,Law2022}. Crucially, traditional surface-water BMPs
can be inadequate over tile drainage: seven years of strip-till, grassed
waterways, and stream buffers produced no measurable change in nitrate,
phosphorus, sediment, or hydrology because they did not intercept subsurface
flow, which must therefore be a design priority \parencite{Lemke2011}.
Finally, monitoring design strongly affects inferred loads: high-frequency
sampling captures storm dynamics that grab sampling misses, and significant
nitrate declines are far more detectable in flow-weighted concentration than
in load or yield \parencite{Drake2018,Danalatos2022,Irvine2019}.

% ----- knowledge/08_other-water-quality-dimensions-pesticides-microb.tex -----
\section{Other Water-Quality Dimensions: Pesticides, Microbes, Trace Elements, and Greenhouse Gases}

Beyond nitrogen and phosphorus, subsurface drainage governs the field-to-stream
transport of pesticides, manure-borne microbes and resistance genes, trace
elements, and dissolved greenhouse gases, and it reshapes agricultural ecology.
Across these dimensions the evidence converges on a common logic: the same
preferential-flow pathways and hydrologic controls that dominate nutrient loss
also govern contaminant delivery, and management---application method, timing,
and edge-of-field engineering---frequently matters more than drainage per se.

\subsection{Herbicides, Pesticides, and Emerging Contaminants}

The timing of runoff and drainage relative to application is the single most
powerful control on peak pesticide concentrations and losses. Synthesising more
than thirty North American studies, \textcite{Kladivko2001} established that
pesticides are commonly detected in drainflow during the first drainage event(s)
after application, concentrations decline rapidly thereafter, and most mass
losses are very small---generally less than $0.5\%$ and often below $0.1\%$ of
applied---while tile concentrations and losses are usually much lower than in
surface runoff, frequently by an order of magnitude. Site-specific values
corroborate this: Illinois long-term atrazine losses were $0.11$--$0.22\%$ of
application \parencite{Kalita2006}; New York 2,4-D and atrazine losses were
$0.08\%$ and $0.02\%$ \parencite{Shalit1996}; and less than $1\%$ of applied
herbicide was recovered in tile drains over two years
\parencite{Fortin2002,Siimes2006}. Early events dominate: most atrazine and
metolachlor loss occurred within the first $61$ days \parencite{SOUTHWICK1990},
and over $80\%$ of catchment losses occurred during just two rain events
\parencite{Leu2004}. In cold regions most loss can occur during snowmelt and
thaw rather than immediately after application \parencite{Siimes2006}, and
two-phase seasonal behaviour is common \parencite{Novak2001,Dousset2004}.

The mechanism reconciling near-immediate arrival with small overall losses is
preferential (macropore) flow. Tracer breakthrough curves were typical of
preferential flow, with herbicides arriving after less than $1$~cm of net
drainage, interpreted as transport through only $\sim3\%$ of pore volume
\parencite{Fortin2002,Dousset2004}. Soil structure, rather than sorption or
degradation, dominated between-soil differences, with greater metolachlor losses
in a clayey Vertic Cambisol than a loamy Luvisol \parencite{Novak2001,Novak2003}.
Although mass losses are small, transient concentrations can exceed thresholds:
Ontario maxima reached $287$~$\mu$g~L$^{-1}$ (atrazine) in tile discharge, with
the atrazine MCL exceeded in $60\%$ of tile samples \parencite{GAYNOR1995}, and
Iowa subsurface drains had glyphosate maxima above $250$~$\mu$g~L$^{-1}$
\parencite{Coupe2012}. Degradates contribute meaningfully and can be more mobile
than parents: deethylatrazine accounted for $10$--$27\%$ of combined loss
\parencite{GAYNOR1995,MUIR1978,JAYACHANDRAN1994}. Management levers include
banding, residue-conserving tillage with intercrops (reducing loss
$46$--$49\%$), and application-timing restriction via a soil-wetness index
\parencite{Gaynor2001,Julien2023,Moorman1999}, whereas corn-stover harvest can
worsen drainage quality \parencite{Shipitalo2016}. Woodchip and biochar-amended
bioreactors removed atrazine (up to $93\%$) mainly abiotically without inhibiting
denitrification \parencite{Hassanpour2019,Ilhan2011}, and modelling frameworks
(PESTDRAIN, HYDRUS-2D, RZWQM, SWAT+) reproduce drainage pesticide dynamics when
macroporosity and sorption kinetics are represented
\parencite{Branger2009,Boivin2006,Wendell2024,Novak2003}. Tile drainage also
transports emerging contaminants---PPCPs, hormones, PBDEs, perfluoroalkyl acids,
and neonicotinoids---whose behaviour mirrors the timing and preferential-flow
patterns above, generally at low, transient concentrations
\parencite{Gottschall2012,Gottschall2013,Gottschall2017,Ghirardini2019,Schaafsma2019}.

\subsection{Manure-Borne Microbial Contaminants and Resistance Genes}

Manure and organic amendments introduce fecal microorganisms, resistance genes,
hormones, and trace elements that drainage can rapidly route to receiving waters,
with application method and soil moisture---rather than manure presence
alone---largely governing transport. Curli-producing \textit{E.\ coli} were far
more abundant in tile drains from farms that surface-applied manure than one that
incorporated it, indicating that application method itself shapes the transported
phenotype \parencite{Truhlar2015}. Tillage governs pathogen delivery mainly
through macroporosity: zero tillage generally produced higher subsurface flow and
greater \textit{E.\ coli} loads than conventional tillage
\parencite{Thiagarajan2007,Jenkins2009}. Antecedent soil moisture and water-table
position exert first-order control: effluent on wet soil produced drain-flow
\textit{E.\ coli} of $10^{5}$--$10^{7}$ per~100~ml versus $\sim10^{4}$ under a
greater moisture deficit, and where preferential flow occurred, drainage
concentrations approached those of the applied effluent
\parencite{Collins2007,Ding2021}. Manure is also a reservoir of antibiotic
resistance genes that persist through drained soils into drainage water: shotgun
metagenomics identified $320$ unique ARGs, with manure treatment shifting soil
and drainage-water resistomes and many manure-derived genes associated with
mobile genetic elements in food- and water-borne pathogens \parencite{Smith2019b}.
Steroid hormones were detected only transiently and generally did not differ
significantly among treatments \parencite{Jenkins2009}. Engineered receiving
features can attenuate microbial loads---a constructed wetland reduced nitrate and
\textit{E.\ coli} mass loads $67.6\%$ and $63.3\%$ \parencite{Haverstock2017}---
though some wetlands act as net \textit{E.\ coli} exporters \parencite{Stott2023}.
Drained soils, bioreactors, and wetlands host distinct, dynamic microbial
communities structured by depth and season, and favourable conditions can be
engineered to develop effective denitrifier assemblages
\parencite{Andrus2014,Porter2015,Hathaway2015,Knorr2023,Schilling2017}.

\subsection{Trace Elements, Selenium, and Acid Sulphate Hazards}

Where drainage concentrates trace elements, it can create acute ecological
hazards. The Kesterson Reservoir remains the definitive case: waterborne selenium
averaged $0.122$~mg~L$^{-1}$ ($\sim230$ times reference), biomagnified through the
food chain, and caused dead and deformed bird embryos, with integrated studies
fulfilling the equivalent of Koch's postulates and identifying selenium as the
sole causative agent \parencite{SAIKI1987,Ohlendorf2002,BENSON1993}. The
regulatory legacy is durable: open-water evaporation is now prohibited in
California, requiring rock-media systems \parencite{Ayars2014,BRADFORD1991}, and
selenium-removal wetlands retained about $59\%$ of inflow Se \parencite{Gao2003}.
In evaporation basins, evapoconcentration accumulates extreme salinity and trace
elements \parencite{WESTCOT1988}, and redox chemistry jointly controls arsenic
and uranium mobility \parencite{Ryu2010,Gao2007,Duff1997}; boron isotopes and
other tracers can distinguish contaminant sources
\parencite{Tayfur2010,Komor1997,Zhou2023}. At the opposite extreme, intensive
drainage can drive negative micronutrient balances and net soil depletion
\parencite{Gul2013}, while central-Iowa drainage ionic composition was dominated
by bicarbonate, calcium, and magnesium \parencite{Zimmerman2017b}. Drainage of
acid sulphate soils is a distinct hazard: oxidation of pyritic sediments drives
stream pH declines to near or below $4.0$ and elevates aluminium and metals, with
controlled pipe drainage recommended over in-stream liming
\parencite{Osterholm2005,Åström2005,Virtanen2017,Salo2023}.

\subsection{Greenhouse-Gas Dimensions}

Drainage nitrate is coupled to greenhouse-gas fluxes, and drainage water is
commonly supersaturated with dissolved N$_2$O, CH$_4$, and CO$_2$. In Japan,
indirect emissions via subsurface drainage were equivalent to $4.6\%$, $58.3\%$,
and $2.5\%$ of direct emissions, with N$_2$O-N/NO$_3$-N ratios below the IPCC
factor \parencite{Sawamoto2003,Minamikawa2010,Minamikawa2013}. In Canada,
dissolved N$_2$O losses approached surface emissions, were strongly seasonal, and
declined along tile length, implying an in-drain sink and likely underestimation
of the IPCC coefficient \parencite{Burton2024}; a global synthesis found
subsurface drains had the highest mean emission factor among artificial waters,
with the stationary IPCC factor tending to overestimate N$_2$O production
\parencite{Webb2021}. Conversely, installing drainage sharply reduces methane in
paddy systems by improving aeration: drained fields emitted $71\%$ less CH$_4$
than non-drained fields \parencite{Shiratori2007,Furukawa2008}. Management
modulates N$_2$O: controlled drainage kept cumulative N$_2$O below $1\%$ of
applied N, and saturated riparian buffers removed nitrate without increasing
N$_2$O \parencite{Smith2019,Davis2019,Nash2015e,Carstensen2021}. Carbon export
from drained fields is dominated by inorganic carbon through tiles ($\sim96\%$ of
total carbon export), so reducing water throughput supports carbon retention
\parencite{Schilling2023,Driscoll2025}.

\subsection{Ecological Responses to Subsurface Drainage}

Drainage reshapes agricultural ecology from soil fauna upward. \textit{Lumbricus
terrestris} abundance was typically higher above drains than between them (median
counts twice as high, fresh mass five times), so drainage pattern affects
population density and should inform sampling design
\parencite{Nuutinen2001,Nuutinen2011}. Drainage type is associated with
farmland-bird outcomes, with skylark and meadow pipit favouring open-drained over
subsurface-drained fields \parencite{Marja2013}, while tile-drainage wetlands
carried cumulative pesticide concentrations $4.7$ times higher than reference
wetlands (neonicotinoid-dominated), with measurable neurological effects in
juvenile amphibians \parencite{Campbell2023}. Subsurface drainage of meadows was
a leading cause of decline in semi-aquatic snails \parencite{Dreyfuss2018}.
Restored aquatic habitats can benefit: in restored Iowa oxbows, tile drainage had
negligible effects on physical or biotic variables, providing promising evidence
that multipurpose oxbows can enhance biodiversity without ecological penalty, with
outcomes improved by sediment excavation
\parencite{Leberg2024,Winikoff2020,Harris2018}. Crucially, management practices
often exert clearer effects than tile drainage itself: no-till and diverse
rotations altered bulk density, organic matter, and aggregate stability while
cover crops and drainage had minimal effect in some settings
\parencite{Saha2024,Hermawan1996,Sherbine2023}, and additional practices such as
polyacrylamide amendment and anti-weed mesh with drainage provide further
quantified nutrient-retention and production benefits
\parencite{Jiang2010,GiménezValero2024,GiménezValero2025}.

% ----- knowledge/09_agronomic-benefits-yield-salinity-reclamation-an.tex -----
\section{Agronomic Benefits: Yield, Salinity Reclamation, and Socio-Economics}

Beyond its hydrologic and water-quality functions, subsurface drainage is a
production technology whose agronomic and economic value is well established.
The evidence---multi-year paired-field trials, long-term monitoring, and
calibrated simulation across humid, arid, monsoonal, and cold-climate
settings---supports four consolidated conclusions: drainage and its managed
variants raise crop yields, most strongly on poorly drained soils and in wet
years, while maintaining or improving water-use efficiency; subsurface systems
reliably leach salts and reclaim saline and waterlogged land; saline and
brackish water can be reused safely where salt balance is maintained; and
drainage investment delivers favourable financial and social returns. Meta-
analyses anchor the headline magnitudes: a global synthesis of 891 experiments
confirmed drainage significantly reduces soil salinity and raises yield
\parencite{Yang2025}, and a U.S. meta-analysis of 154 comparisons found
controlled drainage and subirrigation raised yield $8.0\%$ over free drainage
\parencite{Singh2022b}. Spacing, depth, and outlet-control mechanisms are
developed in the design section; here we own the yield, reclamation, reuse, and
socio-economic outcomes, together with the salinity-specific design optima and
salt-balance dynamics.

\subsection{Crop yield and agronomic outcomes}

\subsubsection{Magnitude of yield gains}

Direct drained-versus-undrained comparisons repeatedly document substantial
yield increases, primarily by relieving early-season waterlogging. Artificial
drainage raised soybean yield $275$~kg~ha$^{-1}$ ($8\%$) in experiments and
$157$~kg~ha$^{-1}$ ($4\%$) in producer fields \parencite{Mourtzinis2021}, and in
a long-term Indiana experiment corn yields rose progressively in all drained
treatments across 24 corn years while the undrained control stayed flat, though
soybean was unaffected by spacing \parencite{Rui2024,Kladivko2005}. The largest
relative gains occur in heavily stressed environments: Nile-Delta drainage
doubled cotton and rice yields and raised wheat and clover $50\%$
\parencite{ELMOWELHI1988}; Indian acid saline ``kari'' soils gained
$1.36$~t~ha$^{-1}$ ($61\%$) \parencite{Mathew2001}; Bangladeshi waterlogged
saline clay gained $30$--$64\%$ \parencite{Islam2025,Islam2022}; and an Alberta
system raised barley grain yield from $1310$ to $3900$~kg~ha$^{-1}$
\parencite{BENNETT1982}. More moderate but robust gains dominate temperate
systems: shallow drainage in North Carolina raised corn $4\%$ and soybean
$14.3\%$ \parencite{Watkins2024}; Iranian paddies exceeded controls by
$13.4$--$31.3\%$ (up to $66$--$70\%$) \parencite{DarziNaftchali2017b,DarziNaftchali2014};
and Chinese blind-ditch and pipe drainage raised rice $8.0\%$ and $6.0\%$
\parencite{Qi2025b}.

\subsubsection{Controlled drainage, subirrigation, and recycling}

Managed practices deliver benefits beyond free drainage when they retain water
for crop use in dry periods---quantified at $8.0\%$ average yield gain across
154 comparisons \parencite{Singh2022b}. In Missouri, free drainage and drainage
water recycling raised corn yield $20\%$ and $26\%$ over a nondrained control
(recycling reaching $48\%$ in some years), avoiding drought losses
\parencite{Kaur2021}; in North Carolina recycling raised corn and soybean $8\%$
and $17\%$ \parencite{Moursi2024}. Subirrigation is strongly site-dependent,
governed by lateral hydraulic conductivity: across three Ohio sites it improved
relative yield to $80$--$93\%$ at two sites but had almost no effect where
conductivity was very low, and simulations projected it increasingly
outperforming free drainage under future climate \parencite{Gunn2018,Gunn2019}.
Controlled drainage and subirrigation kept soil water optimal in both wet and
dry years \parencite{Niaghi2019}, improved claypan soybean yield while reducing
interannual variability and allowing earlier planting \parencite{Nelson2011},
and combined with foliar pyraclostrobin produced synergistic soybean gains up
to $1070$~kg~ha$^{-1}$ \parencite{Nelson2011b}.

\subsubsection{Mechanisms, water-use efficiency, and boundary conditions}

The proximate mechanism is relief of excess-water stress and restoration of
root-zone aeration \parencite{BENNETT1982,DarziNaftchali2014}, with growth-stage
sensitivity (cotton flowering and boll formation most sensitive; surface
waterlogging worse than subsurface) \parencite{Qian2021}. Improved aeration
propagates through physiology---raising emergence, root vigour, chlorophyll, and
dry matter in cotton \parencite{Yang2022}, lowering leaf Na$^+$/K$^+$ in
sunflower \parencite{Islam2022}, and raising rice nitrogen-use efficiency
$16.1\%$ \parencite{DarziNaftchali2022b}. Controlled drainage repeatedly
improves irrigation water-use efficiency without yield loss: simple management
of existing Egyptian systems improved WUE $15$--$20\%$ \parencite{Wahba2005},
controlled drainage cut applied irrigation $14$--$16\%$ \parencite{Mahmoud2021},
and Chinese buried pipes raised water productivity, highest near $1.4$~m burial
\parencite{Liu2025}. Yield response is soil- and weather-specific: only
undrained controls show negative rainfall responses, wet years favour narrower
spacing and dry years wider \parencite{Rui2024,Kladivko2005,Acharya2019}, and
drainage advanced sowing $7$--$8$ days and avoided $29$--$52\%$ harvest
population losses in wet years \parencite{Mourtzinis2021,Nelson2011}. The
evidence also delineates limits: where surface drainage is adequate, subsurface
drainage did not raise grain yield \parencite{Wiersma2010}; medium-textured
soils are most responsive while fine and coarse show no clear effect
\parencite{Singh2022b}; over-drainage can induce drought
\parencite{Cordeiro2015,DeBoer2006}; an optimum water-table depth exists that
varies with soil, crop, and permeability \parencite{HOOGHOUDT1952,Bigelow2001};
and installation without maintenance underperforms \parencite{Malota2016}.
Adoption is a leading adaptive practice---about three-fourths of surveyed U.S.
producers report drained cropland \parencite{Morton2015}, and nearly $90\%$ of
Finnish farmers rate subsurface drainage important, ahead of irrigation
\parencite{PeltonenSainio2021}.

\subsection{Salt leaching and reclamation of saline and waterlogged soils}

\subsubsection{Reclamation efficacy}

Subsurface drainage reliably lowers the water table below the root zone and
removes salt from the upper profile, converting saline and waterlogged land into
productive cropland, with topsoil salinity reductions commonly $35$--$80\%$ and
corresponding yield gains, validated at production scale
\parencite{Yang2025,Derby2024}. In India, drained-block salinity fell from
$7.18$ to $4.62$~dS~m$^{-1}$ (mean $35.7\%$) while undrained soils rose, and
yields increased $18.75$--$27.62\%$ \parencite{Sharma2006}; wide-spacing
reclamation lowered topsoil salinity from $8.4$ to about $2.1$~dS~m$^{-1}$ and
raised cropping intensity from $143\%$ to $191\%$ \parencite{Manjunatha2004}.
Chinese reclamations are similarly large: drip-plus-pipe cut EC from about
$15$ to $2.9$--$6.9$~dS~m$^{-1}$ and raised seed-cotton yield $52\%$
\parencite{Heng2021}, and after four years the $0$--$1.0$~m layer fell from
$27.6$ to $4.1$--$7.2$~dS~m$^{-1}$ \parencite{Zhang2021}. Desalinization
is surface-concentrated---$0$--$20$~cm rates exceed $80\%$ while deeper layers
lag and the meta-analytic reduction is nonsignificant below $80$~cm
\parencite{Han2023,Yang2025}. Comparable reductions are reported for Turkey,
Tunisia, and Canada, the last raising barley yield up to twentyfold
\parencite{Bahçeci2009,HajAmor2019,Steppuhn2006,Steppuhn2017}, and reclamation
homogenises the field and improves porosity and infiltration
\parencite{Yu2015,Yu2025,Sharma2000}. Where the water table stays shallow,
however, combined waterlogging and salinity are more damaging than waterlogging
alone, with yields failing above $12$~dS~m$^{-1}$ \parencite{Kahlown2002}.

\subsubsection{Design parameters governing salt removal}

Closer spacing and greater depth shorten the desalinization period and increase
salt removal, though depth benefits saturate. A DRAINMOD analysis found salinity
fell $55.5\%$ in two years at $20$~m spacing but only $5.9\%$ in five years at
$200$~m, a slightly-saline target within a decade requiring spacing below $30$~m
\parencite{Chen2025}; the meta-analysis generalises this into climate-specific
layouts (spacing $>30$~m and depth $>1.2$~m in humid climates; $10$--$20$~m
spacing in arid climates with shallow water tables) \parencite{Yang2025}. Field
and column studies identify optima near $15$--$20$~m
\parencite{Wang2020,Han2023,Chen2022}, with desalination decaying beyond about
$30$~m \parencite{Zhu2025b}. Depth effects are pronounced but bounded---pipes
shallower than $1.0$~m are essentially invalid for discharging salt
\parencite{Qian2021b}, yet returns diminish beyond about $1.4$~m
\parencite{Bahçeci2007,Bahçeci2006}, and an intermediate $0.8$--$1.2$~m depth is
often preferred under saline irrigation, saving $10$--$15\%$ water
\parencite{Feng2017,Muhammad2021}. Envelope and filter design is a first-order
control: the envelope factor dominates pipe material, a double-layer geotextile
achieving $>95\%$ salt removal while triple layers clog
\parencite{Li2022,Wang2024d,Elzoghby2021}. Economic analyses converge on roughly
$20$~m spacing at $1.2$~m depth as the maximum-return layout \parencite{Chen2025}.

\subsubsection{Salt balance, temporal decline, and leaching structure}

Because tile drains intercept salt-laden water, effluent salinity is
characteristically higher than irrigation-water salinity and declines over years
as stored salt is depleted \parencite{Li2022,Karegoudar2022}, with ECdw/ECiw
ratios of $11$--$17$ in Iran \parencite{Akram2013}. Salt-balance accounting is
the standard diagnostic: an oasis exported about $94\%$ of imported salt while
discharging $42\%$ of irrigation water \parencite{Askri2022}, though clogging can
cut removal to $15$--$27\%$ \parencite{BenAissa2013,HajAmor2019}. Most residual
salt is removed within ten years, effluent-salinity--drain-age relationships
being stable across decades \parencite{JOHNSTON1993}, though modeling indicates
$15$--$50$ years to full equilibrium, with drainage-water salinity more sensitive
to leaching fraction than to spacing \parencite{Kelleners2000}. Salt export
increases with both depth and spacing, depth having the larger effect where
groundwater quality declines with depth \parencite{GRISMER1993,KAMRA1991,KAMRA1994}.
Leaching is vertically stratified (upper $>$ middle $>$ lower) and ionically
selective, with Cl and Na leached most efficiently ($91\%$ and $88\%$ at
$0$--$20$~cm) and Ca persisting as precipitated gypsum
\parencite{Chen2022,Qian2025,Yu2016}. Re-accumulation is common, following a
``wash-out--recovery--accumulation'' pattern with surface re-salinization in dry
periods \parencite{Zhu2025b,Dou2022,Shi2021b}. Mechanistically, under complete
ponding flow concentrates near the drain ($\sim80\%$ of flow through $20\%$ of
soil volume), so progressive or partial flooding, low-permeability barriers, and
surface mulches raise efficiency substantially (a $10$~m barrier saved
$52$--$82\%$ of leaching water; a mulch cut effluent volume $81\%$)
\parencite{RAO1991,Youngs2000,Yu2018,Zhang2022}, whereas macropore bypass and
fine texture prolong leaching \parencite{Xin2016,MANN1982}.

\subsection{Saline and brackish water reuse}

Subsurface drainage enables safe reuse of saline and brackish water where salt
balance is maintained by adequate leaching. Blended drainage water up to
$9$~dS~m$^{-1}$ produced about $90\%$ wheat yield in India, cyclic use generally
outperforming continuous use \parencite{Sharma2004}; five-year data gave relative
wheat yields of $95.8$, $90.3$, $83.7$, and $77.8\%$ at $6$, $9$, $12$, and
$18.8$~dS~m$^{-1}$ with root-zone salinity kept below $4$~dS~m$^{-1}$
\parencite{Sharma1998}. In China, saline ($6$~g~L$^{-1}$)
irrigation matched freshwater wheat yields with continued salinity decline
\parencite{Gao2023}, and drainage kept saline irrigation viable up to
$6.25$~dS~m$^{-1}$ for maize \parencite{Feng2017} and $8.0$~dS~m$^{-1}$
under cut-soiler drainage \parencite{Neha2022}, with cotton optima at
$2$--$4$~g~L$^{-1}$ \parencite{Li2024,Lu2019}. Reuse nonetheless carries sodicity
and accumulation risks: cyclic reuse caused profile salinity fluctuations and
sodic surface soils (SARe $>11$) \parencite{Mitchell2000,Jayawardane2001},
though modeled reuse over subsurface drains caused only marginal groundwater
deterioration \parencite{KARAJEH1994c}. Integration with complementary practices
raises performance: combined pipe-plus-ditch systems reached $170\%$ of a single
ditch's efficiency \parencite{Yao2024}; timing drainage to snowmelt more than
doubled surface salt leaching \parencite{Yu2016}; cut-soiler shallow drainage cut
salinity $50$--$60\%$ \parencite{Neha2022,Neha2022b}; and regional coupling of
optimised drip irrigation with tailored drainage reduced district-wide irrigation
$10\%$ and salinity $15\%$ \parencite{Liu2026}. At basin scale, persistent
waterlogging and salinity in Central and South Asia reflect deteriorated
infrastructure and inadequate effluent disposal, underscoring that drainage must
be managed within an integrated salt balance
\parencite{Dukhovny2007,Qureshi2008,Ritzema2016}.

\subsection{Socio-economic returns and investment viability}

\subsubsection{Yield, income, and productivity gains}

Drainage investment generates multi-dimensional returns---higher yields, cropping
intensity, farm and household income, land value, and employment
\parencite{Gopalakrishnan2007,AbdelDayem2007}. Drainage typically raises average
yields $10$--$30\%$ and lifts net farm income where leaching water and
sustainable disposal are assured \parencite{AbdelDayem2007,Datta2000,Wichelns2002}.
In the Nile Delta, net incomes rose by $180$--$375$ per hectare with yield gains
accruing progressively from $50\%$ in year one to $100\%$ by year four
\parencite{Wichelns2002}. In India, the Uttar Pradesh project reclaimed over
$175{,}000$~ha of sodic land, raising farm income over $150\%$, cropping
intensity from $49\%$ to over $200\%$, and household income $53\%$
\parencite{Gopalakrishnan2007}, and Haryana pilots raised cropping intensity and
wheat yields as salinity fell from about $50$ to $5$~dS~m$^{-1}$
\parencite{Datta2002,Datta2000}. Long-term drainage also improves the soil
resource base, raising the soil quality index $76$--$136\%$ with carbon-
sequestration co-benefits \parencite{Mukhopadhyay2023,Mukhopadhyay2023b}, and in
Pakistan a project raised cropping intensity from $102\%$ to $157\%$ and living
standards $30\%$ \parencite{Ghumman2011}.

\subsubsection{Financial appraisal, scale, and cost-reducing design}

Formal appraisals confirm favourable performance. The RAJAD project
($14{,}925$~ha) reported a benefit--cost ratio of $1.28$--$2.87$, payback of
$4.5$--$6.9$~years, and an internal rate of return of $18$--$35\%$
\parencite{Gopalakrishnan2007}; Nile-Delta drainage gave farm-level benefit--cost
ratios of $10.3$ (non-saline) and $5.5$ (saline) \parencite{Wichelns2002}, and
small-scale manually installed drainage yielded IRRs of $8$--$11.5\%$ where water
for leaching and disposal are available \parencite{Datta2000,Datta2002}, though
not all appraisals are high \parencite{Ghumman2011}.
National programmes reflect this importance: Egyptian drainage serves about
$2.4$~million hectares \parencite{AbdelDayem2007}, India reclaimed about
$1.39$~million ha \parencite{Gopalakrishnan2007}, and globally $12$--$17\%$ of
cultivated area is drained \parencite{Valipour2014}. Careful design
sharply reduces cost: Nile-Delta design charts gave minimum network costs of
$205$--$374$ per hectare at $300$~m lateral length ($37\%$ below $100$~m)
\parencite{Gabr2024}; corrugated plastic pipe and factory pre-wrapping of
envelopes reduced installation costs by roughly $50\%$ \parencite{Yannopoulos2020};
and mechanization with laser grade control and improved envelopes lowered costs
and raised quality \parencite{Ritzema2006,Ritzema2007,TeKrony2004}. Disposal
infrastructure follows economies of scale, with economic outcomes strongly
conditioned by disposal costs \parencite{Singh2001,Bethune2008,HANSON1994,Napierala2024}.

\subsubsection{Policy and externality context}

Because drainage water carries disposal costs and externalities, its economics
interact with policy. In California an unregulated regime produced applied-water
and drainage volumes above efficient levels, whereas a direct drainage fee
induced efficiency (net social benefit $1132$~US\$/ha)
\parencite{DINAR1989}. Experience underscores that economics alone do not govern
outcomes: management must be technically sound, economically viable, and socially
acceptable \parencite{Letey1999}, and basin-wide accounting shows combination
strategies can cut soil salinity and river salt loading while raising yields
\parencite{Burkhalter2006,Quinn2009,McDowell2021}. Because artificial drainage
shifts diffuse phosphorus toward the subsurface pathway
\parencite{Lucci2014,Allaire2011,Dadfar2010}, and because deeper, wider drains
raise drain-water salinity and salt load, drainage water quality is increasingly
advanced as a necessary design criterion alongside production, with active
water-table management and decision support the key future directions
\parencite{Ayars2015,Yannopoulos2020,Ritzema2023}.

% ----- knowledge/10_modeling-measurement-and-mapping-of-drainage-sys.tex -----
\section{Modeling, Measurement, and Mapping of Drainage Systems}

The physical findings summarised in the preceding sections rest on a mature
toolkit for simulating, measuring, and locating drainage systems. Across
site-scale validations, benchmark comparisons of analytical solutions,
pipe-scale experiments, watershed representations, formal calibration and
uncertainty analyses, and geospatial detection, a robust set of consolidated
findings has emerged: process models reproduce drainage flow and nitrate
reliably after calibration but with skill that improves strongly with temporal
aggregation; treating drainage as purely saturated flow over-predicts
discharge; saturated conductivity and evapotranspiration dominate model
sensitivity; data-driven and remote-sensing methods now complement process
models; and no single geophysical or imaging technique is universally best for
mapping buried networks. This section consolidates that methodological
evidence so the physical and management conclusions can stand without model
mechanics.

\subsection{Field-scale drainage and water-quality models}

\textsc{DRAINMOD} is by far the most extensively tested field-scale
water-management model and is repeatedly judged a reliable design and
evaluation tool for poorly and artificially drained lands, its process-based
structure with independently measurable inputs describing drainage hydrology
across many climates when soil hydraulic properties are adequately
characterised \parencite{Singh2019b,Gowda2012}. A defining, cross-cutting
result is that skill is scale-dependent in time: water-table depth and
drainage are reproduced with accuracy improving markedly from daily to monthly
to annual scale. In North Carolina, water-table and daily-drainage
Nash--Sutcliffe efficiencies were 0.49--0.73, monthly drainage 0.77--0.87, and
annual volumes within 5\% of observations \parencite{Gowda2012}; Iowa and
eastern-Canada cumulative drainage fell within 2--10\% and about 30~mm of
measurements \parencite{Singh2006b,Dayyani2009}. The persistent weak point is
daily peak flow, which is systematically underestimated
\parencite{Morrison2014,Singh1996}, and uncertainty widths grow from daily to
yearly scale while annual predictions remain within bounds \parencite{Wang2006}.
The model has been validated well beyond its humid origin---cold Canadian
snowmelt ($R^2$ up to 0.93), semi-arid India, and Illinois claypan soils
\parencite{Dayyani2009,GUPTA1993,MOSTAGHIMI1989,Abduljaleel2023,Awad2022b}---but
delineated failure modes include unaccounted deep and lateral seepage
\parencite{Sinai2006}, absent time-variable depression storage
\parencite{Abduljaleel2023}, and its parallel-drain geometry requiring
plot-by-plot calibration in irregular landscapes \parencite{He2002}. Where
site data are scarce, pedotransfer estimation (ROSETTA, van Genuchten fits)
with automated calibration (PEST) drives it reliably
\parencite{Singh2006b,Qi2015,Malota2024,Malota2015}, and enhancements target
the daily-peak weakness through macropore flow (capturing $\sim$75\% of peak
events) and neural-network surrogates \parencite{Askar2020,Awad2022,Awad2022b};
the evapotranspiration algorithm is a pronounced sensitivity, the best ET
method not necessarily best for water-table prediction
\parencite{Malakshahi2020}.

Richards-equation and comparative codes generally perform comparably.
HYDRUS~(2D/3D) predicts fluxes and water tables closely (EF 0.91--1) and
clarifies mechanisms such as conductivity control of drawdown in compacted
soils \parencite{Karandish2018,DarziNaftchali2018,Mante2018,Guedessou2021,Ebrahimian2015};
CSUID, WAVE, SWAP, \textsc{WaSim}, CATHY, RZWQM, WEPP, ADAPT, and the
few-parameter \textsc{SIDRA} family have all calibrated satisfactorily, with
\textsc{DRAINMOD} often giving the most accurate discharge among them
\parencite{Alzraiee2013b,Alzraiee2013,ElSadek2001,Prasher1996,Verma2014,Malota2022,Muma2017,Malone2014,Xian2017,RevueltaAcosta2021,Mendez2004,Zimmer2023,Henine2022,LESAFFRE1988}.

For water quality the same pattern recurs across the model families
(\textsc{DRAINMOD} suite, RZWQM/RZWQM2, SWAT/APEX, APSIM, DNDC, HYDRUS): after
calibration they reproduce annual and cumulative drainage and nitrate loss
well---frequently within $\pm5$--$20\%$---and capture the direction of
management effects, but predictive skill falls with temporal disaggregation,
nitrate is predicted less accurately than flow, and daily nitrate
concentrations, crop N uptake, and residual soil N are recurrent weak points
\parencite{Youssef2018,Negm2014,Singh2020b,Dayyani2010,Ma2001,Liang2018,Ikenberry2017,Marjerison2016,Ojeda2018,Qi2025,Ma2007,Ahuja2007}.
\textsc{DRAINMOD-N~II} and RZWQM2 both satisfactorily predict tile nitrate,
and \textsc{DRAINMOD-DSSAT} predicts drainage, nitrate, and yield very well
\parencite{Du2017,Qi2012,Negm2017,Singh2022c}; cropping-system models
(DSSAT, RZWQM, DNDC, APSIM, SWAT) likewise reproduce cover-crop and
rotation effects on nitrate loss \parencite{Jeong2018,Ahmed2007b,Ale2012b,Pasley2021,Motarjemi2021}.
Phosphorus extensions are advancing but remain limited: RZWQM2-P first
simulated daily dissolved reactive P and its pathway partitioning, while
\textsc{DRAINMOD-P} predicted discharge and total P but poorly predicted DRP,
and few standard models yet represent tile P
\parencite{Sadhukhan2019,Grewal2025,BinShokrana2023b,Askar2021,Askar2025,Radcliffe2015,Wang2020b}.
Cold-season snowmelt and freeze--thaw errors motivated dedicated developments
(\textsc{DRAINMOD}~5.1, DRAIN-WARMF, RZ-SHAW), and calibration studies
repeatedly identify lateral saturated conductivity in the tile layer and field
capacity as controlling predictions
\parencite{Wang2006c,Dayyani2010b,Li2024b,Ma2007c,Malone2017}.

\subsection{Flow physics: saturated and unsaturated flow and preferential-flow representation}

A central mechanistic result is that treating drainage as purely saturated
flow overstates discharge. A saturated-flow model over-predicted seepage and
drainage time by roughly an order of magnitude relative to experiment, an
analytical solution, and the SUTRA Richards-equation model, because a vadose
zone retaining substantial water persists even at steady state
\parencite{Xin2016b}. Saturated--unsaturated drainage is governed by
gravitational, pressure, capillary, and air-pressure potentials, so saturated
conductivity alone is an inadequate design basis \parencite{Li2021}. Field
diagnosis in Japanese paddies showed that vertical, not horizontal, resistance
limits the maximum drainage rate \parencite{Nishida2020}, accounting for
vertical head loss improved water-table prediction \parencite{Talukolaee2022},
and retaining unsaturated-zone buffering matched numerical models closely
whereas ignoring it distorted transients \parencite{VandeCraats2021};
equivalent-medium representations reproduce integrated outflow well but can
misrepresent internal dynamics where fast vertical flow occurs
\parencite{Carlier2007}. At field-to-hillslope scale the representation of
macropores is decisive for reproducing rapid drainage. Dual-permeability and
dual-porosity models reproduce hydrographs well ($R^2=0.84$; groundwater
$r=0.98$) while single-porosity models reach acceptable accuracy with far
fewer parameters, though all fail to match low dry-season discharge
\parencite{Varvaris2018}; dual-permeability formulations enable low matrix
conductivity while still reproducing fast events
\parencite{Warsta2014,Salo2017,Koivusalo1999}, adding macropore flow improves
peak simulation without changing total drainage \parencite{Zhou2013}, and
bimodal retention functions predict preferential flow far better than unimodal
ones \parencite{Mohanty1997}. Fully coupled surface--subsurface--tile models
are recommended for detailed dynamics, with equivalent-porous-medium or
seepage-node simplifications reproducing outflow almost exactly for watershed
integration \parencite{DeSchepper2015,Boico2022,Xian2017}.

\subsection{Analytical drain-spacing solutions and pipe-scale hydraulics}

Unsteady-state (transient) spacing equations yield markedly narrower spacings
than steady-state equations and are preferred under variable rainfall---about
30~m versus 55~m in Iranian paddies---with the appropriate transient model
depending on rainfall duration and design criterion
\parencite{OOSTERBAAN1988,DarziNaftchally2014,Moshayedi2020,Pali2014,Nédélec2010}.
Incorporating evapotranspiration and unsaturated-zone flow generally permits
wider, more economical spacing: ET-aware solutions allow 9--18\% increases,
bare-soil evaporation 23--74\%, and unsaturated horizontal flow up to 30\% for
coarse soils near a shallow impervious layer
\parencite{NIKAM1992,Singh2010,Upadhyaya2000,AbdelFattah2005,Yousfi2014}.
Bi-level and sloping-aquifer solutions show asymmetric water tables and
slightly to moderately larger spacings, and dual-depth systems always remove
water more effectively than single-depth ones
\parencite{Verma1998,Kankarej2023,Upadhyaya2013,Hartani2001,Kirkham1997}.
Analytical solutions for multi-layered and ponded systems have been verified
against experiments and MODFLOW, establishing strong three-dimensionality and
sensitivity to anisotropy near drains, while flat-land theory remains
applicable on slopes up to about $0.15$~m~m$^{-1}$
\parencite{Chakrabarti2020,Sarmah2015,SHARMA1991,Afruzi2014,Lin2025,CHU1977,FIPPS1989,LORRE1994};
no single linear-diffusion model, however, can describe all stages of drawdown
\parencite{Kacimov2000}. The most complete descriptions solve the Richards
equation, but corrected Boussinesq approximations remain useful for design when
a variable drainable porosity or effective specific yield is used
\parencite{TANG1977,Xin2016c,VANDELEUR1959,SINGH1980,Bouarfa2000}. Single-tile
solutions using a variable region of influence eliminate the spacing parameter
for irregular systems \parencite{Cooke2001}, and space-fractional and inverse
formulations improve fits in heterogeneous soils and circumvent the awkward
drainable-porosity parameter
\parencite{FathAllah2025,Mehdinejadiani2013,Mehdinejadiani2020,Mehdinejadiani2022,Samani2007,Chahar2010,SINGH1992,KUMAR1994,Onyejekwe1999,Bahceci2008}.

At the pipe scale, entrance resistance and effective radius are now quantified
across perforation and envelope geometries: 3--4-row pipes averaged 0.6~cm
effective radius and 8-row reached 1.9~cm, while a knitted sock raised it to
5.7--6.0~cm and nearly eliminated entrance head loss, increasing perforation
rows being the most efficient improvement
\parencite{Ghane2022b,Ghane2022}. Numerical and experimental work shows
rectangular slots are more hydraulically advantageous and structurally safer
than circular holes, and screen envelopes reduce resistance and prevent rapid
clogging
\parencite{Gaj2020,Oyarce2017,Gaj2021,Stuyt2006}. Within-conduit flow is a
spatially varied flow problem in which Manning's equation alone is inadequate,
local head loss is non-negligible, and resistance rises 40--50\% over
gradually varied flow
\parencite{Bayat2011,Graber2004,Kouchakzadeh2011,Nédélec2008,Wu2022,Gupta2019}.

\subsection{Watershed representation, sensitivity, calibration, uncertainty, and evaluation}

How drains and subsurface heterogeneity are conceptualised strongly controls
simulated catchment behaviour: two grid conceptualisations of a Danish
catchment produced drain contributions of 13\% versus 31\% despite both
matching total streamflow within 0.1\% \parencite{Karan2021}. Seepage-node
representations are faster and better suited to unknown drain locations
\parencite{Boico2022}, distributed HydroGeoSphere reproduced flashy heads
better than MIKE SHE \parencite{DeSchepper2017}, machine learning links
streamflow signatures to simulated drain fraction \parencite{Schneider2022},
embedding drainage in integrated models reproduced streamflow without
site-specific calibration \parencite{Rathore2024,Muma2016}, and equivalent-medium
representations match explicit tile simulations within 2\%
\parencite{Thomas2016}; standard SWAT under-represents subsurface drainage
without explicit adaptation \parencite{Manguerra1998,Golmohammadi2016,Rahman2014}.
For regional assessment, MODFLOW matched observed heads closely, and new
Hooghoudt-based subsurface-drainage packages markedly outperform the
conventional DRN package (RRMSE $<5\%$) by calibrating only conductivity
\parencite{Mirlas2012,Ghorbani2016,Yang2022b,Han2025,Qian2025b}.

Formal analyses converge on saturated hydraulic conductivity and
evapotranspiration estimation as the dominant sources of sensitivity and
uncertainty, with van Genuchten shape parameters, drain spacing, and depth
also prominent, and parameter uncertainty propagating into predicted salinity
and yield
\parencite{NabiSichani2012,Malakshahi2020,Wang2006,Sedaghatdoost2018,Timalsina2023,Alzraiee2013}.
Objective-function design strongly affects identifiability: genetic algorithms
generally outperform PEST, combining discharge with water-table depth already
gives excellent results, and variational data assimilation cut a five-year
cumulative-volume discrepancy below 1\%
\parencite{Sedaghatdoost2019,Chelil2022b,Schoups2005}. The review literature
concludes no universal model exists and that automatic calibration is
preferable because it supports uncertainty estimation \parencite{Singh2019b}.
Crucially, evaluation choices are themselves substantive: the Kling--Gupta
efficiency consistently outperformed and could reverse rankings relative to
Nash--Sutcliffe, and regression-based tests exposed gross over-prediction
invisible to $R^2$
\parencite{Bruhn2025,STONE1995,SAVABI1995,Wessel2022}. Measurement protocols
matter equally: calibrated V-notch and compound weirs, flow meters, and
flow-proportional sampling minimise bias in discharge and load estimates that
otherwise propagates into nutrient budgets
\parencite{Christianson2019,BinShokrana2021,Napierala2023,Katuwal2024,Enamorado2007,Dialameh2022,Young2023}.

\subsection{Data-driven, remote-sensing, and geostatistical methods}

Data-driven methods frequently match or outperform process models for the
specific outputs on which they are trained while requiring less data and
computation. Neural networks predicted real-time water-table control and
saturated conductivity more accurately than pedotransfer functions
(the ANN prediction of saturated conductivity reaching $R^2$ up to 0.96),
random forests outperformed linear regression
for agronomic prediction, and an Attention~U-Net surrogate reached test $R^2$
of 0.997--1 for head, drainage volume, and salt discharge while drastically
reducing computation
\parencite{Yang2000,Ellafi2021,Djiemon2019,Lei2026,Sarangi2006,Bakhsh2005b,Bakhsh2007b}.
Parsimonious surrogates (Cubist, SVM, NIT-DRAIN) can match process models for
flow while training far faster, though process models remain preferable for
climate extrapolation
\parencite{Sharma2003,Negm2016,Frederiksen2023,Chelil2022,Li2025}. Remote
sensing combined with deep learning now maps drained extent at regional to
national scales with accuracies commonly exceeding 85--96\% and totals within a
few percent of census estimates: AgTile-US reported 22.34~Mha ($-0.35\%$ versus
USDA) and SEETileDrain achieved 0.96 testing accuracy
\parencite{Valayamkunnath2020,Wan2024,Moller2018}, while convolutional
segmentation networks reliably delineate visible drain lines and outperform
classical edge detectors \parencite{Redoloza2023,Song2021,Woo2023,Naz2009}.
Ancillary sensing infers drainage-relevant surface conditions through SAR,
Landsat SWIR, and vegetation indices
\parencite{Meunier2001,Rijal2013,Xu2018,Khanal2020}. Finally, kriging-family
methods are consistently the most reliable interpolators for the spatially
variable soil properties governing drain spacing: in the Nile Delta,
conductivity was isotropic with a 5.6~km range and kriged values were
recommended as design inputs, while accounting for variability widened spacing
about 5\% and saved substantial pipe, and neglecting $K_{sat}$ variability
could bias design spacing by $-27\%$ to $+3\%$
\parencite{GALLICHAND1991,GALLICHAND1992,GALLICHAND1993,Moustafa1998b,Moustafa2000,Moustafa1998,Akis2015,Sepaskhah2004}.

\subsection{Mapping drainage infrastructure and quantifying extent}

Because buried networks are frequently undocumented, a substantial body of work
establishes that non-invasive remote sensing and near-surface geophysics can
detect and map subsurface drainage, with a consistent conclusion that no single
method or spectral band is universally superior and that detection depends
strongly on soil moisture, soil composition, drain depth and diameter, and
survey timing relative to rainfall and field operations
\parencite{Carlsen2024,Wienken2024,Koganti2021}. In the largest multi-site
assessment (29 sites), UAV visible-color imagery detected drain lines at 48\%
of sites, multispectral at 59\%, and thermal infrared (TIR) at 69\%, so
acquiring all three bands is best practice \parencite{Allred2020}. TIR is the
most consistently effective single band, detecting roughly 60\% of
infrastructure even under dry conditions when visible imagery failed, though it
exploits apparent temperature differences below 1.0\,\textdegree C, requires
careful timing (optimally within 12~h of $\sim$2.5~mm rain and late morning to
early afternoon), and can fail under low contrast, where RGB imagery guides
interpretation
\parencite{Allred2018,Woo2019,Allred2021,Allred2022,Kratt2020}. Aerial visible
and multispectral photography likewise reveals functional drains through soil
and vegetation contrasts and has repeatedly exposed systems absent from
official registries \parencite{Tlapáková2017}. Ground-penetrating radar senses
pipes directly, with 250--1000~MHz antennas effective, antenna orientation
interacting strongly with soil moisture, and RTK/GNSS-positioned surveys
resolving system type in only four to six hours without dense grids; 3D-GPR
succeeds where soil electrical conductivity is low enough to limit attenuation
\parencite{Allred2005,Allred2013,Allred2018b,Koganti2020}. Geoelectrical
methods coupling time-lapse ERT with injected water or pressurised air locate
networks by induced resistivity change
\parencite{Jouen2018,Henine2021}. The best-supported strategy is combining
methods with complementary sensitivities---particularly UAV RGB/TIR imagery
with GPR, which adds pipe depth and orientation---since each has distinct
favourable conditions and linear farm features are otherwise easily confused
with drain lines \parencite{Koganti2021,Carlsen2024,Wienken2024,Allred2020}.

Beyond field scale, machine learning and hydrologic filtering extend mapping to
watersheds and regions: satellite big data with random forests map potential
drainage, and combining remote sensing, decision-tree classification, and
groundwater-table information improves watershed delineation
\parencite{Cho2019,Yimer2023}. Once layouts are known, GIS/GPS and
DEM-based tools support design and management; a suite of QGIS plugins tested in
Illinois streamlined identifying landscape features, producing topologically
sound flow networks, and optimising pipe depth and size in under ten minutes
versus over two hours for proprietary software
\parencite{Shedekar2018,Falasy2025,Falasy2024,Falasy2025b}. Quantifying extent
frames the significance of drainage effects: regional mapping estimated about
25\% of Germany drained (30--40\% on hydromorphic soils), U.S. datasets diverge
at watershed scale so best estimates remain locally conditioned, and this
extent rests on validated salt- and water-balance models, empirical
effluent-quality relationships, and geophysical (EMI/apparent-EC) monitoring
that together allow design principles to transfer with confidence to new sites
\parencite{Hirt2005,Hirt2011,Jame2022,Chen2025,Tian2023,Bahçeci2006,Skaggs2006,Narjary2024,Qin2025}.

% ----- knowledge/11_conclusions.tex -----
\section{Conclusions}

Across observational, tracer, geophysical, experimental, and modelling lines
of evidence, a coherent and strongly supported picture of agricultural
subsurface drainage now emerges. Its foundation is hydrological: tile networks
are the master control on the water balance of drained land. By intercepting
water that would otherwise recharge shallow groundwater or leave as surface
runoff and routing it rapidly to streams, drains commonly carry $40$--$65\%$ of
streamflow in humid settings and dominate field-scale water export
\parencite{Schilling2019}. This
repartitioning reshapes the flow regime in consistent ways---elevating low
flows and the baseflow fraction, exerting only modest and scale-dependent
effects on the largest peaks, shortening groundwater travel times, and
homogenising catchment response so that drained basins behave far more
uniformly and scale-invariantly than their undrained counterparts. Relative to
undrained or surface-drained land, improved subsurface drainage generally
reduces surface runoff, sediment, and phosphorus even as it increases nitrate
delivery. Because this hydrology governs everything downstream, the physical
behaviour of drainage is the necessary starting point for understanding its
water-quality and agronomic consequences.

That hydrological primacy is clearest for nitrate, the dominant nitrogen
species leaving drained cropland and the best-quantified consequence of
drainage. The subsurface pathway carries the great majority of nitrogen loss;
concentrations routinely approach or exceed the $10~\mathrm{mg~L^{-1}}$
drinking-water limit; and loads are governed first by drainage discharge, so
export is strongly seasonal, concentrating in the non-growing season and in a
few high-flow episodes \parencite{Kladivko2021,Speir2021}. Concentration--discharge behaviour is broadly
chemostatic, and the nitrate is largely organic-derived and delivered rapidly
through macropores, with denitrification providing real but only partial
attenuation. Phosphorus tells a parallel but distinct story: although it moves
in far smaller mass, the long-standing assumption of a filtered ``clean''
subsurface pathway has been overturned. A tiny, hydraulically active volume of
biopores and cracks delivers particulate phosphorus and sediment to drains
within hours, so tiles frequently dominate edge-of-field phosphorus loads
despite lower concentrations than surface runoff, with dissolved losses tightly
coupled to soil phosphorus status and legacy phosphorus sustaining loads at
watershed scale \parencite{Pease2018,Sharpley1995}.

Against these losses the evidence supports a coherent, tiered hierarchy of
mitigation with quantified effects. In-field, nitrogen-rate optimisation toward
the economic optimum, cover crops---especially cereal rye, clustering around
$40$--$60\%$ nitrate-load reduction \parencite{Malone2018,Malone2014b,Malone2023}---diversified rotations, and perennials act
in predictable directions, while tillage is a weak and inconsistent lever. At
the design scale, spacing and depth trade drainage volume against loss, and
controlled drainage and drainage water recycling reduce loads largely in
proportion to the flow they retain, delivering water-quality benefits together
with yield resilience. At the edge of field, constructed wetlands, denitrifying
woodchip bioreactors, reactive phosphorus media, and saturated buffers
intercept loads through well-characterised controls---residence time,
temperature, loading, and pH \parencite{Christianson2013,Jaynes2008}. Two cross-cutting lessons recur: these practices
redistribute nitrogen and can trade dissolved for particulate phosphorus or one
loss pathway for another, so single-pathway evaluation misleads; and their
benefits are scale-attenuated, diluted between the tile and the watershed
outlet. Meeting downstream targets therefore requires stacking in-field and
edge-of-field measures, and traditional surface-water practices that fail to
intercept subsurface flow are largely ineffective over drained land. The same
preferential-flow and hydrologic controls extend to pesticides, manure-borne
microbes and resistance genes, trace elements, and dissolved greenhouse gases,
where transient concentrations can be high even as mass losses stay small and
management often matters more than drainage itself.

Alongside these environmental costs, drainage is an established and valuable
production technology. It reliably raises crop yields---most strongly on poorly
drained soils and in wet years---leaches salts to reclaim saline and
waterlogged land, and enables safe reuse of saline and brackish water where salt
balance is maintained. Controlled variants add further yield and water-use
benefits, and formal appraisals across many regions confirm favourable
financial and social returns, with careful design and mechanisation
substantially reducing cost. Underpinning the whole synthesis is a mature,
validated toolkit: process models such as \textsc{DRAINMOD} reproduce drainage
flow and nitrate reliably after calibration, with skill improving from daily to
annual scale ; the physics of saturated--unsaturated and preferential flow is
well understood; analytical spacing solutions, data-driven surrogates, and
remote sensing complement one another; and geophysical and machine-learning
methods now map buried networks and drained extent at field to national scales.

The strength of these conclusions rests on their convergence. Independent
methods, disparate climates, and separate scales point to the same mechanisms
and comparable magnitudes, so what is known about tile drainage is known with
some confidence. The unifying message is that drainage is fundamentally a
hydrological intervention: it sets the pathways, timing, and volume of water
movement, and water-quality outcomes, agronomic gains, and the effectiveness of
every management practice follow from that control. Recognising drainage water
quality as a design criterion alongside production, and combining complementary
practices across the field-to-watershed continuum, is the consolidated
implication of the evidence assembled here.


\part{Part II --- Research Frontier}
% ===== track: frontier -- Tile Drainage: Research Frontier =====

% ----- frontier/00_abstract.tex -----
\section{Abstract}

Agricultural subsurface (tile) drainage underpins crop production on poorly drained and irrigated land, yet the same conduits that make these soils farmable also short-circuit the landscape's natural buffering, delivering nitrate, phosphorus, salts, and a widening list of constituents to ditches, streams, and coastal waters. Decades of field experiments, monitoring, and modelling have produced a detailed literature on drained-system behaviour and on the practices---nutrient management, controlled drainage, cover crops, edge-of-field treatment, salinity reclamation---that might reconcile productivity with water quality \parencite{Randall2001}. This review is the research-frontier component of a larger synthesis, framed deliberately around what remains uncertain, untested, or unstudied rather than around settled results. Its recurring meta-finding: ideas have outrun the data able to test them--- headline results rest disproportionately on short records, single sites, and narrow agro-ecological windows, and progress in every domain is bounded by the length, density, replication, and consistency of the underlying observations. We proceed from enabling to domain frontiers, then to consequences and integration. The enabling foundations are an observational base of longer, denser, and harmonised monitoring; process-based drainage and water-quality models whose credibility is limited less by their formulation than by the narrow envelope in which they have been validated; and the coupled problems of mapping the largely unmapped drainage network and scaling instrumented-field knowledge to the watersheds where water-quality targets are set. On this footing we treat the domain frontiers of nitrogen, phosphorus, and cropping-system losses; active water-table management and drainage-system redesign; edge-of-field treatment such as bioreactors, wetlands, and phosphorus-removal structures; and salinity reclamation, drainage-water reuse, and disposal. Because removing a constituent from one pathway can shift it to another and altering the field water regime can trigger slow biogeochemical and pedological change, we then consolidate the unintended consequences and multi-objective tradeoffs---greenhouse-gas coupling, pollution swapping, emerging contaminants, trace-element legacies, and drainage-induced soil evolution---that nitrate- and yield-centred accounting has overlooked. Finally, we draw the integrative strands together: no single practice or model is a silver bullet, and decisive advances will come from coupling mechanistic understanding with spatial targeting, economics, policy, and shared decision-support infrastructure. The unifying message is that the next decade of progress will be secured less by new equations or novel practices than by sustained, well-instrumented, openly reported, and geographically diverse field campaigns, coordinated data standards, and interdisciplinary integration tying mechanism to tools and tools to adoption.

% ----- frontier/01_introduction.tex -----
\section{Introduction}

Agricultural subsurface (tile) drainage is one of the foundational
technologies of modern crop production: by lowering the water table on
poorly drained land it secures trafficability, root aeration, and yield
across vast areas of the temperate and irrigated world. Yet the same
engineered conduits that make these soils farmable also short-circuit the
landscape's natural buffering, delivering nitrate, phosphorus, salts, and
an expanding list of other constituents rapidly to ditches, streams, and
ultimately to coastal waters. Decades of field experiments, monitoring
campaigns, and modelling have produced a large and detailed literature on
how drained systems behave and on the practices---nutrient management,
controlled drainage, cover crops, edge-of-field treatment, salinity
reclamation---that might reconcile productivity with water quality. The
distinctive feature of that literature, however, is that its most secure
and most frequently repeated statements are not about what is known but
about what remains unknown. This review is organised around that
observation.

Read across the field, the recurring signal is a precondition rather than
a conclusion. Headline results are drawn disproportionately from short
records, single sites, and narrow agro-ecological windows, and their
authors caution explicitly against premature generalisation. The blunt
appeals are pervasive---that ``more research is needed''
\parencite{SHALIT1995}, that ``the complexity of this subject requires
more research'' \parencite{Malone2014}, that ``further research in this
area is needed'' \parencite{Woo2023}, and that ``further evaluations are
warranted before any concrete conclusions could be drawn''
\parencite{GUPTA1993}---and they mark genuine gaps in evidence rather than
rhetorical closings. \textcite{Rogger2017} pitch the same concern at the
level of the whole discipline, calling for complex-systems thinking that
links processes across time scales, long-term experiments on
physical--chemical--biological interactions, attention to connectivity and
pattern across spatial scales, and multidisciplinary consortia that share
data and terminology under known boundary conditions. The binding
constraint that emerges from these appeals is not a single missing
experiment but a coupled deficit: coordinated, long-term, spatially
replicated, and harmonised data on the one hand, and better process
representation in models on the other, must advance together, because
neither longer records nor more elaborate models resolve the field's open
questions alone.

This paper is the research-frontier component of a larger review of tile
drainage, and it is deliberately framed around uncertainty. Where the
companion component consolidates settled knowledge, our purpose here is to
gather, structure, and sharpen the open questions---the priorities, the
hypotheses that remain to be tested, and the methods, instrumentation, and
datasets still needed to test them. We therefore do not attempt to resolve
any of the gaps we identify; we set them up, tracing how each is posed in
the primary literature and how the strands interlock. The aim is a map of
where the field should go next, in which the confidence of a claim about
what is \emph{not} yet understood is treated as a finding in its own right.

The synthesis proceeds from enabling frontiers to domain frontiers and
then to their consequences and integration. We begin with the cross-cutting
foundations on which every later domain depends: the observational
foundation of longer, denser, and harmonised monitoring; the process-based
drainage and water-quality models whose credibility is limited less by
their formulation than by the narrowness of the conditions under which they
have been tested; and the linked problems of mapping the largely unmapped
physical drainage network and scaling instrumented-field knowledge to the
watersheds where water-quality targets are actually set. On this footing we
turn to the domain frontiers---managing nitrogen, phosphorus, and
cropping-system losses; active water-table management and the re-evaluation
of drainage-system design; edge-of-field treatment technologies such as
denitrifying bioreactors, wetlands, and phosphorus-removal structures; and
salinity reclamation, drainage-water reuse, and disposal in arid and
irrigated systems. Because removing a constituent from one pathway can
shift it to another and altering the field water regime can set slow
biogeochemical and pedological change in motion, we then consolidate the
unintended consequences and multi-objective tradeoffs---greenhouse-gas
coupling, pollution swapping, emerging contaminants, trace-element
legacies, and drainage-induced soil evolution---that a nitrate- and
yield-centred accounting has historically overlooked. Finally, we draw the
integrative strands together, because translating frontier science into
outcomes depends on how practices are bundled and spatially targeted, how
they perform under a changing climate, and whether they are economically
and institutionally adoptable.

A single thread runs through all of these: no individual practice or model
is a silver bullet, and the decisive advances will come from coupling
mechanistic understanding with spatial targeting, economics, policy, and
shared decision-support infrastructure. Underlying that ambition is a data
imperative---better, longer, and more spatially representative
datasets---that recurs as the foundational requirement for validating
models, resolving tradeoffs, and delivering site-specific, climate-resilient
guidance \parencite{Ma2025,SKAGGS1994,Deng2011}. The sections that follow
elaborate that agenda frontier by frontier.

% ----- frontier/02_the-observational-foundation-longer-denser-and-h.tex -----
\section{The Observational Foundation: Longer, Denser, and Harmonised Monitoring}

Across the drainage literature, the field's most secure statements are
about the limits of its own evidence. Headline results are drawn
disproportionately from short records, single sites, and narrow
agro-ecological windows, and their authors caution explicitly against
premature generalisation. The recurring signal is therefore not a settled
conclusion but a precondition: records long enough, dense enough,
spatially replicated enough, and consistent enough to separate a
management signal from the natural variability of weather, hydrology, and
slow soil change. \textcite{Rogger2017} pitch this at the level of the
whole discipline, calling for complex-systems thinking that links
processes across time scales, long-term experiments on
physical--chemical--biological process interactions, attention to
connectivity and patterns across spatial scales, typologies and
meta-analyses that generalise findings, and multidisciplinary consortia
sharing data and terminology under known boundary conditions. The sheer
frequency of blunt appeals---that ``more research is needed''
\parencite{SHALIT1995}, that ``the complexity of this subject requires more
research'' \parencite{Malone2014}, that ``further research in this area is
needed'' \parencite{Woo2023}, and that ``further evaluations are warranted
before any concrete conclusions could be drawn''
\parencite{GUPTA1993}---marks genuine gaps in evidence rather than
rhetorical closings. Because every later domain in this review inherits
these limits, the observational foundation is treated here as a frontier
in its own right.

\subsection{Extending the temporal record}

The single most widely voiced need is for records long enough to
disentangle treatment effects from climatic noise and from slow,
directional change. \textcite{Clifford2023} emphasise truly long-term
experiments, and \textcite{BUHLER1993} caution that a lag of several years
may precede any objective evaluation of an altered practice. Minimum
horizons are proposed---at least five years, with continued observation as
systems are extended in length or depth over time \parencite{Ali2013}, and
lengthening national time series expected to sharpen analysis as records
grow \parencite{Ali2001}. Many studies attach conclusions to explicit
requirements for further years: to quantify a denitrification wall's
effectiveness and its long-term cost per unit nitrate removed against
natural variability \parencite{Jaynes2008}, to establish how long treatment
benefits persist \parencite{Dalhem2019,Osterholm2005}, and to test the
replicability of results in time and confirm that observed differences
endure under comparable conditions
\parencite{ROLLIN1987,Reinhart2020,Rijal2012,GiménezValero2025}, all
echoing the disciplinary case for long-term monitoring at multiple scales
\parencite{SPALING1995}.

Nutrient and hydrological studies converge on the same horizon.
Year-round measurement of flow and concentration for at least several years
on new sites is deemed necessary to set management priorities
\parencite{Kladivko2001}, and long-term records are singled out to close
water and nutrient balances and study legacy pools \parencite{Sharpley2019}
and to disentangle land-management from climate-change effects, with the
corollary that policy evaluation demands patience
\parencite{Bauwe2020,Kelley2017,Gillingham2000,Lemke2011}. Smoothing
year-to-year variability requires evaluating both phases of a rotation
\parencite{Randall2003,Randall2005}, assessing responses on an annualised,
multi-year basis \parencite{Waring2024,Pittelkow2017}, and freeing
drainage-water-management yield effects from confounding
\parencite{Cooke2012}; only long-term records across a range of climatic
conditions, gathering plant, soil, and water measurements together, yield
appropriate guidelines \parencite{RANDALL1995}. Data spanning at least a
normal, a wet, and a dry year are needed to resolve losses and transport
routes \parencite{vanderSalm2012}, and as long a record as possible is
required for trend estimation because water quality may lag the
management-driven change \parencite{Ulén2007}. Long records are equally the
only way to detect slow processes---soil organic-carbon recovery
\parencite{Kumar2014}, discharge-driven organic-matter export
\parencite{Manninen2023}, and whether drainage-induced soil evolution
continues or stabilises over decades \parencite{Frison2009,Frankl2023}.
Record length alone is insufficient: robust seasonal conclusions require
year-round coverage, especially in the non-growing season
\parencite{Klaiber2020,VanEsbroeck2016}, high-frequency records without
gaps and instruments checked as often as possible \parencite{Hodaj2017},
and monitoring timed to infrastructure interventions such as drain repair
\parencite{Wesström2015}.

\subsection{Broadening spatial and geographic scope}

A second frontier is testing whether site-specific findings hold across
other soils, climates, crops, and management regimes. Authors repeatedly
defer broad adoption until similar data are collected from further sites and
years \parencite{Menegaz2025,Zimmerman2017b,Terrado2015}, from additional
agro-ecological zones \parencite{Ali2001}, from other settings for
practices such as the two-stage ditch \parencite{Hodaj2017}, and from
adverse or compacted conditions where a technology's merit is untested
\parencite{Bigelow2001}. The most explicit prescription is that
controlled-drainage and subirrigation studies be conducted outside the U.S.
Midwest and Canada across diverse climates, crops, and soils, and---
crucially for generalisation---that site characterisation (slope,
particle-size distribution by horizon, hydraulic conductivity, drainage
class, and moisture) be reported with measures of variability
\parencite{Singh2022b}. Repetition at larger scales and on contrasting soils
is stressed as a key modulating factor
\parencite{Bilotta2008,Goeken2015}, including for process studies such as
earthworm activity on soils of differing texture and rainfall
\parencite{Nuutinen2001,Nuutinen2011}. Climate-focused transferability
demands high-resolution scenarios and multi-model uncertainty analysis at
additional sites \parencite{Rasmussen2018,Jeantet2022,Jin2008}, and testing
across regions, soils, and climates recurs throughout as the precondition
for localised design guidance
\parencite{Singh2020b,Malota2024,Salazar2008,Holbak2021}. At the level of
synthesis, robust conclusions will require geographically diverse and
methodologically harmonised studies so that meta-analyses can disentangle
management effects \parencite{Wu2025}.

\subsection{Monitoring infrastructure, instrumentation, and data tools}

A distinct set of needs concerns the apparatus that makes long-term data
collection feasible and reliable. \textcite{Kanwar1999} describe how
automatic samplers collect discrete alongside continuous samples, how data
loggers can be extended to soil temperature and moisture, and how relays or
pressure switches could detect sump-pump operation, advising that
equipment-to-logger distances be minimised, data downloaded automatically to
reduce loss, and annual maintenance---cleaning check valves and calibrating
flowmeters---be performed. Diagnostic methods are themselves research
targets: reducing pressure variations by doubling equipment and injecting
smoke from both ends of a drain section could enable work with longer drain
lines, though transport pathways and observed clustering remain unresolved
\parencite{Petersen2012}, and selected pipes, envelopes, perforations, and
surrounding soil should be observed with permanent monitoring of the
sampling process to confirm favourable results are maintained
\parencite{MartínezBeltrán2007}. Closing the water budget requires
high-frequency discharge, runoff, and seepage measurement, process
measurements of microbial activity, denitrification, and redox, and a
publicly available dataset \parencite{Frankenberger2024}. Software and data
infrastructure are an increasingly prominent part of this frontier:
readers are directed to evolving data products such as the National
Elevation Dataset and USGS resources \parencite{Shedekar2018}, and a series
of drainage-network configuration and sizing tools describe themselves as
under continual development and openly invite testing with additional
datasets \parencite{Falasy2025,Falasy2025b,Falasy2024}.

\subsection{Tracers, isotopes, and source apportionment}

Measurement and tracing techniques are repeatedly identified as the means
to resolve pathways that bulk flow-and-concentration records leave
ambiguous. \textcite{Williams2022} call for higher-frequency (sub-daily to
daily) measurements, further work on groundwater dynamics in dissolved-P
transport, and more comprehensive source-apportionment methods.
Water- and nitrate-isotope approaches are advanced to quantify processes
directly: coupling water and nitrate oxygen isotopes and quantifying the
$\delta^{18}$O of nitrification-derived nitrate to constrain evaporative
enrichment \parencite{Hu2024}, monitoring nitrate concentrations and
isotopic compositions with multiple tracers across topsoil, drains, and
ditches to characterise the soil contribution to river nitrate
\parencite{Briand2017}, and boron isotopes to strengthen contaminant
tracing around feedlots and fields \parencite{Komor1997}. High-frequency
tracer measurements are urged to quantify water ages and represent both
celerities and velocities \parencite{Williams2021}, alongside event
sampling and bulk-soil-water isotopes to constrain young-water fractions
\parencite{DimitrovaPetrova2020}, high-resolution specific conductance for
end-member mixing \parencite{Kronholm2015,Schilling2019}, and
conservative-tracer or water-isotope studies to partition quick flow and
baseflow \parencite{Frederiksen2023}. Explicit pathway partitioning---
reporting surface-inlet presence and splitting water and nutrient loads by
pathway via before-after control-impact designs, hydrograph separation, or
tracers \parencite{Flores2021}, evaluating overland-flow versus tile-flow
contributions \parencite{Kokulan2019}, and verifying groundwater flow paths
and plant water sources with stable isotopes
\parencite{Islam2022,YliHalla2022}---recurs throughout. The role of
groundwater outflow in nutrient loading is flagged as under-quantified and
in need of dedicated measurement and modelling
\parencite{Turunen2013,Jokinen2024,Barkle2021}, and finer temporal analyses
of the flow--concentration relationship, including detailed storm-hydrograph
separation, are needed to test event-scale hypotheses
\parencite{Lagzdins2015,Danalatos2022,Kladivko2004}.

\subsection{Standardisation, harmonisation, and comparability}

For accumulating data to support synthesis, methods and reporting must be
harmonised. The value of future meta-analyses is tied directly to
methodologically harmonised, comprehensive, long-term monitoring of the full
suite of environmental impacts \parencite{Wu2025}. \textcite{Zhang2018} set
out the most detailed programme: improving the quality and comparability of
existing indicators through harmonised data collection, monitoring,
interpretation, and reporting; developing new indicators for emerging
concerns while ensuring scientific credibility, sensitivity to change, and
ease of interpretation; establishing standardised evaluation procedures
including thresholds and reference points; recognising data gaps and
improving upscaling and uncertainty assessment; and incorporating
earth-observation data. These priorities complement the call for typologies,
meta-analyses, and shared data and terminology across research consortia
\parencite{Rogger2017}. At the study level, meta-analytic comparability
requires complete hydrologic description of each plot, standardised
analytical methods, adequate monitoring periods, and full reporting of
concentrations with descriptive statistics \parencite{Ghirardini2019}, and
soil biological and soil-health indicators must be sampled at consistent
times under similar moisture and temperature conditions to be comparable
across sites \parencite{Sherbine2023,Li2023b}.

\subsection{Experimental design for detectable effects}

Longer, standardised records are useful only if the design can resolve the
effect sought. \textcite{Schilling2013} specify the ingredients for
paired-basin detection: establishing a suitable correlation among paired
basins across climates, restricting paired sampling to the March--July
high-flow period, delaying post-treatment monitoring to enhance the chance
of detecting change, and adopting practices capable of achieving at least
the minimum detectable change. Effects should be assessed simultaneously at
multiple scales, from field to watershed \parencite{Sunohara2016}, paired
with within-field and within-plot studies to elucidate the mechanistic
processes of export and additional field-scale studies within high-export
sub-watersheds \parencite{Irvine2019}. Larger, more representative trial
areas are sought to test effective-yield formulas \parencite{Xin2016c} and
to refine tracer-based ($^{15}$N) loss observation \parencite{Xiong2025},
and mesocosms are proposed to quantify nutrient-loss responses to
management change under controlled conditions \parencite{Loper2024}.
Because field-scale drainage trials are typically unreplicated, several
authors treat the analysis itself as an open question: whether beta
distributions, time series, or new methods best analyse unreplicated paired
experiments \parencite{Pagliari2021}, and whether directed hypothesis tests
should be favoured over replicate spatial analyses \parencite{Andrus2014}.
Where replication is feasible, replicated field experiments, matched
instrumentation, and improved sampling designs are recommended to support
statistical inference
\parencite{Okuda2020,Wamono2016b,Davis2019,Jacquemin2020}. Together these
design refinements define the bridge between a lengthening, harmonised
observational record and the detectable, transferable conclusions the rest
of the field depends upon.

% ----- frontier/03_advancing-process-based-drainage-and-water-quali.tex -----
\section{Advancing Process-Based Drainage and Water-Quality Models}

Process-based drainage and water-quality models---DRAINMOD and its
derivatives, SWAT, APEX, RZWQM2, HYDRUS, MIKE~SHE, and numerous salinity
codes---are expected to carry much of the predictive burden as field trials
cannot span every combination of practice, soil, and climate. Yet across the
literature their credibility is constrained less by their formulation than by
the narrowness of the conditions under which they have been tested: models are
typically calibrated at a single site, over a handful of seasons, for one
soil--climate--management combination, and their authors are nearly unanimous
in flagging broader empirical validation, missing process data, and weak
algorithmic components as the outstanding priorities
\parencite{Li2006,Qi2015,Youssef2006,Radcliffe2015,SKAGGS1994}. This section
gathers the model-algorithm and model-structure gaps that recur throughout the
field.

\subsection{Broadening the empirical basis for validation}

The most frequently articulated need is validation data spanning a wider range
of conditions than any single study can supply. \textcite{Youssef2006} describe
rigorous field testing across differing management, soils, and climates as a
need common to nearly all simulation models, and warn that assumptions must be
re-verified whenever a code is transferred to a new area
\parencite{Kelleners2000}; similar calls for broader-condition testing,
sensitivity analysis, and parameter refinement recur for pesticide, salt, and
flow models \parencite{Branger2009,Hirekhan2007,Yang2008b,Xian2017}. Models are
repeatedly reported as requiring testing over longer periods and additional soil
types \parencite{Salazar2009,Li2026}, and their effectiveness under new
controlled-drainage or erosion conditions is judged critically needed before
higher credibility can be claimed \parencite{Negm2017,Turunen2020}.

A recurring, more ambitious hypothesis is that empirical coefficients could be
replaced with functions of measurable site attributes. \textcite{Li2006} commit
to testing across climate zones, soils, and management while anticipating that
coefficient constants might become functions of soil texture, converting
site-specific tuning into transferable structure; \textcite{Chelil2022} cast the
same concern as spatial robustness, arguing that additional fields are needed to
confirm that parameters such as nitrate-transfer velocities are transferable
rather than calibration artefacts. Verification at independent sites is treated
as a precondition, not optional confirmation
\parencite{Azevedo1997,Qi2015,Kalita1998,Sadhukhan2019b}, with inconsistencies
producing conflicting parameter values resolved first \parencite{Dar2022} and
empirical codes transferred only through re-parameterisation with local soil,
geology, crop, and climate data, or region-specific coefficients where reference
values exist \parencite{Azevedo1997,FathAllah2025,Rijal2013}. The open question
is how much local data such transfer requires. Extending validation across the
climatic and design envelope is closely related: humid-region models need
testing under semi-arid regimes and after prolonged dry spells
\parencite{Prasher1996}, longer monitoring records give firmer conclusions
\parencite{Ebrahimian2011,Grazhdani1996}, and design geometry demands
experiments at spacings wider than 30~m and regional answers on
drainage-efficiency loss and design storms
\parencite{Mathew2001,Wiskow2003,MULLER1988}. Long simulations of 50--100 years
are argued necessary to judge whether a design meets its criteria
\parencite{Mendez2004,Wang2006b,Ale2009,Häggblom2019}. Portability, in short, is
a hypothesis to be tested, not a property to be presumed
\parencite{Hornbuckle2012,Shipitalo2016,Yousfi2014}.

\subsection{Field and process data still needed to constrain models}

Beneath the demand for broader validation lies a scarcity of process-level
measurements. Nitrogen dynamics emerge as the least well-constrained component:
\textcite{David2009} argue that systematic denitrification measurements are
required to test and refine models, their ideal being a multi-year,
tile-drained corn--soybean study measuring every major N flux, with a pragmatic
fallback of a few key fluxes plus denitrification. Field- and
laboratory-grounded parameterisation, denitrification foremost, is likewise
judged necessary \parencite{Salazar2009}, together with routine evaluation and
reporting of soil-N fluxes in SWAT \parencite{Ikenberry2017}, long-term
carbon--nitrogen estimates to reconcile plot- and subbasin-scale discrepancies
\parencite{Singer2011}, and resolution of the lag between practice
implementation and measurable reduction \parencite{Law2022,Tao2021}.
High-frequency process measurements of microbial activity, denitrification,
soil-P pools, and redox are urged for drainage water management, alongside
closure of water and nutrient balances \parencite{Frankenberger2024,Sharpley2019}.

Phosphorus is similarly data-limited, and its dominant pathway---preferential
flow---is the hardest to observe. Many phosphorus-index factors need
investigation of their effect on field-to-stream transfer, and index rankings
must be shown to reflect actual transfer risk
\parencite{Bechmann2007,Beauchemin2003}. Confirming P-movement trends requires
quantitative tracer measurement of preferential flow
\parencite{Salazar2011,Tiemeyer2009}, structural deficiencies around surface
runoff and tile P transport need repair \parencite{Apostel2021}, and accurate
capture of flashy export demands automated samplers, real-time sensors, and
high-resolution topographic survey of ditch sediment before and after storms
\parencite{Dialameh2022,Warsta2014}. Supporting empirical work includes further
testing of CaCl$_2$-P as a dissolved-reactive-P estimator \parencite{McDowell2001b},
regression models linking drainage parameters to P losses \parencite{Morrison2013},
tests using actual effluent under fluctuating flow \parencite{Sanford2016}, and
long-term monitoring of P sources \parencite{Jia2012}. Sampling design is itself
a constraint: weekly sampling may be needed to capture nitrate spikes before
they pass below suction-cup samplers \parencite{Zotarelli2007}, more frequent
early-season and summer-event sampling is advised
\parencite{Kalita1998,Tiemeyer2009}, and denser sampling geometry sharpens
semi-variograms and optimal sampling density \parencite{GALLICHAND1992,Upadhyay2019}.

\subsection{Preferential and macropore flow representation}

The representation of preferential flow is perhaps the most frequently cited
frontier, the recurring theme being that its contribution to drains must be
characterised as a function of soil type, water content, and practice before
models extrapolate. \textcite{Gaillot2023} and \textcite{Fox2004} call for
mechanistic characterisation of macropore--drain connectivity and of the
``express fraction'' under varying saturation above the drain, and
\textcite{Askar2020} stress that rigorous testing of DRAINMOD's empirical
macropore approach requires water-table-depth measurements alongside drainage
and runoff, direct macropore-flow inputs, and crack geometry and dynamics.
Process improvements span better matrix--macropore partitioning calibrated
against sub-hourly drainage in RZWQM2 \parencite{Li2023d}, dynamic connectivity
or effective macroporosity at watershed scale \parencite{Poon2021}, micropore
diffusion and site-specific conductivities in dual-permeability descriptions
\parencite{Villholth1998,Villholth1998b}, hysteresis and biopore-flow initiation
\parencite{Holbak2021}, manure-driven macropore development \parencite{Zhao2001},
crack re-establishment in EPIC \parencite{Wang2022}, and calibrated macropore
conductivity in three-dimensional dual-permeability models \parencite{Varvaris2018}.
Direct pathway identification is repeatedly urged---ground-penetrating radar and
tracers with explicit flowpath representation \parencite{Tähtikarhu2025,Salo2017},
dye-staining of routes reaching tile lines \parencite{Nielsen2015}, and
near-surface pathway characterisation \parencite{Frey2012,Allaire2011}. Central
quantitative questions remain the fraction of phosphorus and solutes delivered
via preferential versus matrix flow \parencite{King2015,Williams2016}, the range
of macropore flows across soils and management \parencite{Reid2012}, and by-pass
flow through cracks for full-year application
\parencite{LESAFFRE1988,BACKLUND1995,Ruark2009}. Runoff simulation under
crusting, cracking, and preferential flow \parencite{Smith2020}, explicit
representation of actual drain configuration \parencite{Stamm2002}, temporal
variability in hydrological-pathway activation for cracking clays
\parencite{Kokulan2019}, and seasonal, time-dependent shrink--swell and
hydraulic-conductivity measurement \parencite{Boivin2006,Chelil2022b,Jeantet2021}
complete this agenda.

\subsection{Coupled hydraulics, seepage, and soil heterogeneity}

A recurrent structural gap is the absence of models coupling soil pore-water
flow with surface flow inside drains and ditches. \textcite{Xin2016b} call for
coupling drain surface water with pore-water flow to represent the time-varying
head around the drain, and \textcite{Nousiainen2015} propose simulating in-pipe
flow to capture pressurisation and clogging, adding open-ditch networks.
Pressurised-flow dynamics are a specific frontier: coupling pipes to the water
table could introduce a validated temporary-retention function beyond design
capacity \parencite{Henine2010,Henine2014}, with pressurised periods analysed
separately from pressure-transducer data \parencite{Saadat2018}. Entrance
resistance, near-drain pressure heads, and rapid trench-backfill versus slower
inter-trench transport remain unresolved for heterogeneous systems
\parencite{Messing2006,Rimidis2003}, and model dimensionality is itself a
question---two- or three-dimensional treatment of sloping fields, collectors,
and regional management \parencite{Alakukku2010,FIPPS1989,Carlier2007,Mohanty1998}---within
broader integrating frameworks \parencite{Manguerra1998,SAVABI1993,RevueltaAcosta2021}.

Nearly every modelling study attributes residual error to unresolved
heterogeneity in hydraulic properties, especially conductivity. Authors seek
detailed representation of macropores, fractures, and sand lenses with specific
storage and residual saturation \parencite{DeSchepper2017}, depth-varying and
dual-porosity parameterisations in layered and paddy soils
\parencite{Wesström2003,Karandish2018}, and denser, depth-resolved
sampling---at least one conductivity measurement per hectare with
post-construction verification \parencite{PalaciosVélez2004}, more transect
points \parencite{Messing2006}, and horizontal as well as vertical
conductivities with hydrophobicity and hysteresis \parencite{Tähtikarhu2025}.
Anisotropy and long-term property change form a related sub-frontier
\parencite{Akis2015,Moustafa1998,Koivusalo1999}, and accurate saturated
conductivity and drainage-volume/upward-flux relationships, supported by
SSURGO-linked subsoil porosity, are emphasised as model inputs
\parencite{DeBoer2006,Schilling2015,DeSchepper2015}. Closing the water
balance---particularly seepage and groundwater exchange---remains persistent:
lateral seepage should vary during simulation \parencite{Feng2018}, groundwater
contributions could enter through double- or residual-storage models
\parencite{DaneshYazdi2016}, and groundwater-table data in objective functions
would sharpen inverse recovery of soil properties
\parencite{Sedaghatdoost2018,Sedaghatdoost2019}. Improved calibration strategy is
widely urged---parameter-estimation software and multiple conceptual models
\parencite{Boico2022}, pilot-point conductivity fields with proxy
machine-learning predictions \parencite{Schneider2022}, Kling--Gupta objective
functions \parencite{Brighenti2023}, and transient tests of steady-state
assumptions \parencite{Schilling2022,Gates2002,Burkhalter2006}. Contemporary
needs also span unsaturated modules for disconnected groundwater
\parencite{Kompanizare2024}, independent layer-to-layer drainage methods under
preferential flow \parencite{Ojeda2018}, and dynamic frozen-soil models with
measured structure \parencite{Ulén2016}.

\subsection{Nutrient routines, cold-season, and groundwater coupling}

Phosphorus and its transport through macropores emerge as the least mature
component across water-quality models. \textcite{Radcliffe2015} provide the most
comprehensive agenda: nested paired field experiments and structural additions
across codes---SWAT and ADAPT routing macropore flow directly to drains, a
macropore component for APEX, DRAINMOD simulating P losses and sediment-attached
P, and HYDRUS parameterising simpler models---with every study quantifying
predicted-P accuracy and uncertainty. Consistent with this,
\textcite{Guo2020b} recommend physically based tile subroutines, macropore
nutrient pathways, and particulate-P losses, \textcite{Ford2017} seek
understanding of how macropore flow and P respond to tillage, cover crops, and
drainage water management, and \textcite{BinShokrana2023b} identify improved
fertiliser-P partitioning and P-only performance statistics, with
\textcite{Gowda2012} anticipating a DRAINMOD phosphorus module. Within SWAT,
APEX, and RZWQM2, \textcite{Williams2016} argue that incorporating identified
flow pathways is critical, \textcite{Qi2018} recommend a soil-matrix-potential
Richards equation and revised particulate-P algorithm, and
\textcite{Moriasi2012} call for comparing dynamic and constant surface-depression
storage. Nitrogen-cycle needs include surface-residue and mobile/immobile-water
phases for nitrate transport \parencite{Malone2017}, complete N budgets in
calibration \parencite{Gillette2018}, and temporal decay coefficients
\parencite{Rudra2005}.

Soil N-process representation is repeatedly the weakest link. Priorities include
improving nitrification, denitrification, and mineralisation within a better
hydrology framework \parencite{Smith2020,Fang2012,Hu2007}, and recurring RZWQM
agendas---subsurface lateral flow below tiles, simulated soybean N fixation,
improved plant N uptake, and reduced winter nitrification
\parencite{Ma2007,Ma2008,Ma2001,Bakhsh2004}---together with revised winter
processes in the PNM model \parencite{Marjerison2016} and a more physically based
tile-nitrate model using ``soft data'' and a passive humus pool
\parencite{Qi2025}. Cold-season representation is singled out as inadequate:
better freeze/thaw and spring-thaw algorithms \parencite{Feyereisen2006b},
DRAINMOD-N improvement for freezing seasons \parencite{Yang2007,Saadat2020,Jiang2020},
a frozen-soil impedance factor and calibratable snow--rain threshold for RZ-SHAW
\parencite{Li2024b}, and continued RZWQM winter development
\parencite{Ahmed2007b,Motarjemi2021}. Coupling drainage models to groundwater is
a distinct need: SWAT-MODFLOW-RT3D for reactive transport once salinity is
incorporated \parencite{Addab2022}, water-table-fluctuation and
groundwater--surface-water exchange in SWAT to represent riparian nitrate
retention \parencite{Hoang2017}, joint spatial variation in leaching and
reduction across systems \parencite{Hojberg2017}, and defensible water balances
splitting artificial drainage from groundwater flow
\parencite{Barkle2021,Elsayed2023,Yang2008}. In arid settings, coupling of water
and salt dynamics with unsaturated models, root-growth modules, and
large-dataset validation is pressed
\parencite{Ebrahimian2015,Sarangi2006,Han2025,Liu2022,Guo2025}.

\subsection{Calibration, sensitivity analysis, and uncertainty}

How models are calibrated and their parameters interrogated is inseparable from
validity. Multiobjective frameworks integrating statistical and graphical
diagnostics tailored to the hydrological system are advocated
\parencite{Bruhn2025}, reliable soil-hydraulic data for the drain layer is
emphasised with interacting parameters calibrated and negligible ones fixed
\parencite{Wang2006}, and independent calibration against observed discharge is
recommended so interception and storage are not conflated with conductivity and
drainable porosity \parencite{Henine2022}. Pedotransfer estimates such as
Rosetta conductivity should be used cautiously, with further calibration so both
water-table depth and discharge are simulated \parencite{Malota2022,Timalsina2023}.
A common thread is identifiability---separating parameters that genuinely act
from those merely correlated. Automated and global approaches are widely urged:
automatic calibration with ensemble Kalman filtering for parameter identification
\parencite{Warsta2014,Sadhukhan2019b,Singh2019b}, global sensitivity analysis to
capture interactions that local one-at-a-time methods miss
\parencite{Timalsina2023,Chelil2022,Kacimov2000}, established bounds for the
SWAT sensitivity adjustment factor \parencite{Moriasi2013b}, and spatially
distributed streambed conductance with exchange fluxes as calibration targets and
drainage proxies \parencite{Karan2021}. Explicit uncertainty treatment is a
distinct, underdeveloped frontier: propagation of parameter and measurement error
\parencite{Alzraiee2013b}, the role of input-data and weather-data quality
\parencite{Upadhyay2019,Tao2021}, explicit estimation of P-model prediction
accuracy \parencite{Radcliffe2015,Beauchemin2003}, structural uncertainty via
multiple water-quality models \parencite{Kujawa2020}, probabilistic
drainage-water-management effectiveness \parencite{Shedekar2021}, independent
evaluation against withheld data \parencite{Schoups2005}, and better statistical
treatment of unreplicated paired experiments \parencite{Pagliari2021,Andrus2014}.
The recurring message is that models are too often reported as point predictions
when their inputs warrant propagated bounds.

\subsection{Model structure, coupling, and machine learning}

Beyond recalibration, structural refinements and integrative strategies are
judged necessary. Coupling complementary models is a common, largely untested
proposal: linking biogeochemical to spatial hydrological models such as SWAT or
MIKE~SHE for regional prediction \parencite{Li2006}, coupling nitrate-transport
and drainage models with variational assimilation of spring-fertilisation
information \parencite{Chelil2022}, and integrated use of more than one model as
more practical than a single code \parencite{Singh2019b,Apostel2021}.
Machine-learning models raise their own validation requirements the literature
is only beginning to confront: hyperparameter tuning and diverse training data
for drain detection \parencite{Song2021}, neural networks needing training that
includes extreme situations and periodic retraining \parencite{Yang2000}, and
machine-learning surrogates trained across diverse climatic and hydrogeological
conditions with transient processes and real-time control
\parencite{Lei2026,Schneider2022}. Process-model deficiencies are named
directly---detailed site characterisation, added non-linear regression
dependences, and improved submodels for mineralisation, crop N uptake, and
growth under N limitation \parencite{Ahuja2007,Ikenberry2017}, with added
complexity the price of predicting snowmelt-thaw drainage timing
\parencite{Feyereisen2007}. Physical and analytical refinements form a rich set
of concrete proposals: numerical models incorporating anisotropy and varying
bottom slopes verified against analytical solutions \parencite{Lin2025,Kacimov2000},
reanalysis with the Laplace rather than Boussinesq equation \parencite{Bouarfa2000},
recession equations extended to spatially varying properties \parencite{Cooke2001},
three-dimensional treatment at heterogeneous sites \parencite{Mohanty1997},
distinct mixing-layer depths and partition coefficients for solute transport
\parencite{Shalit1996}, reconsidered deep-percolation and shallow-water-table
assumptions \parencite{Kolars2019}, and strengthened index-system construction
with additional soil indicators \parencite{Hou2018}.

\subsection{Instrumentation and measurement-method validation}

Because model validation is only as sound as the data feeding it, uncalibrated
flow structures and drifting sensors can masquerade as model inadequacy.
In-house calibration of flow-measurement structures is recommended repeatedly:
weir calibration for individual applications \parencite{Christianson2019,Napierala2023},
in-house calibrations representing field flow conditions for control structures
\parencite{Katuwal2024}, stage--discharge equations developed before monitoring
for cylindrical structures \parencite{BinShokrana2021}, and post-hoc sensor
calibration against grab samples during tracer tests \parencite{deOliveira2023b}.
Emerging technologies carry explicit validation requirements: a computer-vision
flow method needs testing under variable real-world conditions with the
Manning--Strickler equation to compensate for uneven velocities
\parencite{Young2023}, hybrid RTK--rotating-laser systems address
less-than-ideal survey conditions \parencite{Freeland2014}, and sensor design
and scintillometer humidity-stability effects require refinement
\parencite{Enamorado2007,Zhang2010}. Hydraulic-loss characterisation in drainage
conduits remains conspicuously incomplete, limiting design-equation realism:
head loss in corrugated-pipe junctions and manholes under free-surface flow
\parencite{Bayat2011}, friction loss and transitional cavity-flow conditions
\parencite{Nédélec2008}, declining-bend slope effects on air discharge
\parencite{Gurklys2008}, appropriate Manning coefficients for design
\parencite{Kouchakzadeh2011}, and strain-gauge characterisation of non-linear
HDPE-pipe behaviour \parencite{Gaj2021}. The consistent message across the
modelling frontier is that the next decade of progress depends less on new
equations than on sustained, well-instrumented, and openly reported field
campaigns needed to test the equations already in hand.

% ----- frontier/04_from-field-to-watershed-mapping-networks-and-sca.tex -----
\section{From Field to Watershed: Mapping Networks and Scaling Response}

\subsection{Overview}

Two linked spatial problems anchor this frontier. First, the physical drainage
network---where tiles lie, at what depth and spacing, and how mains connect
fields to ditches---is largely unmapped, so uncertainty about the network
itself propagates into every downstream hydrologic and water-quality model.
Second, even where field processes are well understood, the field is not the
scale at which water-quality goals are set; translating instrumented-plot
knowledge to the watersheds where nutrient and sediment targets apply remains a
genuinely open methodological problem. The two are related: because networks
are unknown and because routing, connectivity, and in-stream processes are
weakly represented, effectiveness estimates for practices such as drainage
water management and phosphorus mitigation cannot yet be aggregated with
confidence from field to basin. The priorities that recur across the literature
are therefore concrete---non-invasive detection methods with usable
practitioner guidelines, shared drainage databases with standardised reporting,
intensively monitored small catchments as upscaling bridges, and mechanistic
treatment of the channel-network, connectivity, and depressional processes that
mediate the field-to-watershed link.

\subsection{Locating, mapping, and cataloguing drainage networks}

Because the drainage network is often unknown, detection and mapping form a
fast-moving frontier in which methods work under some conditions but lack
systematic guidelines. Actual tile locations and measured discharge would
sharpen model calibration \parencite{Amado2018}, and predictive mapping from
imagery needs validation against GPS surveys plus methods that avoid
misidentifying noise and that reconnect broken segments
\parencite{Naz2009,Schilling2019}. The most explicit data agenda is
institutional: report tile diameter, depth, spacing, and mains, and build a
shared drainage database jointly maintained by agencies, researchers,
landowners, and contractors, with repositories updated as records accrue
\parencite{Zhou2025,Valayamkunnath2020}.

\subsubsection{Remote sensing and machine learning}

Remote-sensing work seeks urgent development of non-invasive methods for aging
systems, including UAV flights combined with ground-penetrating radar
\parencite{Wienken2024}, and multispectral, thermal-infrared, and
machine-learning agendas that define the site conditions under which drain
lines and outlets can be mapped and that yield decision-tree guidelines for
practitioners \parencite{Allred2020,Allred2021,Allred2022}. \textcite{Carlsen2024}
recommend reconsidering the often-cited 2.54~cm minimum rainfall for imagery
acquisition and investigating deep learning, multitemporal acquisition, and
synthetic aperture radar while expanding public drainage-map datasets---echoed
by calls for transformer-based architectures and higher-resolution imagery
integrated into hydrologic frameworks \parencite{Wan2024} and for reduced
tracing uncertainty through multiple interpreters, ground truth, and richer
covariates \parencite{Redoloza2023,Harris2018,Yimer2023}. Optical and radar
covariate studies further seek more frequent imaging, repeated SAR with ground
measurements, and UAV cover-crop mapping
\parencite{Tlapáková2017,Meunier2001,Xu2018}.

\subsubsection{Geophysics}

Geophysical evaluations remain condition-dependent and call for data from more
sites and soils \parencite{Allred2005}, systematic testing of GPR antenna
configuration and orientation \parencite{Allred2013}, and prediction of GPR
suitability via electromagnetic simulation driven by EMI. Concurrent EMI/GPR
surveys and drone or magnetic surveys are recommended, with performance shown to
depend on soil, residue, tillage, and wetness
\parencite{Koganti2020,Koganti2021,Jouen2018,Kratt2020,Allred2018b}. The
open question is less whether any single method can detect drains than how to
specify, in advance, which method suits a given soil, cover, and moisture
state.

\subsection{Scaling from field to watershed}

The defining open problem is the disconnect between field-scale process
knowledge and watershed-scale response. Quantifying the relationships between
field processes and watershed response is repeatedly identified as critical for
water-quality strategy, motivating explicit field-to-watershed linkage and
upscaling frameworks \parencite{Williams2015,Williams2017,Wang2020b}. Small,
intensively monitored catchments are proposed as the empirical bridge for
upscaling and model validation \parencite{Kyllmar2023,Basu2010}, and
effectiveness estimates for drainage water management and phosphorus mitigation
should be scaled up from field to watershed
\parencite{Wang2021,Ross2016,Shedekar2021}.

The hydrologic effects of tiling remain inadequately characterised at watershed
scale and must be pursued through combined modelling and field data: attending
to seasonal streamflow recessions in tiled versus non-tiled watersheds
\parencite{Schilling2008,Schilling2012}, to events of all sizes across
landscapes \parencite{Sloan2017}, to routing between spatially variable
depressional stores and flow accumulation across river networks
\parencite{Sloan2016,Fonley2021}, and to the decoupling of row-crop cover from
artificial drainage and of surface from subsurface contributions
\parencite{BolandBrien2014}. Streamflow prediction in tiled landscapes is
itself difficult, raising the unresolved question of whether curve-number or
more process-based methods are required
\parencite{Wessel2022,Thomas2016,Rahman2014}. Whether a representative
catchment scale transfers elsewhere must be established against additional
watersheds \parencite{Basu2010,Thomas2016}, and newly proposed physically based
tools such as the Advanced Terrestrial Simulator and HYDRUS-2D warrant scaling
and transferability tests across soils, topography, land-use change, and
climate \parencite{Rathore2025,Naderi2025,Basu2010}.

\subsection{Channel-network, connectivity, and drain-to-stream processes}

Interactions between drains and the receiving channel network are a frontier in
their own right. Open questions include the fate and transport of tile effluent
after it reaches surface drains \parencite{Clement2017}, drain--river and
drain--ditch interactions \parencite{Grangeon2021}, and field-scale monitoring
of how tile flow affects downstream discharge, stream power, and channel
geomorphology \parencite{Zhou2025,Zhou2013}.

Channel-network, in-stream, and sediment processes are underdeveloped in
watershed models. Networks should be modelled as interdependent reaches with
explicit flow paths and travel times \parencite{Giraud1997,Williamson2025};
in-stream biogeochemistry and hillslope--riparian interaction need stronger
representation \parencite{Schilling2012b,Basu2010}; and gully, streambank, and
non-field sediment sources require mechanistic treatment within catchment-scale
sediment budgets \parencite{Maalim2013,Dalzell2018,Walling2002}.

Connectivity metrics are increasingly flagged as management-relevant. Mapping
surface ditches may be more efficient than locating in-field drains, and future
phosphorus-risk models should incorporate ditch connectivity and lag times
\parencite{Moloney2020,Opoku2024}, supported by improved connectivity
algorithms and monitoring-based thresholds for lowland catchments
\parencite{Calsamiglia2018}. Finally, drained closed depressions and potholes
are singled out as poorly understood: they demand process study at additional
locations, quantification of runoff--tile--antecedent-moisture relationships,
and data on blind-inlet longevity
\parencite{Martin2019,Williams2020,Feyereisen2015b}. Across these threads, the
common gap is a lack of coupled, connectivity-aware representations linking
individual drains to the reaches and depressions that ultimately govern
watershed water quality.

% ----- frontier/05_managing-nutrient-losses-nitrogen-phosphorus-and.tex -----
\section{Managing Nutrient Losses: Nitrogen, Phosphorus, and Cropping Systems}

In-field nutrient management remains the primary lever for reducing nitrate
and phosphorus (P) losses from tile-drained land, yet the levers themselves are
better described than optimised. The frontier is less about which practices
help than about the metrics, mechanisms, crediting rules, and temporal targets
that would let managers set the right rate, form, and timing for a given soil,
climate, and cropping system. \textcite{Ma2025} frame the programme broadly,
arguing that far more work is needed to understand the tradeoffs among
productivity, nutrient-loss reduction, and ecosystem services in tile-drained
systems. This section synthesises the open questions for nitrogen, for
phosphorus, and for cropping-system change; model-algorithm gaps,
greenhouse-gas tradeoffs, edge-of-field treatment, and spatial bundling are
treated only in passing here and developed elsewhere in the review.

\subsection{Refining nitrogen management and the timing of loss}

Refining fertiliser nitrogen (N) is the most consistently identified priority.
\textcite{ROSSI1991} call for careful evaluation of fertiliser timing, rate,
form, and placement, and \textcite{Zhao2000} for strategies that raise N-use
efficiency in corn through optimal rate, timing, and form. \textcite{Randall2001}
add crediting of previous legumes and manure and improved manure management,
while \textcite{Dattamudi2020} conclude that application schemes need
reevaluation toward lower corn rates, precision agriculture, and split, properly
timed applications---echoing \textcite{Elmi2002}, who see reductions from
matching rate and timing to crop uptake. Optimising N rates is expected to
retain primary research and regulatory focus, with accurate 4R implementation
requiring site-specific knowledge to balance agronomic and environmental goals
\parencite{Christianson2015b}.

\subsubsection{Next-generation rate tools and decision support}

\textcite{Bandura2025} propose concrete standards for the next generation of
N-rate tools: developed against the economically optimum N rate (EONR),
adequate if they predict need within 10~kg~N~ha$^{-1}$ of EONR, and
environmentally unacceptable if they exceed it by more than 30~kg~N~ha$^{-1}$ at
an N:corn price ratio of 5.6. \textcite{Liu2021} call for recommendations that
weigh environmental and economic benefits, case-by-case cost analyses, and
coordinated agency action to promote 4R adoption, and \textcite{Banger2018} for
experiments across wider soils and management to identify yield tradeoffs from
rate reductions---shifting emphasis toward cover crops and edge-of-field
practices where yield-loss risk appears. Comprehensive tools integrating weather
forecasts, management history, and local field parameters, with regional
agronomic-optimum-N-rate benchmarks \parencite{Maas2022}, precision application
by soil-N availability \parencite{Gupta2023}, better crediting of non-fertiliser
N and variable-rate models incorporating remote sensing \parencite{Dinnes2002},
and low-cost proxy variables for soil N-supply capacity \parencite{Cassman2025}
are all advanced. Local decision approaches also need refinement: late-spring
nitrate-test programmes are framed as a start to be replaced as research emerges
\parencite{Jaynes2004}, with grain-protein monitoring and risk insurance for
adopters \parencite{Karlen2005} and weather-based tools to reduce soil-testing
needs \parencite{Malone2010}. \textcite{Jaynes2013} seek to pinpoint where
between the V12 and V16 stages sidedress N ceases to benefit, and
\textcite{Jaynes2015} note sidedressing can be paired with soil tests or crop
N-deficiency sensing for real-time adjustment; comprehensive volatilisation
studies and data from data-poor regions are further gaps \parencite{Chatterjee2020b}.

\subsubsection{Timing, form, and application method}

Method-level questions remain open. \textcite{Malone2007b} highlight quantifying
split versus single spring applications, and \textcite{HernandezRamirez2011}
note it is unknown whether reduced-rate fall manure could match sidedress
environmentally, pointing to nitrification inhibitors and more complex
rotations; \textcite{Elmi2005} stress quantifying all major N loss pathways
simultaneously. Product-level frontiers include enhanced-efficiency fertiliser
with a nitrification inhibitor tested alongside drainage water management
\parencite{Kaur2024c}, NH$_4^+$-N stabilisers and cover-crop effects on
NH$_4^+$-N dynamics \parencite{Vincent2025}, and fertigation of silage corn with
its crop-injury and pathogen tradeoffs \parencite{Gamble2018}. Evening out
manure application in high-livestock catchments \parencite{Bechmann2014} and
comparing N-use consequences under global-warming projections
\parencite{Haghnazari2020} round out this cluster.

\subsubsection{Crediting residual, mineralisable, and legacy nitrogen}

A recurring, more specific gap is crediting N supplied by soil organic matter,
prior legumes, manure, and land-use legacies. \textcite{Bakhsh2007,Bakhsh2005}
call for assessing proper credit from mineralisation of the organic-matter pool
and soybean residue before applications, and \textcite{Randall2001} for improved
soil-N tests to estimate mineralisable and carryover N---echoed by a spring test
after dry years \parencite{RANDALL1995} and pre-sowing mineral-N measurement
\parencite{Conrad2009}. Systems transitioning from perennial cover need
second-year credits after grass conversion \parencite{Huggins2001} and care
after alfalfa \parencite{ROBBINS1980}, while irrigated vulnerable soils call for
reduced early-year N with residual-nitrate crediting \parencite{Derby2009}.
Legacy N is itself underexplored: a clearer definition is needed to apportion
losses \parencite{Williams2022}, and the coupled biogeochemical and hydrological
processes regulating legacy-N formation and its effect on nitrate export require
dedicated study \parencite{Hu2024}. Reconsidering fall application under warmer
winters \parencite{Gentry2014}, resin-membrane strips to characterise soil-N
cycling \parencite{Masri2024}, and better long-term flow-weighted-N predictions
for benchmark soils \parencite{Skaggs2012} complete the agenda.

\subsubsection{The temporal window of loss}

A distinctive frontier is the temporal concentration of losses and the need for
management targeted to that window. \textcite{Bakhsh2007} argue strategies must
reduce early-spring nitrate leaching, and \textcite{Bakhsh2010} extend this to
March--June while flagging the temporal distribution of growing-season rainfall.
Assessment of nitrate status and tile loss during April, May, and June warrants
study \parencite{Nance2007}, and a winter-hardy cover crop ahead of corn with a
shift of fall N to spring may best cut loss in intensively tiled watersheds
\parencite{Gentry2024}. For outlet management, \textcite{Skaggs1995} recommend
selectively managing outlets to reduce peak nitrate losses, especially after
drought-reduced yields, and \textcite{Elmi2004,Elmi2002} note that keeping
controlled drainage operational in early spring and late fall---when losses are
severe and drainage is not needed for production---could achieve further gains.

\subsection{Phosphorus metrics, mechanisms, and transport}

Phosphorus in subsurface drainage is repeatedly flagged as under-researched
relative to N. \textcite{Hertzberger2019} call for greater effort on P transport,
and \textcite{Yuan2022} for research on conservation-practice effectiveness in
reducing dissolved-P loss and the causes of its annual variability.
\textcite{Pease2018} identify pedologic controls on P transport as key needs,
and \textcite{deOliveira2022} recommend refining loss frameworks through
within-field studies and a better understanding of the ``where'' and ``when'' of
drainage P loss. A first-order requirement is to resolve partitioning and
sources: \textcite{King2015} call for enhanced P routines and simultaneous
measurement of P losses in surface runoff and subsurface drainage to determine
partitioning, \textcite{Williams2016} for differentiating soil versus fertiliser
sources of tile-drained P, and \textcite{Qi2018} for measuring the size
distribution of particulate-P carriers.

\subsubsection{Beyond agronomic soil-test P: saturation and sorption metrics}

A consistent theme is that agronomic soil-test P (STP) is insufficient for
environmental risk. \textcite{Sharpley1995} argue that environmental tests
estimating bioavailable P and sorption capacity should replace agronomic tests
on high-risk soils. \textcite{Hooda1999} propose that degree of P saturation
(DPS) may better estimate loss potential than extractable P such as Olsen-P, and
\textcite{Nguyen2002} call for critical DPS values calibrated across sediment-P
tests and equilibrium P concentrations. \textcite{Osterholz2020} recommend
assessing fields in the 100--400~mg~kg$^{-1}$ STP range and examining sorption
saturation, macropore formation, and mineralogy, while P-adsorption capacity
should precede developing organic soils for cropping \parencite{MILLER1979}.
\textcite{McDowell2001} propose that P isotope-exchange kinetics could explain
and eventually predict loss across a wide range of STP.

\subsubsection{Desorption, colloidal, and particulate transport}

The sorption/desorption interface and particulate pathways are poorly resolved.
\textcite{Mumbi2025} regard quantifying when and how varying water sources and
fluxes desorb and adsorb P as critical to predicting subsurface DRP transport,
and \textcite{Flores2021} note the need to understand particulate settling and
P desorption. \textcite{Uusitalo2003} argue that mitigating bioavailable
particulate P requires reducing particle mobilisation across the whole field
surface and that controlling diffuse DRP is a longer-term task than reducing
erosion, with the availability of eroded P to algae still limited
\parencite{Ulén2007}. \textcite{Wilson2019} highlight quantifying the colloidal-P
contribution to DRP, and \textcite{Norgaard2013} recommend including both
dissolved and colloid-associated leaching parameters in field-scale risk
assessment. Sediment-P release under alternating waterlogged--drying regimes
\parencite{Nguyen2002}, P fractionation under polyacrylamide \parencite{Jiang2010},
and combining columns, plots, and catchments to link particulate loss to
macroporosity \parencite{Laubel1999} are further priorities.

\subsubsection{Management effects, events, and legacy phosphorus}

\textcite{Uusitalo2018} note that assessing the pace of P stratification needs
more data and leave open whether no-till adoption as a eutrophication measure is
justified. \textcite{Ni2020} call for studies of optimal STP under changing
climate to address STP--precipitation interactions, and \textcite{Wang2018} ask
whether pre- or post-application tillage decreases tile DRP and whether findings
extend to manure injection, with immediate post-application manure risk still
poorly understood \parencite{Madison2014}. \textcite{Brendel2019} recommend
designing practices that reduce P losses during frequent small events, and
\textcite{Schilling2017b} refining organic-P/total-P ratios and quantifying
wastewater-treatment contributions. \textcite{Jarvie2017} ask how much of a 40\%
reduction target 4R practices alone could achieve and warn that drawing down
legacy P may be increasingly costly and intrusive; cost-effective hierarchies
prioritising lower Olsen P and targeted alum \parencite{McDowell2015} and P
removal by grass harvesting \parencite{Zhang2015b} are advanced. Fragmented
reporting also hampers synthesis: \textcite{Christianson2016} call for field
studies documenting multiple P species and further developing the Drain Load
database, \textcite{Macrae2019} for understanding hydroclimatic event drivers
and tillage or placement effects, \textcite{Tanga2025} for monitoring
particulate and dissolved forms with site-specific stoplog depths,
\textcite{Zhang2017} for new soil-P testing under no-till, and \textcite{Spiese2023}
for novel practices addressing nitrate and DRP simultaneously.

\subsubsection{Phosphorus in organic soils}

Phosphorus in cultivated organic soils is singled out as poorly understood.
\textcite{Grenon2022} call for mechanistic studies of the pathways, hot spots,
and hot moments of P transfer from agricultural organic soils, and
\textcite{Grenon2021} propose a concrete agenda using isotope tracing, XANES,
and $^{31}$P~NMR alongside sequential fractionation, microbial-turnover
assessment, and long-term drainage-water-management studies. Nutrient leaching
from abandoned peat fields, and the true nature of H$_2$SO$_4$-P, remain poorly
studied \parencite{Höyhtyä2025}.

\subsection{Cover crops, perennials, and rotational diversification}

\subsubsection{Cover crops: agronomic, environmental, and modelling frontiers}

Cover crops---especially winter rye---attract the largest concentration of open
questions. \textcite{Kladivko2014} call for research and technical assistance to
overcome adoption obstacles and for field experiments across more Midwestern
soils, and \textcite{Birru2023,Zhao2016} stress that existing evidence is drawn
from too few, geographically clustered sites and must be extended across more
species, soils, and climates. Agronomic priorities include tools to promote rye
in colder regions \parencite{Waring2020}, establishment strategies under
temperature and precipitation extremes \parencite{Bawa2023}, evaluation of more
cultivars to reduce yield variability \parencite{Kaspar2007}, breeding for cold
hardiness \parencite{Dinnes2002,Liu2019}, whether yield drag reflects poor
planter performance through residue \parencite{Dougherty2020}, and better
management of kura clover as living mulch \parencite{Qi2011b,Qi2011}. On
environmental mechanisms, authors seek research on the spring release of
scavenged N and its synchronisation with crop growth \parencite{Lacey2015},
adjustable seeding rates \parencite{Roth2018}, combined cover-crop and manure
effects on subsurface DRP \parencite{Askar2023}, and optimisation of cover-crop
services against disservices to the water balance \parencite{Meyer2019}.
\textcite{Liu2019} call for coordinated common experiments to quantify
winter-cover P contributions, longer-term studies remain scarce
\parencite{Kaspar2012,MartinezFeria2016}, growing cover crops longer in spring
may prevent N$_2$O increases \parencite{PrezaFontes2023}, and because over-winter
effects differ for N versus P, practices should target individual pathways
\parencite{McNeill2026}---with \textcite{Krueger2013} recommending practices
beyond rye where subsurface nitrate stayed high and rye had no effect on DRP.
Cover-crop modelling is a distinct frontier: validated rye growth and plant-N-
uptake algorithms \parencite{Feyereisen2006,Malone2007c}, dedicated winter-rye
modules and calibration \parencite{Malone2014b,Chatterjee2025}, and correction
of cover-crop and frozen-soil weaknesses \parencite{Dietzel2016,Gupta2022,Gupta2023b}
are sought, with detailed algorithm gaps developed in the modelling section.

\subsubsection{Perennials, bioenergy, and rotational diversification}

Perennial and land-use-change options feature prominently. Poorly drained
depressions are proposed for perennial bioenergy crops such as
\textit{Miscanthus $\times$ giganteus} or wetland restoration \parencite{Huang2023},
and keeping soil covered year-round with perennial ley is judged the most
effective way to mitigate N losses, with the P focus on decreasing surface
runoff \parencite{Norberg2022}. Additional data on alfalfa's effect on subsurface
discharge during large events \parencite{Arrueta2025}, a clear causal link
between perennial crops and nutrient loads with policies rewarding environmental
services \parencite{Arrueta2022}, robust water-quality datasets for perennialised
systems \parencite{Wilson2023}, and large-scale bioenergy-scenario assessment
with cost--benefit analysis to alleviate Gulf hypoxia
\parencite{Guo2018,Woodbury2018} extend the agenda. Rotation and diversification
questions recur: accounting for losses on a rotation rather than annual basis
\parencite{Zhao2016}, the paucity of research on desirable rotation crops in
two-crop systems \parencite{Pasley2021}, partial conversion of cropland to
pasture or wetland to meet nitrate targets \parencite{Randall2001,Gowda2007,Nangia2008},
and explicit joint evaluation of controlled drainage, rotation, cover crops, and
N management \parencite{Ma2007b}. Novel options---triticale with red clover
\parencite{Nance2007}, camelina--soybean relay-cropping \parencite{Emmett2022},
and fertilised, harvested rye for revenue \parencite{Malone2018,Malone2023}, and
diverse systems benchmarked against the two-year rotation \parencite{Gentry2025}
---are all advanced. \textcite{Dinnes2002} call broadly for perennial cash-cropping
systems that maintain returns while reducing leaching, and \textcite{Jaynes2001}
conclude that dramatic management change on artificially drained soils will be
required before corn production can be considered sustainable.

% ----- frontier/06_water-table-management-and-drainage-system-desig.tex -----
\section{Water-Table Management and Drainage-System Design}

The frontier in drainage water-quantity management is defined by a single
shift: from passive, one-way conveyance toward active, location-specific
control of the water table and the coupled irrigation--drainage balance. The
enabling ideas---water-table-control structures, automation, forecast-based
operation, and re-evaluated physical design---are largely in place, but each
remains a hypothesis awaiting the long-term, multi-site field and simulation
evidence needed to quantify its yield and water-quality trade-offs. What
follows organises these open questions around active water-table control, the
control systems that would automate it, and the design criteria that would
support it.

\subsection{Controlled drainage, subirrigation, and drainage water management}

Integrating controlled drainage (CD) with irrigation is repeatedly framed less
as settled practice than as a research frontier: the recurring call is to move
from passive, ``unilateral'' drainage toward active control of the water table
and the coupled irrigation--drainage balance
\parencite{Ayars2006b,Kesicka2023b,Ritzema2023}. \textcite{Ayars2006b} argue
that water-table-control structures should be required in new designs and
considered for retrofits, and that new criteria and management methods are
needed---claims that require long-term field studies, possibly combined with
modelling, to validate---echoed by calls to build in control options and expand
changeable CD networks \parencite{Dukhovny2018,Mahmoud2021,Strock2011}. The
quantitative payoff of active management remains open: \textcite{Watkins2024}
hypothesise that active outlet management during extremely wet periods could
minimise wet-stress and raise yields; \textcite{Saadat2017} emphasise the
cost--benefit trade-off; and \textcite{Helmers2012} suggest that CD yield
reductions could be mitigated by greater management while weighing installation
cost, management burden, and applicable acreage. Seasonally differentiated
control---free drainage in wet seasons, CD in dry seasons, with supplemental
irrigation for extremes---is an explicit hypothesis to test
\parencite{Asgari2021}.

Coordinating irrigation with drainage control is a persistent unsolved problem:
irrigation must be reduced when CD is implemented, requiring coordination among
authorities and farmers \parencite{Wahba2005,Ritzema2011,Ritzema2016}, and CD
requires genuinely integrated management with careful site selection
\parencite{Dukhovny2018}. Subirrigation embedded in regional water systems is a
distinct, largely unresolved frontier: regional water needs should be known
before responsible incorporation, with optimised pumping and cautious use of
alternative resources \parencite{deWit2022}, and SWAP-type modelling can
pre-screen sites and forecast consequences before implementation
\parencite{deWit2024}. Seepage effects on CD performance need coordinated
experimental and modelling research, including guard lines to evaluate lateral
seepage \parencite{Youssef2021,Pease2017,Ale2009}, while excess-moisture stress,
spring drawdown, and spring nitrogen loading remain priorities
\parencite{Craft2018,Lahdou2019}. Finally, longer-term studies must disentangle
whether soil-property changes under drainage--subirrigation arise from the
practice itself or from irrigation-water quality \parencite{Kaur2023}, and
complete water and nutrient budgets within and across field boundaries---plus
effects on runoff, preferential flow, and recharge---remain to be tracked
\parencite{Shedekar2021,Ross2016,King2022}.

\subsection{Automation, forecast-based, and data-driven control}

Automation and data-driven control are identified as the next technological
step, but the control systems themselves are still to be built and validated.
\textcite{Youssef2023} call for a new generation of control structures and
decision-support tools, sensor-based systems tailored to shallow-water-table
soils, and systems that adapt to changing weather; \textcite{Ayars2015} extend
this to adaptive-control decision support. Concrete methodological agendas
include dynamic modelling and data assimilation of the field water balance,
feasible formulations of the control problem, and fusion of control automation
with physical infrastructure \parencite{Ikonen2023,Isomäki2024}, automatic
water-depth control for aerobic rice \parencite{Arif2021}, and---anticipated
decades earlier---farm decision aids that jointly incorporate water quantity,
quality, crop growth, and economics \parencite{THOMAS1995}.

Anticipatory, forecast-based operation is the emerging synthesis of these
threads. Forecast-driven management should incorporate evapotranspiration,
subirrigation, and spatially variable seepage, and must explicitly address
forecast uncertainty \parencite{VandeCraats2021}; collector-ditch water levels
and control schedules should be optimised \parencite{Salla2022,Strock2025b};
and control should be paired with local crop and climate characteristics
without overlooking the risk that rapid water-table decrease raises irrigation
demand in dry summers
\parencite{Güngör2023,Dou2023,Golmohammadi2016,Kim2023,Wang2020c}. The open
questions are thus as much about coordinating forecasts, sensors, and structures
into a single adaptive system as about any individual component.

\subsection{Novel structures, design criteria, and drainage intensity}

Novel structures and materials define a first set of open questions: improved
HYDRUS-2D boundary conditions and field verification of siphon drainage
\parencite{Ren2023}, long-term measurement of working time against construction
cost \parencite{Tao2025b}, and new structures that adjust water-holding capacity
while inhibiting rapid vertical and bypass flow \parencite{Li2021}. Multi-level
and bi-level configurations warrant study of actively managed shallow and deep
drains \parencite{Hornbuckle2007,Verma1998,Nelson2011b}, and envelope,
perforation, and screen design remain unresolved across gradients, soils, and
textures \parencite{Alavi2021,Wu2022,Oyarce2017,HajAmor2019}. Optimal layout and
burial-depth determination in saline and problem soils recur in recent studies,
which stress ditching depth, restrictive-layer breaking, amendments, and
construction cost within multi-scheme decision models
\parencite{Tian2023,Tian2024,Tian2023c,Zhang2024d,Zhai2023,Han2023,Qian2025b,Ren2021,Zheng2025,SerranoCoronel2018,Zegeye2021}.
Analytical advances are also sought---formulas for sloped and multilayer soils
\parencite{Tao2016,Tao2019}, partially penetrating and trapezoidal ditches
\parencite{Chahar2010,Afruzi2014}, and mass-transfer and three-dimensional flow
effects \parencite{Zhang2024}---alongside CAD/GIS-based design software and
standardised drainage-coefficient definitions to enable comparison
\parencite{GALLICHAND1992b,Skaggs2017}.

Many site studies conclude that established design criteria themselves require
re-evaluation. New criteria that protect crops while minimising drain-water
salinity and volume are sought \parencite{Christen2001}, along with updated
depths and spacings at specific sites
\parencite{Muhammad2021,Malota2016,Bahçeci2009,Kahlown2002}, review of
recommended spacings paired with controlled-traffic farming
\parencite{Mante2022,Mante2018}, cost-effectiveness of narrow spacings in
claypan soils \parencite{Nelson2017}, and standard depth-and-spacing guidelines
tested in situ before extension \parencite{HajAmor2019}. Design should further
account for changing crop water needs, climate, and irrigation modernisation
\parencite{Rui2024,Akram2013,Ziaeifar2025,Ritzema2023,Ritzema2007,Ritzema2006,Ritzema2011}.

Drainage intensity emerges as a systems question rather than a per-field
specification. Reduced spacing would require increased network capacity and
heighten reliance on best management practices
\parencite{Singh2009,Dialameh2025}; a systems approach should integrate
drainage, crop--soil processes, and nutrient-loss reduction
\parencite{Castellano2019}; and DRAINMOD-based analysis should develop optimum
rates and spacings across locations \parencite{Strock2011}. Which soil
parameters are most critical, and what confidence intervals each requires,
remain fundamental unanswered questions
\parencite{Ayars2015,SHIRMOHAMMADI1995,HARDJOAMIDJOJO1982}. Underlying all of
these is a methodological prerequisite: design guidance must rest on long
records, since drainage design should not rely on short-term simulations
\parencite{Wang2006b}, and long runs with hourly input are needed to assess
wider spacings, interannual variability, and shallow-water-table control
\parencite{Ale2009,Häggblom2019,Malota2015,Salo2019,Kladivko2023,Chang2012}.

% ----- frontier/07_edge-of-field-treatment-technologies.tex -----
\section{Edge-of-Field Treatment Technologies}

\subsection{Overview}

Edge-of-field treatment structures---denitrifying woodchip bioreactors and
their companions (constructed wetlands, saturated riparian buffers, and
phosphorus-removal structures)---have accumulated abundant performance
demonstrations yet remain defined more by open design questions than by settled
rules. Bioreactors remove nitrate by microbially mediated denitrification, but
their performance is bounded by low water temperature, variable hydraulic
retention time (HRT), and the biogeochemical evolution of the woodchip medium as
it ages. The frontier clusters into interlocking lines: how to size, scale, and
standardise design; how to resolve the internal hydraulics and tracer methods on
which design rests; whether the microbial engine can be deliberately enhanced;
how to select, amend, and manage reactive and phosphorus-sorbing media; how to
characterise multi-decadal aging, clogging, and construction effects; how
pumped and coupled configurations extend the design space; and how companion
practices complement or substitute for bioreactors. A recurring meta-message is
the need for standardised design procedures, broader performance-metric
reporting, and long-term data across climates, landscapes, and maturation stages
so that watershed-scale nutrient-reduction goals can be met
\parencite{Christianson2021,Maxwell2022,Bell2015}. The evidence is drawn
overwhelmingly from site-scale field and laboratory studies concentrated in the
U.S.\ Midwest; the by-products of reductive treatment (greenhouse gases,
dissolved organic carbon export, methylmercury) are treated here only as design
constraints, their downstream accounting deferred to the consequences synthesis.

\subsection{Sizing, scaling, and standardised design}

A central frontier is reconciling site-level effectiveness with watershed-scale
load removal. \textcite{Christianson2026} argue for better balancing bioreactors
that are as effective as practical at the site while still removing as much load
as possible; they flag the untested performance of \emph{undersized} bioreactors
serving larger drainage areas, propose greater design flexibility to scale the
practice, and caution that additional empirical data are needed to strengthen the
loading-density relationship and that performance data from full-size bioreactors
beyond Illinois are required before prediction models are used for design. For
hydrologically flashy sites, \textcite{Dougherty2020b} propose sizing designs
upward and repeatedly call for simpler, standardised ``plug-and-play'' design
procedures; standardisation is echoed at the policy level by
\textcite{Christianson2021}, who recommend refining the Conservation Practice
Standard and streamlining the design process. The adequacy of residence-time
($T$)-centric design is itself questioned: \textcite{deOliveira2023} advocate
field-scale assessment of baffles in wide bioreactors to raise hydraulic loading
while preserving efficiency, and more tracer testing of non-standard media
coupled to nitrate monitoring and modelling. Foundational validation remains
incomplete---further field-scale validation of the design procedure against beds
of differing design \parencite{Christianson2013}, validation of hydraulic models
at additional denitrification-bed sites with conservative isotopic tracers
\parencite{Ghane2015}, region-specific parameterisation of flow, seasonality, and
microbial limits for new settings such as Lithuania \parencite{Povilaitis2020},
and sorption experiments underpinning tracer-based design \parencite{Ghane2019}.

\subsection{Hydraulics, retention time, and tracer methods}

HRT is repeatedly identified as the master variable linking hydraulics to
removal, yet its treatment in design guidance is imperfect. \textcite{Moorman2015}
note that research on the HRT--nitrate-removal relationship would sharpen sizing
criteria, and \textcite{Povilaitis2018} call for proper optimisation of HRT
effects through both laboratory- and field-scale study across flow conditions. In
a global synthesis, \textcite{Christianson2025c} recommend methodological
advances: assessing varying HRT by changing flow rate versus sampling along the
reactor length; greater transparency in estimating woodchip porosity; more
consistent reporting of dissolved oxygen, pH, and dissolved organic carbon; and
machine-learning refinement of HRT--removal-efficiency estimates as data
accumulate. The internal hydraulics of woodchip media are poorly resolved:
\textcite{Ghane2016} argue that the non-Darcy flow characteristics of woodchips
should be exploited to revise NRCS design guidelines. Tracer methodology is a
specific bottleneck. \textcite{Ghane2019b} report anomalously low bromide
recovery and recommend testing for organobromine in woodchips, microscopic and
excavation-based examination of physical bromide entrapment, and either
accounting for retardation or adopting non-sorbing conservative tracers such as
isotopes; \textcite{Ghane2018} likewise recommend tracer tests to compare the
hydraulic performance of retrofitted beds. Reliable tracer methods are thus a
prerequisite to benchmarking hydraulic performance at all.

\subsection{Enhancing the microbial engine}

The defining enhancement frontier is whether denitrifiers can be inoculated into
woodchip bioreactors to raise removal where cold temperatures limit biological
activity. \textcite{Jang2019} propose that bioaugmentation might increase
low-temperature removal but note that genome completion is needed to identify the
nitric-oxide reductase of a candidate \emph{Cellulomonas} inoculant.
\textcite{Anderson2020} frame the first-order uncertainties: whether inoculated
strains enhance field removal, how long they survive in bioreactors, and whether
carbon addition boosts both inoculated and indigenous denitrifiers. Survival and
persistence recur globally, with \textcite{Wang2025b} calling for field-scale
methods that minimise inoculum loss and for exploration of emerging targets and
consortia---nitrate-dependent anaerobic methane oxidation (n-DAMO), biostimulation
to favour inoculated bacteria, and algal--bacterial granular sludge. Terminal
denitrification steps in fungal strains, including confirmation of nitrous-oxide
reductase and reduction of nitrate to dinitrogen under bioreactor conditions,
remain to be verified \parencite{Aldossari2021}, and microbial amendments require
replicated pilot-scale testing before upscaling
\parencite{Jeglot2022,Ma2021}. Against this enthusiasm,
\textcite{Hathaway2015} offer a testable counter-hypothesis---that optimisation is
more likely from refinements to design and management parameters that shape
environmental conditions than from direct microbial manipulation---a useful null
hypothesis for field trials.

Rational manipulation presupposes mechanistic understanding of the resident
community, which is incomplete. \textcite{Hartfiel2022} argue that bioreactor
microbiome study should draw on soil and wetland ecology; that controls on
cellulose- and lignin-metabolism genes warrant attention; that dissimilatory
nitrate reduction to ammonium and anammox have been observed but understudied;
and that varying substrates to select for differing communities is a promising
route to more consistent performance and reduced by-product formation---the
mechanism by which HRT shapes the community, and thereby by-product production,
being a specific gap. Complementary needs include managing bed pH for acidic
drainage and conceptualising regional controls on in-bed nitrogen removal
\parencite{Bock2018}, routine monitoring of nitrite to avoid its accumulation in
effluent \parencite{Kouanda2021}, and characterising the selective pressure of
oxygen versus nitrate load on community assembly \parencite{Schaefer2021}. The
emissions accounting of these by-products is treated in the consequences
synthesis; here they define a design and monitoring constraint.

\subsection{Reactive media, carbon sources, and structural materials}

Beyond conventional woodchips, alternative fill and carbon-management strategies
are incompletely characterised. \textcite{Law2023} encourage field studies of
corncobs capturing spatial and temporal variability, the nitrogen-removal-rate
range, and whether a lower maximum HRT is needed, while
\textcite{Porter2015} call for study of organic-carbon availability at different
depths and moisture contents, so that substrate choice or internal water levels
could be tuned. Because woodchip carbon fuels denitrification, its release
kinetics and pre-treatment are active frontiers: the longevity of
composting effects and the biodegradability and release rates of woodchip DOC
remain to be quantified \parencite{Kouanda2021}, and startup carbon losses raise
questions about recycling the initial flush or pre-washing chips before
installation \parencite{Abusallout2017}. Synthetic-carbon dosing is unresolved---
the source, frequency, and amount needed to optimise performance under dynamic
inflow while avoiding excess DOC and biofilm clogging \parencite{Rivas2020}, and
the hybrid application of synthetic carbon in low-temperature-tolerant reactors
\parencite{Zhang2021b}.

A distinct materials frontier concerns drain envelopes and structural fill.
Rice-husk envelopes require field testing of combination ratios and long-term
assessment given husk degradation and acidification
\parencite{Kaboosi2012,Navabian2017}; aggregate sizing needs in-situ testing in
clay soils \parencite{Byrne2023}; and synthetic geotextile envelopes must be
evaluated individually because of varying physical properties, with
filter-sock-plus-aggregate combinations not yet recommendable
\parencite{Byrne2023b}. Structural concrete alternatives---fly-ash- and
coal-gangue-based concrete---remain untested for durability under high
groundwater \parencite{Wang2016c}.

\subsection{Phosphorus-removal media and filter design}

Phosphorus-sorption materials (PSMs) constitute one of the most heavily populated
gap areas, spanning selection, longevity, regeneration, and hydraulic
integration. \textcite{Li2020} argue that adsorbent choice must weigh phosphorus
concentration, interfering ions, temperature, flow rate, and sediment load;
prioritise selectivity at neutral pH; be testable for reuse and recovery; and be
economical at field scale. \textcite{Mendes2022} add that the apparently
synergistic retention mechanism of filter media at high flow rates is unknown and
that regeneration of spent media should be pursued. Iron- and steel-based media
dominate specific proposals: optimising an iron-oxyhydroxide regeneration process
and evaluating low-cost goethite such as bog iron ore \parencite{Allred2017};
determining minimum contact time for zero-valent iron via higher-flow column and
isotherm tests \parencite{Allred2014}; regenerating steel by-products by alkaline
washing \parencite{Hua2016}; connecting a steel-slag filter to a bioreactor only
after the initial DOC flush and mixing fine steel to preserve conductivity
\parencite{Sellner2019b}; optimising removal, retention time, slag mass, and cost
\parencite{Wang2020d}; controlling sediment influx and improving flow
distribution \parencite{Shedekar2020}; and recognising that ligand-exchange PSMs
suit tile drainage only where structures allow high-frequency media replacement
(every 4--6 months) \parencite{Penn2020}. Fly-ash and pellet media raise
saturation-timing questions---saturation time must be quantified before field use
so pellets are replaced before shifting from phosphorus sink to source
\parencite{Li2017,Li2018b}---and filter blends remain unoptimised with respect to
media configuration, contact time, and high-flow contact
\parencite{Agrawal2011,Agrawal2011b,Kumari2025,Kankarej2023}. Integrated
woodchip-plus-PSM designs need multi-scale, multiyear evaluation of removal
effectiveness, media placement, lifespan, and regeneration \parencite{Husk2018}.

Novel phosphorus dynamics are emerging. \textcite{Zhang2024b} call for
experiments on enhanced orthophosphate uptake by woodchips undergoing
anoxic--oxic cycling and identification of polyphosphate-accumulating organisms,
while redox-driven phosphorus release remains a monitoring gap requiring start-up
flush and aging assessment with high-frequency dissolved-reactive-phosphorus
(DRP) monitoring and drawdown plates \parencite{Bailon2025}, multi-season
confirmation of DRP removal \parencite{Rivas2020}, more reactive passive-sampling
resins \parencite{Li2020b}, and characterisation of dissolved organic phosphorus
lability and bioavailability \parencite{Liu2024b}.

\subsection{Aging, longevity, clogging, subsidence, and construction}

The multi-decadal service life and physical stability of treatment structures are
among the least-studied dimensions. \textcite{Ghane2018} recommend longer study
of woodchip particle-size reduction, investigation of whether woodchip carbon
becomes recalcitrant to the point of inefficiency before particle-size reduction
throttles flow, and studies beyond ten years. \textcite{Christianson2020}
similarly call for long-term documentation of biofilm formation, siltation, wood
decay, and the effects of a decreasing carbon-to-nitrogen ratio and increasing
lignocellulose index, to inform woodchip-recharge thresholds; long-term
($>$8~yr) field programmes were underway to establish efficacy
\parencite{Greenan2009}, and spatiotemporal removal should be evaluated across
normal, wet, and drought years alongside woodchip C:N measurements
\parencite{Hartfiel2023b}. Estimating replacement interval and cost depends on
performance, lifespan, and residue cost, with pre-measured feedstock volumes
proposed to control construction-method effects \parencite{Kohn2023}. Chemical
clogging is mechanistically unresolved and needs simulation models for prevention
\parencite{Guo2020,Guo2021}. Subsidence monitoring is a new methodological
frontier: \textcite{Christianson2025b} propose applying a LiDAR method at
additional sites---especially immediately after construction---to build
subsidence time series and assess targeted inlet recharge. Sediment management at
construction is itself a design frontier: \textcite{Feyereisen2023} plan to weigh
the opposing effects of bypass flow and nitrate removal in a sediment-exclusion
system and advise pre-construction investigation of sedimentation sources.
Network-scale diagnostics also need extension---higher-pressure air injection and
time-domain electromagnetic resistivity monitoring for large drainage networks
\parencite{Henine2021}.

\subsection{Pumped, energy-assisted, and coupled configurations}

Pumping to raise hydraulic loading is an emerging but incompletely validated
avenue. \textcite{DíazGarcía2025} call for assessing different pumping
configurations, note that continuous 24-h operation or a longer pumping window
could increase treated flow, and caution that pumping may suit only the most
appropriate applications. \textcite{Hartfiel2024} advocate solar-powered pumped
bioreactors to reduce operating and unit cost, and \textcite{Hartfiel2023}
recommend sizing pumping to a bioreactor's actual rather than design capacity,
timing pumping to source-water nitrate concentration, temperature, and carbon
content, and extending applications to nitrate-laden surface waters, irrigation
waters, ditches, cisterns, and groundwater. Management experiments---partial
inlet-recharge strategies and dual-chamber operation engaging a second chamber
only under high flow---are proposed \parencite{Hartfiel2023b}, and pumped
constructed-wetland-plus-woodchip-bioreactor systems require analysis of
particulate versus dissolved nutrient fractions during intense discharges and
longer monitoring of the pumping effect on phosphorus retention
\parencite{Pugliese2020}. Novel source waters and enhanced configurations broaden
the space: denitrification responses to spring flow and concentration, with
payment-based incentives for spring bioreactors \parencite{Easton2019};
low-temperature bioelectrochemical reactors with horizontal flow, baffled
oxidising/reducing zones, and plate cathodes \parencite{Law2018}; and coupled
phosphorus-removal structures sited at critical source areas, including a proposed
redox-coupled Fe(II)--nitrate reaction for nitrate removal by activated aluminium
\parencite{Penn2024}.

\subsection{Wetlands, saturated buffers, and companion structures}

Frontier questions extend to practices deployed alongside or instead of
bioreactors. For saturated riparian buffers (SRBs), a key challenge is predicting
nitrate removal well enough to design more effective buffers, motivating
unsteady-flow and numerical-transport methods and optimal-width approximations
\parencite{McEachran2020}; replicating SRB results across landscapes, with
biomass harvesting to sustain viability \parencite{Jaynes2014}; more
sophisticated removal and hydrology components coupled to economic objective
functions, and determination of maximum allowable buffer length before optimising
width \parencite{Abdalaal2023}; confirmation that removal increases as vegetation
matures, with multipipe configurations and additional monitored sites
\parencite{Jaynes2019,Jaynes2019b,Singh2024b}; and paired-field determination of
actual performance and the cost-effectiveness of the buffer-flow and
controlled-drainage components \parencite{Ghane2025}.

Constructed and treatment wetlands present persistent performance-limit
questions. Hourly-time-step dynamic modelling is needed to approximate nitrate
removal in tile-fed wetlands with means to buffer and gradually release flows
\parencite{Tanner2011}, and long-term study is required of sustainable
nitrogen-retention efficiency across climates, whether organic-carbon
availability limits denitrification, the value of biomass cutting or slow-release
carbon supplements, and flow-regime effects analysed through residence-time
distributions \parencite{Tanner2005,Tanner2003,Tanner2005b}. A broad agenda for
constructed treatment systems spans driving factors, flow--treatment
relationships, design-parameter effects, long-term field tests, and
winter-efficiency enhancement \parencite{Mendes2021}, alongside maturation of
constructed wetlands for phosphorus retention \parencite{Mendes2020}, siting and
sizing by wetland-to-catchment ratio and seasonal discharge
\parencite{Steidl2019,Tournebize2017}, and subtropical adaptation of
land-treatment (FILTER) systems \parencite{Gardner2001}.

Denitrification-bioreactor and combined-installation refinement remains open:
reactor size, packing materials, elimination of preferential flow, and pH
maintenance \parencite{Aide2016}; separating the contributions of controlled
drainage with and without a bioreactor and evaluating clogging potential
\parencite{Husk2017}; and verification of design retention times and hydraulic
conductivity \parencite{Feyereisen2015}. Controlled and flow-regulating
structures raise further questions---improving the connection of regulator
elements to limit effective flow sizes \parencite{Popek2021}, pipe-riser and
event-scale trapping effects in water-and-sediment control basins
\parencite{Gupta2019}, and soil-surface clogging, storage geometry, and
discharge-capacity effects in drainage-inlet flow control \parencite{Nédélec2010}.
Finally, subsurface attenuation adjacent to drainage needs broader
characterisation: hyporheic-zone processes and buffer widths sufficient to
guarantee adequate residence times across hydrogeologic settings
\parencite{Puckett2005}, low-rate effluent irrigation to reduce direct transfers
through mole-pipe drainage \parencite{Monaghan2007}, and long-term monitoring of
tillage and earthworm effects on phosphorus leaching \parencite{Ulén2018}.
Collectively, these companion-practice frontiers reinforce the section's central
message: moving edge-of-field treatment from site-scale demonstration toward
validated, transferable design guidance capable of meeting watershed-scale goals
across diverse climates, soils, and hydrologic regimes.

% ----- frontier/08_salinity-reclamation-and-drainage-water-reuse-in.tex -----
\section{Salinity, Reclamation, and Drainage-Water Reuse in Arid and Irrigated Systems}

Subsurface drainage is a cornerstone technology for reclaiming salt-affected
land, yet its salt-leaching outcomes accrue slowly, are only partially
observable with current programmes, and generate salt-laden water whose disposal
remains unsolved. The frontier here is defined by this triad---slow and
incompletely measured reclamation, the management of the water that reclamation
produces, and the specialised chemistry of acid sulfate and irrigation-return
systems---together with the sensing, modelling, and design tools needed to
resolve it.

\subsection{Salt-leaching dynamics and long-term reclamation trajectories}

The recurring message is temporal: meaningful salt reduction often accrues so
slowly that short observation windows systematically underestimate long-term
benefit, and only extended drainage with sustained monitoring resolves the true
trajectory \parencite{Almen2021,Daigh2025}. The long-term effects of combined
saline-water irrigation and subsurface drainage on both salinity and yield
require further evaluation \parencite{Feng2017,Feng2021}, desalination rates and
trends remain incompletely quantified \parencite{Guo2024b,Yao2024}, and
controlled drainage introduces unresolved temporal trade-offs, buffering
outflow while risking salt accumulation
\parencite{Jouni2018,Lu2019,Islam2025,Heng2021}. Longer observation records
paired with optimal leaching quotas and cost--benefit analysis are needed to
close these gaps \parencite{Guo2025}, and salt-tolerant crops in reclamation
rotations, alongside quantification of rainfall and irrigation-efficiency gains
relative to salt-removal capacity, remain to be characterised
\parencite{Qadir2001,Castanheira2010}.

A parallel gap is measuring what actually leaves through the drains. Priorities
include greater attention to surface-water monitoring of salt and fertiliser
export \parencite{Mastrocicco2013}, continuous flow meters and
electrical-conductivity probes for oasis-scale balances with evaluation of
impacts on downstream phreatic groundwater \parencite{Askri2022}, nutrient
balances of recipient waters \parencite{Turtola1995}, and cluster wells with
deep sampling to track salinity over time \parencite{Singh2020d,Gafni2001}.
Process understanding of salt mobilisation is itself an active frontier: how
macropores govern transient-state leaching and how their effects might be
minimised \parencite{Xin2016}, the controlling factors in low-permeability and
structured media \parencite{Zhang2022}, salt replacement between pipes under
brackish-water leaching \parencite{Li2023}, rainfall--irrigation interactions
and generalisable desalinisation ratios across climates and soils
\parencite{Yu2015,Chen2025,Wang2024d}, ponding strategies and leaching curves
\parencite{RAO1991,Yang2022c}, and the optimal root-zone matric potential that
links leaching intensity to crop response \parencite{Zhang2021,Heng2021}.
Overarching these, waterlogging and salinity should be assessed together rather
than separately \parencite{Valipour2014}, and the economic feasibility of
drainage as a preventive measure against salinisation remains to be established
\parencite{Datta2000,HalbacCotoaraZamfir2022}.

\subsection{Drainage-water reuse, disposal, and treatment}

Removing salt from soil relocates the problem to the drainage water, whose
management is among the most persistent open challenges. Environmentally
acceptable disposal methods are still needed \parencite{Ayars2014,Steppuhn2006},
evaporation-pond management must be tested at full scale and combined with
volume- and contamination-reduction measures \parencite{BRADFORD1991}, and
trace-element fate---particularly arsenic speciation and uranium recovery or
removal---compounds the disposal burden \parencite{Gao2007,Duff1997,Gao2003}.
Emerging treatment routes remain at the demonstration stage and require testing
under realistic conditions: electrolysis of saline wastewater at varying
salinity \parencite{Yang2023}, sorptive filter materials evaluated with real
rather than artificial drainage water \parencite{Pugliese2023}, and blending of
drainage water with reverse-osmosis product water \parencite{Singh2020d}.
Life-cycle cost assessment, scalable data infrastructure, and low-temperature
crystallisation toward near-zero discharge define the longer-horizon disposal
agenda \parencite{Qin2025}.

Reuse research spans site-specific load--flow relationships, germination-stage
salinity tolerance, and improved scheduling that meets crop evapotranspiration
from a shallow water table while minimising upward salt flow
\parencite{WESTCOT1988}. Additional needs include provisional field guidance for
substituting drainage water on wheat \parencite{Sharma1998}, blending combined
with evaporation-pond storage \parencite{Datta2002}, farm-level reuse embedded in
a basin-wide salt balance \parencite{Qureshi2008}, exploitation of seasonal
groundwater undulation and drip conversion to prevent accumulation
\parencite{Mirlas2021,Dou2022,Singh2019}, sequential reuse across cereals,
grasses, forest plantations, and halophytes with innovative evaporation-pond
designs \parencite{Sharma2004}, and land-treatment/pond combinations under
restricted hydraulic loading with preferential-flow control and pre-treatment
before wastewater application \parencite{Christen2010,Duan2015}. Biochar,
phytoremediation, and related reclamation innovations warrant further
development \parencite{Liu2025f,Steppuhn2017,Youngs2000}. Crucially, economic
analysis indicates that feasible drainage-water recycling requires markedly
higher and less volatile irrigated yields together with lower adoption costs
than currently prevail \parencite{Sellars2022}.

\subsection{Acid sulfate soils and irrigation-return management}

Drained acid sulfate (AS) and boreal soils raise a distinct set of specialised
questions. Short-term hydro-meteorological controls on sulfide oxidation,
ripening effects on hydraulic properties, and secondary-mineral formation should
be incorporated into model applications \parencite{Salo2023}; the role of the
nitrogen cycle in sulfide oxidation is unknown in boreal cultivated AS soils
\parencite{Virtanen2017}; the peat-cover thickness required to hinder oxidation
of sulfidic material needs detailed study \parencite{YliHalla2022}; and more
detailed AS-soil maps, a clearer distinction between anthropogenic and organic
acidity, and cessation of drainage of potential AS soils are urged
\parencite{Åström2005}. Amendment strategies---gypsum, elemental sulphur, and
lime for long-term management of problem soils \parencite{MANN1982,He2015},
biochar selected for low contaminant content with urgent long-term impact
assessment \parencite{Tian2023b}, and optimum rates of organic conditioners that
minimise adverse water-quality impacts \parencite{Zhang2015}---require
longer-term verification.

For reclaimed saline, waterlogged, and Vertisol systems, priorities include how
subsurface drainage integrated with crop and residue management affects soil
quality, carbon sequestration, nutrient budgets, and greenhouse-gas fluxes
\parencite{Mukhopadhyay2023,Mukhopadhyay2023b}, the long-term effects of
controlled drainage on salinity \parencite{ElGhannam2023}, salt-crust formation
and dissolved-species transport \parencite{Colombani2016}, leaching effects on
additional elements and organic modifiers \parencite{SoltaniGerdefaramarzi2020},
monitoring of elemental toxicities with waterlogging-stage-specific drainage
scheduling \parencite{Qian2021}, long-term root-growth resistance and
drainable-porosity determination \parencite{Bahçeci2018,Bahceci2008},
additive-effectiveness and land-replacement characterisation \parencite{Isa2023},
and multi-location, multi-condition testing over longer periods
\parencite{Sadhukhan2019}. On the irrigation-return side, needs span tolerant
rootstocks, halophyte germplasm, and transgenic crops \parencite{Cuevas2019},
land-use planning that steers irrigation toward low-salinity aquifers
\parencite{Bethune2008}, and sediment- and return-flow reduction from
furrow-irrigated fields with improved pond removal and nitrogen-leaching-limiting
management \parencite{Bjorneberg2020}. Systematic evaluation of alternative
field-scale irrigation and drainage strategies under ambiguity, expressed as a
statistical likelihood of achieving goals, is advanced \parencite{Alzraiee2013},
supported by long-term geohydrological and crop data to validate simulations
across irrigation/drainage systems \parencite{Yao2017}. Retired land over shallow
water tables, surrounded by irrigated fields, poses unresolved ultimate-use
questions \parencite{Purkey2001}.

\subsection{Sensing, modelling, and design under salinity}

Salinity-specific sensing is a fast-growing cluster: EMI mapping validated across
soil orders and settings \parencite{Narjary2024}, coupled
apparent-conductivity--water-content models \parencite{Yu2025}, and EM38 with
multispectral and hyperspectral remote sensing integrated through interfaces that
explicitly incorporate tile drains \parencite{Singh2020d,Jayawardane2001}.
Coupled water--salt transport modelling remains immature. Models are needed that
simulate salt and nitrate simultaneously in drainage water and the soil profile
\parencite{Ebrahimian2015}; neural networks can predict effluent salinity from
minimal input, but process codes such as SALTMOD may be preferred for detailed
salt balances \parencite{Sarangi2006}, and further evaluation must precede firm
salt-transport recommendations \parencite{Hirekhan2007}. Several strands press
for tighter coupling of water and salt dynamics with unsaturated-zone models such
as HYDRUS, SWAP, and UBMOD-S, incorporating regional layout parameters and
multi-objective optimisation under climatic, textural, and groundwater
variability \parencite{Han2025,Wang2025c,Wang2024c}, and for root-growth modules,
automatic TDR monitoring, and large-dataset validation
\parencite{Liu2022}, extending to climate-change scenarios, deep-soil salt
accumulation, and real-world scalability \parencite{Liu2026}. Sustainability of
saline-water irrigation warrants coupled HYDRUS--EPIC evaluation
\parencite{Feng2019} and better representation of water-table contributions to
storage and evapotranspiration \parencite{Jia2017}. Earlier salinity series
staged validation and salt-flow analysis across successive papers
\parencite{KARAJEH1994}, and proper representation of regional seepage would
demand a fundamentally different modelling approach \parencite{Kelleners2000}.

Design under heterogeneity is the final gap. Integrated redesign of system,
surface-water, and nutrient management is required \parencite{vandenEertwegh2006},
together with texture-specific understanding of the economics of preferential
subsurface drainage \parencite{Neha2022}, and vertical porous flow pipes or fully
charged segments to enhance salt drainage
\parencite{Wang2024d,Yao2024,Jayawardane2001}. New design criteria that protect
crops while minimising both the salinity and the volume of drain water are sought
\parencite{Christen2001}. The most integrative vision unifies these threads:
zonal design algorithms driven by a multi-scale water-and-salt transport model
coupled to remote-sensing inversion for real-time early warning
\parencite{Zhu2025b}---a target that ties salinity sensing, coupled modelling, and
adaptive design into a single research programme.

% ----- frontier/09_unintended-consequences-and-multi-objective-trad.tex -----
\section{Unintended Consequences and Multi-Objective Tradeoffs}

Subsurface drainage and the reductive treatment systems appended to it are
increasingly recognised for effects that fall outside the nitrate- and
yield-centred accounting that has historically justified them. Removing a
nutrient from one pathway can shift it to another, treating drainage water can
manufacture new by-products, and altering the field water regime can set slow
biogeochemical and pedological changes in motion. This section consolidates
these cross-constituent consequences---greenhouse-gas coupling, pollution
swapping, emerging contaminants, trace-element legacies, and drainage-induced
soil evolution---as a single frontier defined less by settled magnitudes than
by unmeasured tradeoffs. The recurring message is that few studies have
quantified more than one consequence at a time, so the net balance among
nitrate reduction, greenhouse-gas emission, secondary pollutant release, and
long-term soil change remains largely conjectural.

\subsection{Greenhouse-gas coupling and carbon export}

Greenhouse-gas (GHG) coupling is a conspicuous gap in drainage accounting.
Fluxes of carbon sequestration, pre-implementation emissions, and dissolved
nitrous oxide in tile drainage and conservation-practice outflows need
quantification \parencite{Mitchell2023}, and coordinated large-scale studies
are required because few have simultaneously measured N\textsubscript{2}O,
CH\textsubscript{4}, and soil-organic-carbon change at the same sites
\parencite{Hall2025}. Longer-term dynamics of wetland carbon storage,
N\textsubscript{2}O, and CH\textsubscript{4}, and the GHG effects of upgrading
drainage infrastructure, remain unresolved \parencite{Mitchell2022}, and carbon
budgets should be integrated into wetland-design models \parencite{Kadlec2012}.

A central strand couples nitrate leaching to indirect N\textsubscript{2}O.
\textcite{Sawamoto2003} call for study of the nitrate--N\textsubscript{2}O
relationship in groundwater and of indirect N\textsubscript{2}O from subsurface
drainage, \textcite{Minamikawa2015} argue for new calculation methods for
indirect N\textsubscript{2}O from leaching and runoff, and the tradeoff between
indirect N\textsubscript{2}O and heterotrophic-denitrification N removal must be
evaluated across groundwater, drainage, river, and estuary
\parencite{Minamikawa2013,Minamikawa2010}. Open questions include
nitrate--N\textsubscript{2}O linkages when N application and precipitation are
temporally disconnected \parencite{OBrien2022}, the magnitude and fate of
N\textsubscript{2}O losses \parencite{Burton2024}, whether dissolved
N\textsubscript{2}O losses exceed surface emissions \parencite{Smith2019}, and
long-term cumulative dissolved-GHG effects of management \parencite{Oliveira2025}.
\textcite{Webb2021} outline an emissions-accounting agenda---seasonal
emission-factor impacts, process-based modelling, gas-transfer-velocity models,
the DOC/NO\textsubscript{3} ratio, and a global database---while
\textcite{Fontes2019} recommend year-round monitoring of both N\textsubscript{2}O
and nitrate leaching. Net-GHG evaluation in longer experiments
\parencite{PrezaFontes2023}, the fate of waterborne methane alongside phosphorus
balances when selecting measures \parencite{Carstensen2021}, comparative
life-cycle assessment from complete N balances \parencite{HernandezRamirez2011},
and biogeochemical measurement and modelling of methane production and oxidation
in cold, poorly drained soils \parencite{Kumar2014b} extend the frontier.

Dissolved-organic-carbon (DOC) dynamics and their ecological implications remain
largely unknown \parencite{Ruark2009,Furukawa2008}, with dissolved-organic-matter
quality, microbial availability, and episodic drainage influence demanding
paired-watershed designs and scale-appropriate sampling
\parencite{Dalzell2011,Bellmore2015}. Treatment structures add their own GHG and
carbon budgets: \textcite{Brunton2024} recommend fuller assessment of full-size
bioreactor methane dynamics, age effects on GHG consumption and release,
``hot moment'' production, and real-time control on outflow redox potential, and
\textcite{Strock2025} note that a 2-h hydraulic residence time raises GHG
emissions and that DOC export needs monitoring. Carbon-export accounting more
broadly is a monitoring gap---aqueous total-carbon export for accurate
sequestration accounting, multi-year data under variable climate, and the
potential of water-level control to minimise dissolved-inorganic-carbon loss
\parencite{Schilling2023}.

\subsection{Pollution swapping and by-product formation in treatment systems}

Reductive treatment introduces tradeoffs that are only partly quantified.
\textcite{Christianson2021} name reducing pollution swapping---especially at
start-up---and improving understanding of the removal of additional pollutants,
compounds, and biologicals as key research areas. The by-products of concern
span the reduction sequence. Nitrite accumulation should be measured in future
bioreactor studies alongside long-term evaluation of phosphorus transport
\parencite{Bell2015}. Mercury methylation is a monitoring priority that
\textcite{Maxwell2022} suggest minimising through an outflow
nitrate-concentration design criterion rather than a mass metric. Redox-driven
phosphorus release is a recurring gap: \textcite{Bailon2025} call for evaluating
start-up flushes with immediate monitoring, aging bioreactors, and total
phosphorus alongside higher-frequency dissolved reactive phosphorus (DRP)
monitoring and drawdown plates, and \textcite{Rivas2020} for continued
multi-season monitoring to confirm longer-term DRP removal and the factors
governing it. The mechanisms of enhanced phosphorus uptake by woodchips under
anoxic--oxic cycling---including polyphosphate-accumulating organisms and whether
anaerobic conditions are necessary---remain to be elucidated
\parencite{Zhang2024b}. \textcite{Strock2025} tie these together, calling for
work on P-removal-material longevity, cold-condition optimisation, DOC-export
monitoring, and strategies for managing P-treatment materials to minimise
unintended phosphorus release, while \textcite{Schilling2023} note that
water-level control may curb dissolved-inorganic-carbon loss. Across these
by-products the common prescription is the same: longer, higher-frequency,
multi-constituent monitoring to bound tradeoffs that single-season demonstrations
cannot resolve.

\subsection{Emerging contaminants, pesticides, pathogens, and antibiotic resistance}

Emerging contaminants and biological interactions broaden the constituent scope
of drainage assessment. Priorities include better constraint of the mechanism
and lifetime of glyphosate in soil and water \parencite{Spiese2023}, planter
modifications to reduce aquatic exposure to neonicotinoid and other seed-applied
pesticides \parencite{Schaafsma2019}, how watershed microbial communities respond
to and control nutrient fluxes \parencite{Knorr2023}, and compound-source
identification across rainfall-runoff events \parencite{Fraindová2022}.
Particulate silicon transport coupled to warming and diatom dynamics
\parencite{Sethna2022}, nutrient stoichiometry and organic-matter fate in aquatic
ecosystems \parencite{Manninen2023,Manninen2018}, and release of accumulated
compounds and resistance genes in sludge-amended soils \parencite{Ghirardini2019}
extend the list.

\subsubsection{Contaminant transport through inlets and macropores}

How contaminants enter and move through drainage networks is explicitly
under-studied. \textcite{Li2017b} recommend research on how nitrate, phosphorus,
and pesticides move through drainage inlets, and on inlet performance under
crop-specific debris. The effective contributing area may differ for water,
tracers, and reactive contaminants \parencite{Akay2008}, and controls on
tile-sediment flocculation remain unresolved \parencite{Stone1997}. Pesticide
transport modelling forms a coherent cluster: particle-mediated colloidal
algorithms and separate macropore accounting for leaching \parencite{Wendell2024},
transport by mobile particles in convection--dispersion models \parencite{Zehe2001},
a quasi-two-dimensional profile distribution and macropore--drain interconnectivity
in chemical routing \parencite{Fox2007b}, slow-immobilisation and
dissociation-kinetics processes for persistent pesticides \parencite{Novak2003},
multi-season field monitoring across soil types \parencite{Dousset2004}, and
vulnerability functions from field characteristics \parencite{Leu2004}. Analytical
and marker methods---hydroxy-triazine determination \parencite{MUIR1978}, a dicamba
marker \parencite{Heaven2012}, continued linuron monitoring \parencite{Gordon2005},
and over-winter freezing/thawing effects on herbicide fate \parencite{Siimes2006}
---complete this cluster.

\subsubsection{Pathogens, microbial ecology, and antibiotic resistance}

Microbial contaminant transport spans several open questions. Pathogen removal
under lower residence times, \textit{Salmonella} and \textit{E.\ coli} removal
mechanisms, and behaviour at higher flows need investigation
\parencite{Soupir2018}. The ecology of indicator bacteria in treatment
wetlands---survival of faecal versus cryptic ``\textit{E.\ coli}-like'' strains
\parencite{Stott2023} and curli expression and downstream risk
\parencite{Truhlar2015}---is unresolved, as is long-term pathogen leaching after
manure application \parencite{Malik2004}. Fecal-source tracking needs
corroborating markers \parencite{Gottschall2013}, irrigation with wastewater
warrants disinfection and reduced cold-season rates \parencite{Ding2021}, and a
broad catchment-scale microbial agenda spans riparian entrapment, constructed
wetlands, and water as a human-infection vector \parencite{Collins2007}.
Antibiotic-resistance genes are an explicitly emerging frontier: their diversity,
persistence, and dissemination require qPCR to complement metagenomics
\parencite{Smith2019b}, and further study is needed to link resistance genes to
host identity and transfer mechanisms \parencite{Rieke2018}.

\subsection{Trace-element legacies and ecotoxicology}

Trace-element contamination associated with drainage in the western United States
represents a long-horizon monitoring commitment. \textcite{BENSON1993} outline
continued biological and groundwater monitoring, research into 10--20-year
selenium-dissipation strategies, annual management planning, and pilot-scale
microbial-volatilisation demonstration, and \textcite{SAIKI1987} note further
studies to determine whether aquatic animals elsewhere in the San Joaquin Valley
are contaminated by selenium. Ecotoxicological monitoring of other contaminants
extends the frontier: \textcite{Campbell2023} call for mechanistic work on
reduced amphibian cerebellum size, field examination of imidacloprid and its
olefin metabolite, and regular monitoring of neonicotinoids near agricultural
waterways with tighter restrictions to protect non-target organisms, while
\textcite{Dreyfuss2018} note that effects on \textit{Galba truncatula} raise
conservation questions given its role as intermediate host of \textit{Fasciola
hepatica}. These constituents share a common feature: their consequences unfold
over decades and reach organisms far removed from the field, so short monitoring
windows systematically understate them.

\subsection{Drainage-induced soil evolution, health, and biology}

Slower-acting soil and biological responses to drainage are repeatedly flagged as
understudied. The soil-biology response to changing drainage hydrology requires
longer-term evaluation of biological indicators under standardised sampling
\parencite{Sherbine2023}, study of interactions between subsurface drainage and
soil physicochemical properties and microbial activity \parencite{Yang2025}, and
continued research on microbial community structure and function to explain
tillage and drainage impacts \parencite{Kaur2024b,Kumar2014}. Robust soil-health
diagnostics are sought---scalable nematode measurement within a soil-health index
\parencite{Li2023b}, valid and repeatable soil-quality indicators in a systems
framework \parencite{Nakajima2015}, soil-enzyme activities as management
diagnostics \parencite{Singh2007}, proof that no-till reduces nitrate leaching
sustainably \parencite{Masri2024}, and dynamic compaction indicators with
long-term restoration trials \parencite{Caron2025}. Specific biological questions
include \textit{Fusarium} abundance and soybean root rot under long-term drainage
\parencite{Han2017}, disruption of crack-to-burrow linkages by conservation
tillage \parencite{Shipitalo2004}, evapotranspiration quantification to explain
drain-flow differences \parencite{Daigh2014}, cultivar drought tolerance guiding
system selection \parencite{DarziNaftchali2017b}, and forage-sorghum nitrate
testing \parencite{Nash2017}.

A distinctive line treats subsurface drainage as a driver of long-term,
potentially chaotic soil evolution. \textcite{Montagne2009} argue that
quantifying human-induced soil evolution as a function of time---including
drainage's impact on clay-mineral dissolution---should be a research priority,
and that soil evolution under climate change must be explicitly studied and
incorporated in global models \parencite{Montagne2010}; significant material
losses over a few decades make anthropogenic perturbation a powerful tool for
studying soil-evolution timescales \parencite{Montagne2013}. Whether
drainage-induced evolution continues or stabilises over decades, resolved through
measurements along drains at different slope positions and network ages, remains
open \parencite{Frison2009}, and a long classical experiment continues to test
whether soil structure is eventually affected \parencite{HOOGHOUDT1952}.
Empirical prediction of earthworm activity from tillage, residue cover, and
particle size is proposed \parencite{Ouellet2008}, and sustained drainage
operation and cultivation is framed as a management imperative for long-term soil
health \parencite{Singh2007}. Collectively, these consequences reframe drainage
as an agent of slow, coupled biogeochemical and pedological change whose full
cost cannot be read from a few years of outlet monitoring.

% ----- frontier/10_toward-integrated-targeted-and-adoptable-solutio.tex -----
\section{Toward Integrated, Targeted, and Adoptable Solutions}

The preceding frontiers make clear that better process understanding, models,
and monitoring are necessary but not sufficient: translating frontier science
into water-quality and productivity gains depends on how practices are
combined, where they are placed, how they perform under a changing climate, and
whether they are economically and institutionally adoptable. This section draws
the integrative and socio-economic strands of the literature together. Its
recurring message is that no single practice suffices, that biophysical
questions are inseparable from human ones, and that the decisive advances will
come from coupling mechanistic understanding with spatial targeting, economics,
policy, and shared decision-support infrastructure.

\subsection{Bundling, sequencing, and spatial targeting of practices}

A pervasive conclusion is that no single best management practice (BMP) is a
silver bullet \parencite{Sunohara2016}, and that multiple BMPs in series are
needed---with in-field/edge-of-field combinations judged highly promising yet
little studied, and wetland and water-retention measures still to be developed
for sensitive downstream waters \parencite{Lenhart2016}. Attaining water-quality
goals will require combinations implemented at a high level in a targeted manner
\parencite{Christianson2014}. Open questions concern how best to pair
conservation tillage with cover crops and in-field with edge-of-field practices
while accounting for legacy effects \parencite{Makhtoumi2025}, how to formalise
optimisation of synergistic combinations \parencite{Kreig2019}, and how
practices interact and what nitrogen-removal capacity each provides given better
mapping of tile-drained land \parencite{McLellan2015}. Broad programmes are
proposed to compare drained and undrained systems, to bundle bioreactors,
phosphorus reactors, controlled drainage (CD), and low-impact development
alongside internal-loading management \parencite{Baulch2019,Husk2024,Husk2017},
to develop cost-effective BMPs that support cover crops and mitigate warming
\parencite{Gupta2023}, and to test intercrops and hydrology-altering management
for phosphorus reduction \parencite{HernandezRamirez2011,GAYNOR1995b}.

Closely related is targeting practices to the parts of the landscape that matter
most. Identifying hydrologically active runoff-generating areas as mitigation
priorities is recommended \parencite{Melland2008}, but whether elevated loads
originate from many diffuse sources or a few hot spots
\parencite{Irvine2019}, whether targeted or regionally saturated cover-crop
implementation is needed for consistent reductions \parencite{Speir2022}, and how
to delineate zones of higher nitrate leaching \parencite{Bakhsh2005b} remain
open. Targeting critical source areas and hydrologically connected hotspots is
repeatedly stressed \parencite{McDowell2015,Brendel2019}, prioritisation is
expected to improve the return on conservation investment
\parencite{Makhtoumi2025}, and automating spatial-targeting principles into an
accessible GIS tool with whole-basin surveys and cost--benefit analysis is a
priority \parencite{Kvítek2023}. Further scaling needs include combining GIS with
hydrological modelling and saturation risk to find nitrate thresholds across
soil texture and temperature \parencite{Smith2023}, resolving how drainage
density controls stream solute dynamics \parencite{Kennedy2012}, and identifying
critical precipitation periods rather than annual averages \parencite{Salo2006}.

\subsection{Climate change, extremes, and system resilience}

Climate adaptation frames many priorities. The overarching need is to keep
identifying practices and combinations that increase system resilience to
climate extremes across additional soils and longer periods
\parencite{Saha2024}, supported by modelling agendas spanning additional
periods, representative concentration pathways, management options, cultivars,
cover-crop types, and soils \parencite{Malone2020}, and by field-scale
integrated biophysical modelling that resolves within-county variation in
fertiliser and manure inputs under multiple CMIP5/CMIP6 forcings with attention
to perennial-grass co-benefits \parencite{Ferin2023}. A central design
challenge is cropping and drainage systems that function under both drought and
flood, including engineering solutions to capture excess floodwater
\parencite{Dietzel2016}. Adapting drainage design itself to a changing climate is
an explicit priority, tied to understanding weather--soil--groundwater--discharge
interactions under increased precipitation and to evaluating designs over longer
periods and different climates
\parencite{Deelstra2015,Tuohy2018,Malota2024,Deelstra2014,Deelstra2014b}.
Remaining needs span the costs, benefits, and payback of adaptation and
carbon-farming knowledge exchange \parencite{PeltonenSainio2021}, the effect of
extreme events on cover-crop performance \parencite{Gupta2023b}, nitrogen
consequences under warming \parencite{Haghnazari2020}, and testing across
additional rainfall environments, soils, and rotations to confirm
drainage--management interactions
\parencite{Takeshima2023,Liu2025,Grewal2025,Negm2014,BENOIT1967}.

\subsection{Economics, incentives, and cost-effectiveness}

A large share of the frontier is economic. Cost--benefit and cost-effectiveness
analyses are repeatedly flagged as future work---for surface amendments and
energy savings \parencite{Wamono2016b}, comprehensive short- and long-term
subsurface-drainage benefits \parencite{Guo2025}, paddy-field cost and
consumption \parencite{Jia2023}, edge-of-field cost-efficiency
\parencite{Mendes2020}, and the opportunity cost of drainage investment together
with policies that motivate improvement
\parencite{Wichelns2002,THOMAS1995,Kladivko2005,Saadat2017,Manik2019,Helmers2012}.
Identifying long-term savings---such as greater fertiliser-nitrogen-use
efficiency---to offset higher installation costs, alongside incentive and
cost-share design, is recommended \parencite{Castellano2019,Sunohara2015}, and
cost reduction is itself a research target through local manufacture and new
low-cost techniques and materials \parencite{Datta2002,Yannopoulos2020}.
Complementary needs include advanced case-by-case cost analyses to promote 4R
implementation \parencite{Liu2021}, more data on the cost-effectiveness and
ecological benefits of conservation practices \parencite{Yuan2022}, and explicit
incorporation of practice costs into future modelling \parencite{McLellan2015}.
Incentive mechanisms warrant study: continued positive incentives and greater
weight on non-excludable benefits in multifunctional landscapes
\parencite{Christianson2014}, cost-share and risk-insurance instruments
\parencite{Jaynes2012,Karlen2005}, and nutrient-trading programmes for
shallow-drain adoption \parencite{Ghane2021b}. Feasibility analyses caution that
drainage-water recycling requires markedly higher, less volatile irrigated
yields and lower adoption costs \parencite{Sellars2022}, motivating work on the
economics of cover crops \parencite{Waring2020}, long-term cumulative yield and
profit of drainage installation \parencite{Mourtzinis2021}, and integrated
biophysical--economic modelling at field scale \parencite{Ferin2023}.

\subsection{Policy, institutions, adoption, and the human dimension}

Biophysical understanding alone is insufficient; governance, capacity, and
farmer engagement are decisive. A repeated recommendation is to shift from
standardised toward flexible, location-specific solutions that reflect
stakeholder preferences, balancing top-down against bottom-up participatory
management and linking farmers' tacit knowledge with researchers' explicit
knowledge \parencite{Ritzema2011,Ritzema2016,Ritzema2023}. Capacity development,
education, and policy change to lower adoption barriers are emphasised
\parencite{Castellano2019,SHIRMOHAMMADI1995,Kesicka2023b}, along with governance
frameworks, inter-agency coordination, dedicated drainage institutions, and
time-bound action plans
\parencite{Singh2019c,Qureshi2008,Gopalakrishnan2007,Ritzema2008b,Valipour2014,Quinn2009},
and awareness, subsidy, micro-finance, and participatory
operation-and-maintenance programmes
\parencite{Satyanarayana2007,Manjunatha2004,BenAissa2013,MCFARLANE1992}.

Adoption and social dimensions are stressed across contexts. Detailed
cost--benefit analyses, joint farmer--researcher experimentation, and capacity
building are prioritised, with adoption ultimately depending on farmers'
long-term cost--benefit perspective \parencite{Okwany2016}; creating know-how and
farmer confidence before project execution is essential \parencite{Ghumman2011};
and key practice-use issues---where to apply BMPs, how to bundle them, and
transferability to commercial farms---remain open \parencite{Ward2018}.
Stakeholder willingness to adopt must inform proposed measures
\parencite{Cuevas2019}, research integrating ecological, pollution-related, and
economic impacts with models for optimal policy design is urged
\parencite{Saikkonen2014}, and incentivising on-farm research while scaling
individual efforts to meso- and macro-scale policy is called for
\parencite{David2015}. Policy reforms sought include legislative reassessment to
remove barriers to natural-water-retention measures \parencite{Kvítek2023},
national support to retain open ditch-and-field networks \parencite{Marja2013},
producer-participation and pilot-programme information for policy optimisation
\parencite{Malone2023}, and integrated packages---reduced plantings, farmer
compensation, and gradual subsidy liberalisation---in salinity contexts
\parencite{Kotb2000}.

\subsection{Integrative landscape-scale frameworks and decision support}

The most explicit long-range vision is numerically explicit landscape-scale
models that couple the processes governing ditch services, optimise the spatial
arrangement of practices according to connectivity, and define multi-service
maximisation strategies informed by stakeholder-valued rankings---prerequisite
to which is better mapping of ditch hydraulic, vegetation, and sorption
properties \parencite{Dollinger2015,Herzon2008}. Landscape-scale approaches that
combine on-farm change with off-site wetland and riparian restoration are
broadly endorsed \parencite{Blann2009,Goeller2019}, and truly integrated
watershed management requires better information on controlling processes such
as the groundwater-head distribution linking topography and geology to the drain
fraction \parencite{Mahmood2023}. Coupled field--watershed models enable
water-table-management scenario analyses not possible with component models used
separately \parencite{Dayyani2010}, frameworks should remain extensible to
additional processes and dynamically calibrated indices
\parencite{Basu2010,Tao2025,Verma2014}, and modelling should increasingly guide
practice placement and stacking
\parencite{Kröger2013,Tomer2015,Lemke2011,McDowell2021,King2022,Roley2016}.

Emerging sensing and data technologies are widely seen as enabling this next
phase. Broad agendas span optimal conditions for detecting drainage patterns and
extracting drain lines, finer-resolution soil moisture, machine-learning mapping
of soil properties, and within-field nitrogen-stress mapping
\parencite{Khanal2020}, alongside spatial frameworks that combine climate and
soil type with geo-referenced farmer-reported data to estimate field-specific
benefits \parencite{Cassman2025}. Decision-support needs include validating
aggregate edge-of-field losses and delivery ratios \parencite{Reid2018},
testing soil-wetness-index tools across soils \parencite{Julien2023}, integrating
water and nutrient management with precision agriculture \parencite{Jung2024},
and coupling downscaled climate information with an understanding of how farmers
incorporate excess-precipitation experience into decisions \parencite{Morton2015}.
Better indices and shared tools are sought---phosphorus and nutrient indices that
represent subsurface drainage, wetland and riparian-buffer removal, and
stream-network connectivity \parencite{Buczko2007,Pärn2012}, extensible
nutrient-credit-trading tools for drainage water management \parencite{Negm2016},
a Soil Vulnerability Index that distinguishes surface from subsurface drainage
\parencite{Baffaut2020}, and coupling with ecosystem-service tools such as INVEST
to assess tradeoffs \parencite{McLellan2015}. Underlying all of these is the data
imperative: better, longer, and more spatially representative datasets---spanning
water and solute fluxes, greenhouse-gas emissions, root and biomass development,
and management history---are the foundational requirement for validating models,
resolving tradeoffs, and delivering site-specific, climate-resilient guidance
\parencite{Ma2025,SKAGGS1994,Deng2011}.

% ----- frontier/11_conclusions.tex -----
\section{Conclusions}

Read as a whole, the frontier of subsurface-drainage research is defined
less by an absence of ideas than by an absence of the evidence, tools, and
coordination needed to convert ideas into transferable guidance. Across
every domain surveyed here---monitoring, modelling, network mapping,
nutrient and water-table management, edge-of-field treatment, salinity and
reuse, unintended consequences, and adoption---the field's most secure
statements are about the limits of its own evidence. Headline results rest
disproportionately on short records, single sites, and narrow
agro-ecological windows, and their authors caution explicitly against
premature generalisation. The recurring signal is therefore a precondition
rather than a conclusion: progress in each domain is bounded by the length,
density, spatial replication, and consistency of the underlying
observations \parencite{Rogger2017,Kladivko2001}. This makes the observational foundation a frontier in its own
right, and one on which all others depend.

Four converging priorities emerge from that shared diagnosis. The first is
coordinated, long-term, multi-scale data. The single most widely voiced need
is for records long enough to separate management effects from climatic
noise and from slow, directional change in soils and hydrology, dense and
year-round enough to resolve the non-growing-season and event-scale windows
in which losses concentrate, and geographically broad enough to test whether
site-specific findings hold across other soils, climates, crops, and
management regimes. Layered onto record length are harmonisation and
comparability: standardised methods, complete site characterisation,
shared databases, and openly reported campaigns, without which accumulating
studies cannot support the meta-analyses and typologies that would
generalise them. Tracers, isotopes, and high-frequency source
apportionment recur as the means to resolve the pathways that bulk
flow-and-concentration records leave ambiguous, and experimental designs
capable of detecting the effect sought---paired basins, nested scales,
and defensible analysis of the unreplicated field trials that dominate the
literature---form the bridge between a lengthening record and transferable
conclusions.

The second priority is strengthened and validated model components. Because
field trials cannot span every combination of practice, soil, and climate,
process-based drainage and water-quality models are expected to carry much
of the predictive burden; yet their credibility is constrained less by their
formulation than by the narrowness of the conditions under which they have
been tested. Validation across a wider envelope, the replacement of
empirical coefficients with functions of measurable site attributes, and the
process-level measurements needed to constrain the least mature routines---
denitrification and the nitrogen cycle, phosphorus and its dominant
preferential-flow pathway, cold-season dynamics, and coupled seepage and
groundwater exchange---are near-universal priorities \parencite{Youssef2006,Radcliffe2015}. So too are explicit
attention to identifiability and uncertainty, so that models are reported
with propagated bounds rather than as point predictions, and the disciplined
integration of machine learning and coupled model structures, whose own
validation requirements the field is only beginning to confront. A parallel
and equally unresolved gap is spatial: the physical drainage network is
largely unmapped, so uncertainty about the network itself propagates into
every downstream prediction, and translating instrumented-plot knowledge to
the watersheds where water-quality goals are actually set remains a genuinely
open methodological problem requiring non-invasive detection, shared
databases, intensively monitored bridging catchments, and mechanistic
treatment of channel-network, connectivity, and depressional processes \parencite{Amado2018,Naz2009}.

The third priority is bundled, targeted, and climate-resilient management.
Across nutrients, water-table control, drainage design, and edge-of-field
treatment, the levers are better described than optimised. No single practice
is a silver bullet; the decisive gains lie in combining in-field and
edge-of-field practices, sequencing them, and placing them where they matter
most---on the hydrologically active areas, critical source areas, and
connected hotspots that generate a disproportionate share of loads. The
common thread is a shift from passive, uniform prescription toward active,
location-specific control: from one-way conveyance toward managed water
tables, automation, and forecast-based operation; from agronomic nutrient
rates toward site-specific crediting of legacy and mineralisable sources and
management targeted to the temporal windows of loss; from demonstration-scale
treatment structures toward standardised, transferable design across
climates and maturation stages \parencite{Randall2001,Strock2011,Christianson2021}. Underlying all of these is climate
adaptation: cropping and drainage systems, and the design criteria that
govern them, must function under both drought and flood, and must be
re-evaluated against long records rather than the short-term simulations on
which much current guidance rests.

The fourth priority is interdisciplinary work that links mechanism, tools,
and viable adoption, and it is inseparable from a growing recognition of
consequences that fall outside the nitrate- and yield-centred accounting
that has historically justified drainage. Removing a nutrient from one
pathway can shift it to another; treating drainage water can manufacture new
by-products; and altering the field water regime can set slow biogeochemical
and pedological change in motion. Greenhouse-gas coupling, pollution
swapping, emerging contaminants, trace-element legacies, and
drainage-induced soil evolution share a common feature: few studies have
quantified more than one consequence at a time, so the net balance among
competing objectives remains largely conjectural, and their full costs
cannot be read from a few years of outlet monitoring \parencite{Mitchell2023,Kadlec2012}. Resolving these
tradeoffs, and translating frontier science into practice, ultimately
depends on the human dimension---economics, incentives, policy,
institutions, and farmer engagement---and on integrative landscape-scale
frameworks and shared decision-support infrastructure that couple mechanistic
understanding with spatial targeting, cost-effectiveness, and stakeholder
knowledge \parencite{McLellan2015,Christianson2014}. Biophysical insight alone is necessary but not sufficient; without
economically and institutionally adoptable solutions, better science does not
become better outcomes.

These four strands are mutually reinforcing rather than independent.
Longer, harmonised data are what allow model components to be validated and
their uncertainty bounded; validated, coupled models are what allow
field-scale knowledge to be scaled to watersheds and management to be
targeted and made climate-resilient; and targeted, multi-objective
management is what makes the consequences worth quantifying and the
economics worth analysing. The unifying message of this review is therefore
that the next decade of progress in subsurface drainage will be secured less
by new equations or novel practices than by the sustained, well-instrumented,
openly reported, and geographically diverse field campaigns needed to test
the ideas already in hand, by the coordinated data infrastructure and
standards that make those campaigns cumulative, and by the interdisciplinary
integration that ties mechanism to tools and tools to adoption. The task
ahead is one of consolidation and coordination at scale---an agenda in which
data, models, management, and human-dimension advances must progress
together, because each is bounded by the others.


\printbibliography
\end{document}
